% Leçon 39 — Grammaire : verbe modal 可以 ; 的 、 得、 地 — Chinois Tle A — Cours signé OnBuch+
% Programme officiel MINESEC. Compiler avec : tectonic ZH39-grammaire-verbe-modal.tex
\newcommand{\DOCMATIERE}{Chinois}
\newcommand{\DOCNIVEAU}{Tle A}
\newcommand{\DOCMODULE}{Module 5 : Mass médias et télécommunication}
\newcommand{\DOCLECON}{ZH39}
\newcommand{\DOCTITRE}{Leçon 39 — Grammaire : verbe modal 可以 ; 的 、 得、 地}
% OnBuch+ — préambule « cours signés » ENRICHI (compilé avec tectonic / XeLaTeX)
% Système visuel « Pop » d'OnBuch+ : fond crème (jamais blanc), bordures encre
% épaisses, ombres « dures » décalées, titres Archivo Black en capitales.
% Version MATHÉMATIQUES : graphes (pgfplots/tikz), boîtes pédagogiques
% (définition, propriété, méthode, attention, le savais-tu, exercice résolu,
% exercice, TP, bilan), page de garde et signature OnBuch+ ancrée en bas.
%
% Macros attendues dans main.tex AVANT \input{preamble} :
%   \DOCMATIERE  \DOCNIVEAU  \DOCMODULE  \DOCLECON (numéro)  \DOCTITRE
\documentclass[11pt]{article}

\usepackage{fontspec}
\usepackage[french]{babel} % français = langue principale ; le chinois via \zh{..} / textebox (xeCJK)
\usepackage{xeCJK}
\usepackage{pifont} % \ding{51} \ding{55} utilisés par les modèles
\usepackage[a4paper,margin=2cm,top=3.4cm,bottom=2.8cm,headheight=1.4cm,footskip=1.3cm]{geometry}
\usepackage{amsmath,amssymb,amsfonts}
\usepackage{mathtools} % \xmapsto, \coloneqq... souvent utilisés par les modèles
\usepackage{cancel} % simplifications barrées (bilans de forces)
% \abs{z} (module, valeur absolue) et \norm{u} (norme), utilisés par les modèles
\providecommand{\abs}{}\let\abs\relax\DeclarePairedDelimiter{\abs}{\lvert}{\rvert}
\providecommand{\norm}{}\let\norm\relax\DeclarePairedDelimiter{\norm}{\lVert}{\rVert}
\usepackage[table]{xcolor} % [table] : fournit aussi \rowcolors (alternance de lignes)
\usepackage{pagecolor}
\usepackage{tikz}
\usetikzlibrary{arrows.meta,positioning,calc,shapes.geometric,shapes.misc,shapes.arrows,decorations.pathmorphing,decorations.pathreplacing,decorations.markings,patterns,shadows,mindmap,trees,fit,backgrounds,matrix,intersections,angles,quotes}
\usepackage{pgfplots}
\usepackage[version=4]{mhchem} % équations nucléaires \ce{^{235}_{92}U}
\usepackage[european,siunitx]{circuitikz} % schémas électriques
\pgfplotsset{compat=1.18}
\usepgfplotslibrary{fillbetween,groupplots} % aires entre courbes (calcul intégral)
\usepackage{enumitem}
\usepackage{fancyhdr}
% [most] charge toutes les bibliothèques tcolorbox (listings, minted, raster,
% documentation...) et gonfle inutilement la mémoire TeX sur un document long
% (~300 boîtes) : on ne charge que ce qui est réellement utilisé.
\usepackage[skins,breakable]{tcolorbox}
\usepackage{graphicx}
\usepackage{colortbl}
\usepackage{booktabs}
\usepackage{tabularx}
\usepackage{diagbox} % case d'en-tête diagonale des tableaux à double entrée
% Colonnes extensibles centrée / gauche / droite (C, L, R), souvent utilisées par les modèles
\newcolumntype{C}{>{\centering\arraybackslash}X}
\newcolumntype{L}{>{\raggedright\arraybackslash}X}
\newcolumntype{R}{>{\raggedleft\arraybackslash}X}
\usepackage{array}
\usepackage{multicol}
\usepackage{multirow}
\usepackage{siunitx}
% parse-numbers=false : accepte aussi \SI{2,5 \times 10^{-3}}{...}
\sisetup{locale=FR,output-decimal-marker={,},group-minimum-digits=4,parse-numbers=false}
\DeclareSIUnit\ohm{\ensuremath{\symup{\Omega}}} % U+2126 absent de la police
\DeclareSIUnit\division{div} % réglages d'oscilloscope (ms/div, V/div)
\DeclareSIUnit\tr{tr} % tours (vitesse de rotation, ex. \SI{1500}{\tr\per\minute})
\DeclareSIUnit\molar{M} % concentration molaire (ex. \SI{3}{\molar}, \SI{1}{\micro\molar})
\DeclareSIUnit\an{an} % année (le modèle confond parfois avec \year, un compteur TeX) — ex. \SI{5}{\kilo\meter\per\an}
\DeclareSIUnit\FCFA{FCFA} % franc CFA (ex. \SI{500}{\FCFA\per\kilo\gram})
\DeclareSIUnit\degreeBrix{\textdegree Brix} % teneur en sucre d'un sirop/jus (réfractométrie)
\DeclareSIUnit\month{mois} % le modèle utilise parfois \month, un compteur TeX, comme unité de durée
\usepackage{qrcode}
% \degree : commande fréquemment utilisée par les modèles pour un angle ou une
% température en dehors de \SI{}{} (non fournie par les paquets chargés).
\providecommand{\degree}{^{\circ}}
% \qed (fin de démonstration), utilisé par les modèles sans amsthm
\providecommand{\qed}{\hfill$\square$}
% \barre : double barre des valeurs interdites dans un tableau de variations (tkz-tab)
\providecommand{\barre}{\ensuremath{\|}}
% environnement proof (amsthm non chargé) utilisé par les modèles
\ifdefined\proof\else\newenvironment{proof}[1][Démonstration]{\par\noindent\textit{#1.}\ }{\qed\par}\fi
\usepackage{lastpage}
\usepackage{hyperref}
\hypersetup{hidelinks}

% ============================================================
% Palette « Pop » OnBuch+ — mêmes hex que la classe P (plus_ui.dart)
% ============================================================
\definecolor{poporange}{HTML}{F59321}
\definecolor{popdark}{HTML}{E07A0C}
\definecolor{popcream}{HTML}{FFF6E8}
\definecolor{popink}{HTML}{1C1714}
\definecolor{poppurple}{HTML}{7A5AE0}
\definecolor{popgreen}{HTML}{2E9E6A}
\definecolor{poppink}{HTML}{E0457B}
\definecolor{popblue}{HTML}{2F7DE1}
\definecolor{popgold}{HTML}{C98A12}
\definecolor{popmuted}{HTML}{8A7F76}
\definecolor{popline}{HTML}{EADFD2}
\definecolor{popgray}{HTML}{8A7F76}
\colorlet{popgrey}{popgray}
% Alias tolérants (le modèle nomme parfois les couleurs différemment)
\colorlet{popwhite}{white}
\colorlet{popblack}{popink}
\colorlet{popred}{poppink}
\colorlet{popyellow}{popgold}
% Teintes claires (fonds de tableaux/encadrés) — le modèle nomme parfois
% les couleurs avec un suffixe L (ex. popblueL) sans les avoir définies.
\colorlet{poporangeL}{poporange!20}
\colorlet{popdarkL}{popdark!20}
\colorlet{poppurpleL}{poppurple!20}
\colorlet{popgreenL}{popgreen!20}
\colorlet{poppinkL}{poppink!20}
\colorlet{popblueL}{popblue!20}
\colorlet{popgoldL}{popgold!20}
\colorlet{popredL}{popred!20}
\colorlet{popgrayL}{popgray!20}
% Teintes claires pour fonds de tableaux / zones
\colorlet{popblueL}{popblue!10!white}
\colorlet{poporangeL}{poporange!14!white}
\colorlet{popgreenL}{popgreen!12!white}
\colorlet{poppurpleL}{poppurple!12!white}
\colorlet{poppinkL}{poppink!12!white}
\colorlet{popgoldL}{popgold!14!white}

\pagecolor{popcream}

% ============================================================
% Typographie
% ============================================================
\defaultfontfeatures{Ligatures=TeX}
\setmainfont{PlusJakartaSans-Medium.ttf}[Path=../../../fonts/,BoldFont=PlusJakartaSans-Bold.ttf]
% Symboles, API et emojis absents de la police principale : repli sur DejaVu Sans (caractères actifs)
\newfontface\symfont{DejaVuSans.ttf}[Path=../../../fonts/]
\makeatletter
\newcommand\symactive[1]{\catcode#1=13 \begingroup\lccode`\~=#1 \lowercase{\endgroup\def~}{{\symfont\char#1}}}
\newcommand\symblank[1]{\catcode#1=13 \begingroup\lccode`\~=#1 \lowercase{\endgroup\def~}{}}
\newcommand\symhyph[1]{\catcode#1=13 \begingroup\lccode`\~=#1 \lowercase{\endgroup\def~}{\mbox{-}}}
\newcount\symc
\newcommand\symrange[2]{\symc=#1 \@whilenum\symc<#2 \do{\expandafter\symactive\expandafter{\the\symc}\advance\symc by 1 }}
\newcommand\symblankrange[2]{\symc=#1 \@whilenum\symc<#2 \do{\expandafter\symblank\expandafter{\the\symc}\advance\symc by 1 }}
\makeatother
\symrange{"0250}{"02FF}\symactive{"2713}\symactive{"2714}\symactive{"2717}\symactive{"2718}\symactive{"2610}\symactive{"2611}\symactive{"2612}
\symrange{"2600}{"27BF}
\catcode"2705=13 \begingroup\lccode`\~="2705 \lowercase{\endgroup\def~}{{\symfont\char"2713}}\symblank{"26BD}\symblank{"26A1}
\symrange{"2460}{"257F}\symactive{"223C}\symactive{"2248}\symactive{"2260}
\symactive{"01F8}\symactive{"01F9}\symactive{"1E3E}\symactive{"1E3F}
\symhyph{"2011}\symhyph{"2E1A}\symblankrange{"1F300}{"1FAFF}\symblank{"FE0F}
% Caractères chinois : WenQuanYi Zen Hei (sous-ensemble GB2312 + pinyin), gras/italique simulés
\setCJKmainfont{WenQuanYiZenHei-subset.ttf}[Path=../../../fonts/,AutoFakeBold=3,AutoFakeSlant=0.2]
\xeCJKsetup{CJKspace=false}
% Pinyin : voyelles à caron (ǎ ǐ ǒ ǔ ǖ ǘ ǚ ǜ) absentes de la police principale -> caractères actifs
\newfontface\pyfont{WenQuanYiZenHei-subset.ttf}[Path=../../../fonts/]
\newcommand\pyactive[1]{\catcode#1=13 \begingroup\lccode`\~=#1 \lowercase{\endgroup\def~}{{\pyfont\char#1}}}
\pyactive{"01CD}\pyactive{"01CE}\pyactive{"01CF}\pyactive{"01D0}\pyactive{"01D1}\pyactive{"01D2}\pyactive{"01D3}\pyactive{"01D4}\pyactive{"01D5}\pyactive{"01D6}\pyactive{"01D7}\pyactive{"01D8}\pyactive{"01D9}\pyactive{"01DA}\pyactive{"01DB}\pyactive{"01DC}
% Maths en sans-serif (Fira Math) avec chiffres et lettres droites de Plus
% Jakarta, pour que les formules se fondent dans le texte.
\usepackage[math-style=ISO,bold-style=ISO]{unicode-math}
\setmathfont{FiraMath-Regular.otf}[Path=../../../fonts/]
\setmathfont{PlusJakartaSans-Medium.ttf}[Path=../../../fonts/,range={up/num,up/{Latin,latin},\mathup}]
\usepackage{icomma} % virgule décimale sans espace en mode math
% Caractères mathématiques Unicode absents des polices de texte (Plus Jakarta,
% Archivo Black) : rendus par leur équivalent mathématique, en texte comme en
% formule — sinon ils s'affichent en carré vide (ex. « Arithmétique dans ℤ »).
\usepackage{newunicodechar}
\newunicodechar{ℝ}{\ensuremath{\mathbb{R}}}
\newunicodechar{ℤ}{\ensuremath{\mathbb{Z}}}
\newunicodechar{ℂ}{\ensuremath{\mathbb{C}}}
\newunicodechar{ℕ}{\ensuremath{\mathbb{N}}}
\newunicodechar{ℚ}{\ensuremath{\mathbb{Q}}}
\newunicodechar{𝔻}{\ensuremath{\mathbb{D}}}
\newunicodechar{α}{\ensuremath{\alpha}}
\newunicodechar{β}{\ensuremath{\beta}}
\newunicodechar{γ}{\ensuremath{\gamma}}
\newunicodechar{δ}{\ensuremath{\delta}}
\newunicodechar{ε}{\ensuremath{\varepsilon}}
\newunicodechar{θ}{\ensuremath{\theta}}
\newunicodechar{λ}{\ensuremath{\lambda}}
\newunicodechar{μ}{\ensuremath{\mu}}
\newunicodechar{σ}{\ensuremath{\sigma}}
\newunicodechar{τ}{\ensuremath{\tau}}
\newunicodechar{φ}{\ensuremath{\varphi}}
\newunicodechar{ψ}{\ensuremath{\psi}}
\newunicodechar{ω}{\ensuremath{\omega}}
\newunicodechar{Σ}{\ensuremath{\Sigma}}
% majuscules produites par \MakeUppercase dans les titres (α-aminés -> Α)
\newunicodechar{Α}{\ensuremath{\alpha}}
\newunicodechar{Β}{\ensuremath{\beta}}
\newunicodechar{Γ}{\ensuremath{\Gamma}}
\newunicodechar{Δ}{\ensuremath{\Delta}}
\newunicodechar{Ε}{\ensuremath{\varepsilon}}
\newunicodechar{Θ}{\ensuremath{\Theta}}
\newunicodechar{Λ}{\ensuremath{\Lambda}}
\newunicodechar{Μ}{\ensuremath{\mu}}
\newunicodechar{Τ}{\ensuremath{\tau}}
\newunicodechar{Φ}{\ensuremath{\Phi}}
\newunicodechar{Ψ}{\ensuremath{\Psi}}
\newunicodechar{Ω}{\ensuremath{\Omega}}
\newunicodechar{↦}{\ensuremath{\mapsto}}
\newunicodechar{∈}{\ensuremath{\in}}
\newunicodechar{∉}{\ensuremath{\notin}}
\newunicodechar{∧}{\ensuremath{\wedge}}
\newunicodechar{⇔}{\ensuremath{\Leftrightarrow}}
\newunicodechar{⟺}{\ensuremath{\Longleftrightarrow}}
\newunicodechar{⇒}{\ensuremath{\Rightarrow}}
\newunicodechar{⟹}{\ensuremath{\Longrightarrow}}
\newunicodechar{∘}{\ensuremath{\circ}}
\newunicodechar{∪}{\ensuremath{\cup}}
\newunicodechar{∩}{\ensuremath{\cap}}
\newunicodechar{≡}{\ensuremath{\equiv}}
\newunicodechar{⊂}{\ensuremath{\subset}}
\newunicodechar{⊥}{\ensuremath{\perp}}
\newunicodechar{∃}{\ensuremath{\exists}}
\newunicodechar{∀}{\ensuremath{\forall}}
\newunicodechar{∛}{\ensuremath{\sqrt[3]{\,}}}
\newunicodechar{∜}{\ensuremath{\sqrt[4]{\,}}}
\newunicodechar{⃗}{\ensuremath{{}^{\to}}}
\newunicodechar{ⁿ}{\ifmmode^{n}\else\textsuperscript{\ensuremath{n}}\fi}
\newunicodechar{ˣ}{\ifmmode^{x}\else\textsuperscript{\ensuremath{x}}\fi}
\newunicodechar{ᵇ}{\ifmmode^{b}\else\textsuperscript{\ensuremath{b}}\fi}
\newunicodechar{ʳ}{\ifmmode^{r}\else\textsuperscript{\ensuremath{r}}\fi}
\newunicodechar{ᵘ}{\ifmmode^{u}\else\textsuperscript{\ensuremath{u}}\fi}
\newunicodechar{ʸ}{\ifmmode^{y}\else\textsuperscript{\ensuremath{y}}\fi}
\newunicodechar{ᵃ}{\ifmmode^{a}\else\textsuperscript{\ensuremath{a}}\fi}
\newunicodechar{ᵉ}{\ifmmode^{e}\else\textsuperscript{\ensuremath{e}}\fi}
\newunicodechar{ᵏ}{\ifmmode^{k}\else\textsuperscript{\ensuremath{k}}\fi}
\newunicodechar{ᵖ}{\ifmmode^{p}\else\textsuperscript{\ensuremath{p}}\fi}
\newunicodechar{ᴺ}{\ifmmode^{N}\else\textsuperscript{\ensuremath{N}}\fi}
\newunicodechar{ᶜ}{\ifmmode^{c}\else\textsuperscript{\ensuremath{c}}\fi}
\newunicodechar{ʰ}{\ifmmode^{h}\else\textsuperscript{\ensuremath{h}}\fi}
\newunicodechar{ᵠ}{\ifmmode^{q}\else\textsuperscript{\ensuremath{q}}\fi}
\newunicodechar{ₙ}{\ifmmode_{n}\else\textsubscript{\ensuremath{n}}\fi}
\newunicodechar{ₐ}{\ifmmode_{a}\else\textsubscript{\ensuremath{a}}\fi}
\newunicodechar{ᵢ}{\ifmmode_{i}\else\textsubscript{\ensuremath{i}}\fi}
\newunicodechar{ᵣ}{\ifmmode_{r}\else\textsubscript{\ensuremath{r}}\fi}
\newunicodechar{ₖ}{\ifmmode_{k}\else\textsubscript{\ensuremath{k}}\fi}
\newunicodechar{ᵦ}{\ifmmode_{\beta}\else\textsubscript{\ensuremath{\beta}}\fi}
\newunicodechar{ₓ}{\ifmmode_{x}\else\textsubscript{\ensuremath{x}}\fi}
\newfontfamily\popdisplay{ArchivoBlack-Regular.ttf}[Path=../../../fonts/]
\newfontfamily\poplogo{SpaceGrotesk-Bold.ttf}[Path=../../../fonts/]
\newfontfamily\popsemi{PlusJakartaSans-SemiBold.ttf}[Path=../../../fonts/]
\newfontfamily\popxbold{PlusJakartaSans-ExtraBold.ttf}[Path=../../../fonts/]
\newfontfamily\popxboldls{PlusJakartaSans-ExtraBold.ttf}[Path=../../../fonts/,LetterSpace=3.0]

\setlist[enumerate]{leftmargin=1.7em,itemsep=5pt,topsep=4pt}
\setlist[itemize]{leftmargin=1.7em,itemsep=4pt,topsep=4pt,label={\color{poporange}\textbullet}}
\setlength{\parindent}{0pt}
\setlength{\parskip}{5pt}
\renewcommand{\arraystretch}{1.25}

% pgfplots : style Pop par défaut
\pgfplotsset{
  every axis/.append style={axis line style={popink,line width=1pt},tick style={popink},
    grid style={popline,line width=0.5pt},label style={font=\small},tick label style={font=\footnotesize}},
  popcurve/.style={poporange,line width=1.8pt,no markers,smooth},
}
% Même style disponible dans un \draw TikZ simple (hors pgfplots)
\tikzset{popcurve/.style={poporange,line width=1.8pt}}
\tikzset{popgrid/.style={popline,very thin}}
\colorlet{popcurve}{poporange} % le nom du style est parfois utilisé comme couleur

% ============================================================
% Pill / squiggle
% ============================================================
\newcommand{\poppill}[3]{%
  \tikz[baseline=-0.55ex]\node[draw=popink,line width=1.2pt,rounded corners=2.6pt,%
    fill=#2,inner xsep=6.5pt,inner ysep=3pt]{%
    \popxboldls\fontsize{7}{7}\selectfont\color{#3}\MakeUppercase{#1}};%
}
\newcommand{\popsquiggle}[1]{%
  \tikz[baseline=-0.4ex]{
    \draw[#1,line width=2.6pt,line cap=round]
      plot [smooth] coordinates {(0,0.09) (0.34,0.01) (0.68,0.21) (1.02,0.01) (1.36,0.10)};
  }%
}
% Mot-clé surligné (marqueur orange sous le texte)
\newcommand{\cle}[1]{{\popxbold\color{popdark}#1}}

% ============================================================
% En-tête / pied
% ============================================================
\pagestyle{fancy}
\fancyhf{}
\renewcommand{\headrulewidth}{0pt}
\renewcommand{\footrulewidth}{0pt}
\lhead{\raisebox{-0.15ex}{\poplogo\fontsize{13}{13}\selectfont\color{popink}OnBuch\color{poporange}+}}
\rhead{\poppill{\DOCMATIERE\ \textbullet\ \DOCNIVEAU\ \textbullet\ Leçon \DOCLECON}{white}{popink}}
\lfoot{\popxbold\fontsize{7.6}{7.6}\selectfont\color{popmuted}\MakeUppercase{Cours signé OnBuch+}\ \textbullet\ {\popsemi\DOCTITRE}}
\rfoot{\tikz[baseline=-0.6ex]\node[rounded corners=8.5pt,fill=popink,minimum height=17pt,minimum width=17pt,inner xsep=5pt,inner sep=0]{\popxbold\fontsize{8}{8}\selectfont\color{popcream}\thepage\,/\,\pageref*{LastPage}};}

% ============================================================
% Titres
% ============================================================
\newcommand{\chaptitre}[2]{%
  \begin{tcolorbox}[enhanced,colback=poporange,colframe=popink,boxrule=2.6pt,arc=11pt,
      drop shadow=popink,left=15pt,right=15pt,top=12pt,bottom=11pt,before skip=3pt,after skip=18pt]
    \poppill{#2}{popink}{popcream}\par\vspace{9pt}
    {\raggedright\hyphenpenalty=10000\exhyphenpenalty=10000\popdisplay\fontsize{22}{24}\selectfont\color{popcream}\MakeUppercase{#1}\par}
  \end{tcolorbox}
}
% Section numérotée : pastille orange avec le numéro + titre Archivo
\newcounter{popsec}
\newcommand{\coursec}[1]{%
  \stepcounter{popsec}\setcounter{popsub}{0}%
  \par\vspace{16pt}\noindent
  \begin{minipage}{\linewidth}
  \tikz[baseline=-0.7ex]\node[draw=popink,line width=1.6pt,fill=poporange,rounded corners=4pt,minimum size=22pt,inner sep=0]{\popdisplay\fontsize{12}{12}\selectfont\color{popcream}\thepopsec};%
  \hspace{8pt}{\popdisplay\fontsize{15}{17}\selectfont\color{popink}\MakeUppercase{#1}}\par
  \vspace{4pt}\noindent{\color{poporange}\rule{36pt}{3.6pt}}
  \end{minipage}\par\nopagebreak\vspace{8pt}
}
\newcounter{popsub}
\newcommand{\courssub}[1]{%
  \stepcounter{popsub}%
  \par\vspace{9pt}\noindent{\popxbold\fontsize{11.5}{13}\selectfont\color{popink}{\color{poporange}\thepopsec.\thepopsub}\hspace{6pt}#1}\par\nopagebreak\vspace{4pt}
}

% ============================================================
% Boîtes pédagogiques — même recette : fond blanc, bordure épaisse colorée,
% ombre dure encre, badge plein. Titre optionnel [#1].
% ============================================================
\tcbset{popbase/.style={breakable,enhanced,colback=white,boxrule=2.3pt,arc=9pt,
  shadow={3pt}{-3pt}{0pt}{popink},left=14pt,right=14pt,top=11pt,bottom=11pt,before skip=11pt,after skip=15pt,
  fontupper=\popsemi\fontsize{10.6}{14.5}\selectfont\color{popink},
  fontlower=\popsemi\fontsize{10.6}{14.5}\selectfont\color{popink}}}
% Coche dessinée (le glyphe ✓ n'existe pas dans les polices du document)
\newcommand{\popcheck}[1]{\tikz[baseline=-0.5ex]\draw[#1,line width=1.6pt,line cap=round,line join=round](-0.13,0.0)--(-0.04,-0.09)--(0.13,0.1);}
\providecommand{\coche}{\popcheck{popgreen}}
\providecommand{\resultat}[1]{\par\smallskip\noindent\fcolorbox{popgreen}{popgreenL}{\parbox{\dimexpr\linewidth-2\fboxsep-2\fboxrule\relax}{\textbf{Résultat.} #1}}\par\smallskip}
\newcommand{\popboxhead}[3]{\poppill{#1}{#2}{white}\ifx\relax#3\relax\else\hspace{6pt}{\popxbold\fontsize{10.6}{12}\selectfont\color{popink}#3}\fi\par\vspace{7pt}}

\newtcolorbox{aretenir}[1][]{popbase,colframe=popgreen,before upper={\popboxhead{À retenir}{popgreen}{#1}}}
% Texte en chinois (document de lecture, dialogue, extrait) : cadre
% d'étude. Un titre optionnel (source, auteur) s'affiche dans l'en-tête.
\newtcolorbox{textebox}[1][]{popbase,colframe=popblue,before upper={\popboxhead{文本 · Texte}{popblue}{#1}},after upper={}}
% Marques ✓ / ✗ (forme correcte / fautive) souvent utilisées par les modèles
\providecommand{\cross}{\textcolor{poppink}{\ensuremath{\times}}}
\providecommand{\xmark}{\cross}\providecommand{\crossmark}{\cross}\providecommand{\cmark}{\ensuremath{\checkmark}}
% Ligne de réponse / justification à compléter (inventée par les modèles)
\providecommand{\justification}[1]{\par\noindent\textit{Justification~:} \dotfill\par}
% \textipa (package tipa non chargé) : la police ne couvre pas l'API -> simple passage
\providecommand{\textipa}[1]{#1}
% Soulignement (ulem non chargé) : \uline, \ul
\providecommand{\uline}[1]{\underline{#1}}\providecommand{\ul}[1]{\underline{#1}}
% \glossaire{\item ...} inventée par les modèles : liste de glossaire
\providecommand{\glossaire}[1]{\begin{itemize}#1\end{itemize}}
% \alert (beamer) inventée par les modèles : texte mis en évidence
\providecommand{\alert}[1]{\textbf{\textcolor{poppink}{#1}}}
% variantes ASCII d'ordinaux écrites en macro
\providecommand{\iere}{\textsuperscript{re}}\providecommand{\ieme}{\textsuperscript{e}}\providecommand{\ere}{\textsuperscript{re}}\providecommand{\eme}{\textsuperscript{e}}\providecommand{\er}{\textsuperscript{er}}\providecommand{\ie}{\textsuperscript{e}}
% Ordinaux français écrits en macro par les modèles : 1\ière, 3\ième
\newcommand{\ière}{\textsuperscript{re}}\newcommand{\ième}{\textsuperscript{e}}
% \exemple (inventée par les modèles) : étiquette « Exemple : »
\providecommand{\exemple}{\textbf{Exemple~:}\ }
% \square utilisable aussi hors mode math (cases à cocher)
\AtBeginDocument{\let\squareMath\square\renewcommand{\square}{\ensuremath{\squareMath}}}
% Mot ou phrase en chinois à l'intérieur d'un paragraphe français
\newcommand{\zh}[1]{\ifmmode\text{#1}\else{#1}\fi}
\let\eng\zh \let\esp\zh \let\all\zh % alias tolérés
% flèches d'intonation inventées par les modèles
\providecommand{\textsearrow}{}\renewcommand{\textsearrow}{\ensuremath{\searrow}}\providecommand{\textnearrow}{}\renewcommand{\textnearrow}{\ensuremath{\nearrow}}
% \fr{..} : mot français (inventé par les modèles)
\providecommand{\fr}[1]{\textit{#1}}
% zhbox : encadré de texte chinois inventé par les modèles
\newenvironment{zhbox}{\begin{quote}}{\end{quote}}
% \courssubsub : titre de niveau 3 inventé par les modèles
\providecommand{\courssubsub}[1]{\par\medskip\noindent\textbf{#1}\par\smallskip}
% \zhbf : texte chinois en gras inventé par les modèles
\providecommand{\zhbf}[1]{\textbf{#1}}
% \vs : « versus » inventé par les modèles
\providecommand{\vs}{\textit{vs}\ }
% \blank : case à compléter inventée par les modèles
\providecommand{\blank}[1][]{\underline{\hspace{1.6em}}}
\newtcolorbox{exemplebox}[1][]{popbase,colframe=poporange,before upper={\popboxhead{Exemple}{poporange}{#1}}}
\newtcolorbox{definition}[1][]{popbase,colframe=popblue,before upper={\popboxhead{Définition}{popblue}{#1}}}
\newtcolorbox{propriete}[1][]{popbase,colframe=poppurple,before upper={\popboxhead{Propriété}{poppurple}{#1}}}
\newtcolorbox{methode}[1][]{popbase,colframe=popgold,before upper={\popboxhead{Méthode}{popgold}{#1}}}
\newtcolorbox{attention}[1][]{popbase,colframe=poppink,before upper={\popboxhead{Attention}{poppink}{#1}}}
\newtcolorbox{experience}[1][]{popbase,colframe=popblue,colback=popblueL,before upper={\popboxhead{Expérience}{popblue}{#1}}}
% « Le savais-tu ? » : carte violette pleine (contexte, culture, Cameroun)
\newtcolorbox{savaistu}[1][]{popbase,colback=poppurple,colframe=popink,
  fontupper=\popsemi\fontsize{10.4}{14.2}\selectfont\color{white},
  before upper={\poppill{Le savais-tu\,?}{white}{poppurple}\ifx\relax#1\relax\else\hspace{6pt}{\popxbold\color{white}#1}\fi\par\vspace{7pt}}}
% Exercice résolu : énoncé en haut, \tcblower puis la solution sur fond crème
\newcounter{popexo}\newcounter{popexr}
\newtcolorbox{exoresolu}[1][]{popbase,colframe=poporange,colbacklower=popcream,
  segmentation style={popink,line width=1.2pt,dashed},
  before upper={\stepcounter{popexr}\popboxhead{Exercice résolu \thepopexr}{poporange}{#1}},
  before lower={\poppill{Solution}{popink}{popcream}\par\vspace{6pt}}}
% Exercice (non résolu en place) — la correction est dans la partie Corrigés
\newtcolorbox{exercice}[1][]{popbase,colframe=popink,
  before upper={\stepcounter{popexo}\popboxhead{Exercice \thepopexo}{popink}{#1}}}
% Corrigé
\newtcolorbox{corrige}[1][]{popbase,colframe=popgreen,colback=popgreenL,
  before upper={\popboxhead{Corrigé}{popgreen}{#1}}}
% Objectifs / prérequis / bilan
\newtcolorbox{objectifs}{popbase,colframe=popink,colback=poporangeL,
  before upper={\popboxhead{Objectifs de la leçon}{popink}{}}}
\newtcolorbox{prerequis}{popbase,colframe=popink,colback=popblueL,
  before upper={\popboxhead{Prérequis}{popblue}{}}}
\newtcolorbox{bilan}{popbase,colframe=popink,colback=white,boxrule=2.8pt,
  before upper={\popboxhead{Fiche bilan}{poporange}{L'essentiel à connaître pour l'examen}}}
% Figure encadrée avec légende
\newtcolorbox{popfigure}[1][]{enhanced,breakable=false,colback=white,colframe=popline,boxrule=1.4pt,arc=8pt,
  left=10pt,right=10pt,top=10pt,bottom=8pt,before skip=10pt,after skip=12pt,halign=center,
  lower separated=false,halign lower=center,
  fontlower=\popsemi\fontsize{9}{11.5}\selectfont\color{popmuted},
  #1}
\newcommand{\legende}[1]{\par\vspace{6pt}{\popsemi\fontsize{9}{11.5}\selectfont\color{popmuted}#1}}

% ============================================================
% Signature OnBuch+
% ============================================================
\newcommand{\poplogoinline}[1]{{\poplogo\fontsize{#1}{#1}\selectfont\color{popink}OnBuch\color{poporange}+}}
\newcommand{\signatureonbuch}{%
  \begin{tcolorbox}[enhanced,colback=popink,colframe=popink,boxrule=2.2pt,arc=10pt,
      drop shadow=poporange,left=14pt,right=14pt,top=10pt,bottom=10pt,before skip=0pt,after skip=0pt]
    \tikz[baseline=-0.7ex]\node[circle,fill=popgreen,draw=popcream,line width=1.3pt,minimum size=22pt,inner sep=0]{\popcheck{white}};%
    \hspace{9pt}\begin{minipage}[c]{0.62\linewidth}
      {\popxbold\fontsize{12}{13}\selectfont\color{popcream}Signé\ \textbullet\ {\poplogo OnBuch\color{poporange}+}}\par\vspace{2pt}
      {\raggedright\popsemi\fontsize{8}{10.5}\selectfont\color{popline}\MakeUppercase{Équipe pédagogique OnBuch+}\\\MakeUppercase{Conforme au programme officiel MINESEC}\par}
    \end{minipage}\hfill
    {\poplogo\fontsize{22}{22}\selectfont\color{popcream}OnBuch\color{poporange}+}
  \end{tcolorbox}%
}

% ============================================================
% Page de garde
% #1 titre · #2 matière · #3 niveau · #4 module · #5 n° leçon · #6 durée ·
% #7 description / objectifs (texte ou itemize)
% ============================================================
\newcommand{\pagegarde}[7]{%
  \thispagestyle{empty}
  \newgeometry{a4paper,margin=1.7cm,top=1.6cm,bottom=1.4cm}%
  \begin{tikzpicture}[remember picture,overlay]
    \fill[poporange!18] ([xshift=-3.2cm,yshift=-3.6cm]current page.north east) circle (4.6cm);
    \fill[poppurple!14] ([xshift=2.4cm,yshift=7.6cm]current page.south west) circle (3.8cm);
  \end{tikzpicture}%
  {\poplogo\fontsize{30}{30}\selectfont\color{popink}OnBuch\color{poporange}+}\hfill
  \raisebox{0.6ex}{\poppill{Cours signé \textbullet\ Premium}{popink}{popcream}}\par
  \vspace{1pt}\popsquiggle{poppurple}\par
  \vspace{8pt}
  \begin{tcolorbox}[enhanced,colback=poporange,colframe=popink,boxrule=2.8pt,arc=13pt,
      drop shadow=popink,left=18pt,right=18pt,top=12pt,bottom=13pt,after skip=10pt]
    \poppill{#2 \textbullet\ #3}{popink}{popcream}\hspace{5pt}\poppill{#4}{white}{popink}\par\vspace{8pt}
    {\popdisplay\fontsize{14}{14}\selectfont\color{popink}LEÇON #5}\par\vspace{4pt}
    {\raggedright\hyphenpenalty=10000\exhyphenpenalty=10000\popdisplay\fontsize{25}{27}\selectfont\color{popcream}\MakeUppercase{#1}\par}
    \vspace{8pt}
    {\popxbold\fontsize{9.5}{9.5}\selectfont\color{popink}\MakeUppercase{Durée conseillée : #6}}
  \end{tcolorbox}
  \begin{tcolorbox}[enhanced,colback=white,colframe=popink,boxrule=2.2pt,arc=10pt,
      drop shadow=popink,left=16pt,right=16pt,top=10pt,bottom=10pt,after skip=8pt]
    \poppill{Dans cette leçon}{poppurple}{white}\par\vspace{5pt}
    \popsemi\fontsize{9.3}{12.6}\selectfont\color{popink}#7
  \end{tcolorbox}
  \noindent\begin{minipage}[t]{0.487\linewidth}
    \begin{tcolorbox}[enhanced,colback=popgreen,colframe=popink,boxrule=2.2pt,arc=10pt,
        drop shadow=popink,left=11pt,right=11pt,top=10pt,bottom=10pt,height=3.1cm,valign=center]
      \tcbox[colback=white,colframe=popink,boxrule=1.3pt,arc=5pt,boxsep=0pt,nobeforeafter,
        left=3pt,right=3pt,top=3pt,bottom=3pt,valign=center]{\qrcode[height=1.9cm]{https://play.google.com/store/apps/details?id=com.luvvix.onbuch&hl=fr}}%
      \hspace{8pt}\begin{minipage}[c]{0.5\linewidth}
        {\popdisplay\fontsize{11}{12.5}\selectfont\color{white}\MakeUppercase{Télécharge OnBuch}}\par\vspace{4pt}
        {\raggedright\popsemi\fontsize{8.4}{11}\selectfont\color{white}Gratuit \textbullet\ Google Play\\Tuteur IA, annales, concours, résultats\par}
      \end{minipage}
    \end{tcolorbox}
  \end{minipage}\hfill
  \begin{minipage}[t]{0.487\linewidth}
    \begin{tcolorbox}[enhanced,colback=poppurple,colframe=popink,boxrule=2.2pt,arc=10pt,
        drop shadow=popink,left=11pt,right=11pt,top=10pt,bottom=10pt,height=3.1cm,valign=center]
      \tikz[baseline=-0.7ex]\node[circle,fill=white,draw=popink,line width=1.3pt,minimum size=1.6cm,inner sep=0]{\popdisplay\fontsize{24}{24}\selectfont\color{poppurple}+};%
      \hspace{8pt}\begin{minipage}[c]{0.6\linewidth}
        {\popdisplay\fontsize{11}{12.5}\selectfont\color{white}\MakeUppercase{Passe à OnBuch+}}\par\vspace{4pt}
        {\raggedright\popsemi\fontsize{8.4}{11}\selectfont\color{white}Tous les cours signés, Tuteur IA illimité, fiches PDF — onglet OnBuch+ de l'app\par}
      \end{minipage}
    \end{tcolorbox}
  \end{minipage}
  \vspace{8pt}
  \popcontenu
  \vfill
  \signatureonbuch
  \restoregeometry
  \newpage
}

% Rangée « Ce cours contient » (page de garde)
\newcommand{\poptile}[3]{% #1 couleur #2 gros texte #3 légende
  \begin{tcolorbox}[enhanced,colback=#1,colframe=popink,boxrule=2pt,arc=9pt,shadow={2.5pt}{-2.5pt}{0pt}{popink},
      left=6pt,right=6pt,top=9pt,bottom=9pt,halign=center,valign=center,height=1.75cm,width=\linewidth,nobeforeafter]
    {\popdisplay\fontsize{12.5}{13}\selectfont\color{white}\MakeUppercase{#2}}\par\vspace{3pt}
    {\centering\popsemi\fontsize{7.8}{9.5}\selectfont\color{white}#3\par}
  \end{tcolorbox}}
\newcommand{\popcontenu}{%
  \par\noindent{\popxboldls\fontsize{7.5}{7.5}\selectfont\color{popmuted}CE COURS SIGNÉ CONTIENT}\par\vspace{7pt}
  \noindent\begin{minipage}[t]{0.235\linewidth}\poptile{popblue}{Cours}{complet, illustré, pas à pas}\end{minipage}\hfill
  \begin{minipage}[t]{0.235\linewidth}\poptile{poporange}{Méthodes}{et exercices résolus}\end{minipage}\hfill
  \begin{minipage}[t]{0.235\linewidth}\poptile{poppink}{Exercices}{type examen + corrigés}\end{minipage}\hfill
  \begin{minipage}[t]{0.235\linewidth}\poptile{popgreen}{Fiche bilan}{l'essentiel à retenir}\end{minipage}\par}

% Sommaire
\newtcolorbox{sommaire}{enhanced,colback=white,colframe=popink,boxrule=2.2pt,arc=10pt,
  drop shadow=popink,left=18pt,right=18pt,top=13pt,bottom=13pt,before skip=6pt,after skip=20pt,
  before upper={\poppill{Sommaire}{poppurple}{white}\par\vspace{9pt}\popsemi\fontsize{10.6}{14.5}\selectfont\color{popink}}}

% Signature de fin de cours — poussée tout en bas de la dernière page
\newcommand{\findecours}{%
  \par\vfill\nopagebreak
  \begin{center}\popsquiggle{poporange}\end{center}\vspace{4pt}
  \signatureonbuch
}
\AtBeginDocument{\def\llbracket{\mathopen{[\mkern-3.5mu[}}\def\rrbracket{\mathclose{]\mkern-3.5mu]}}} % intervalles d'entiers (glyphe ⟦ absent de la police)
% Glyphes absents de Fira Math : redéfinis à partir de glyphes disponibles
\AtBeginDocument{%
  \def\setminus{\mathbin{\backslash}}\let\smallsetminus\setminus
  \def\vdots{\mathord{\vbox{\baselineskip3.2pt\lineskiplimit0pt\kern4pt\hbox{.}\hbox{.}\hbox{.}}}}%
  \def\ddots{\mathinner{\mkern1mu\raise6.5pt\hbox{.}\mkern2mu\raise3.6pt\hbox{.}\mkern2mu\raise0.7pt\hbox{.}\mkern1mu}}%
  \def\triangle{\mathord{\scalebox{0.9}{$\symup{\Delta}$}}}%
  \def\nabla{\mathord{\raisebox{\depth}{\scalebox{1}[-1]{$\symup{\Delta}$}}}}%
  \def\checkmark{\popcheck{popdark}}%
  \def\star{\mathbin{\ast}}\def\bigstar{\mathord{\ast}}%
  \def\diamond{\mathbin{\scalebox{0.6}{$\Diamond$}}}%
  \def\bigcup{\mathop{\mathchoice{\raisebox{-0.3em}{\scalebox{1.7}{$\cup$}}}{\scalebox{1.25}{$\cup$}}{\cup}{\cup}}\displaylimits}%
  \def\bigcap{\mathop{\mathchoice{\raisebox{-0.3em}{\scalebox{1.7}{$\cap$}}}{\scalebox{1.25}{$\cap$}}{\cap}{\cap}}\displaylimits}%
  \def\lesseqqgtr{\mathrel{\lessgtr}}\def\bigoplus{\mathop{\oplus}\displaylimits}%
}
\providecommand{\ssi}{si et seulement si} % abréviation fréquente
\AtBeginDocument{\providecommand{\wideparen}{\overparen}} % arc de cercle

%% FIGFIX-BEGIN v3 — correctifs globaux des figures (docs/onbuch-generation/tools/figfix)
\usepackage{adjustbox}\usepackage{etoolbox}
% toute figure TikZ/pgfplots plus large que la ligne est réduite à la largeur disponible
\BeforeBeginEnvironment{tikzpicture}{\begin{adjustbox}{max width=\linewidth}}
\AfterEndEnvironment{tikzpicture}{\end{adjustbox}}
% légendes pgfplots lisibles par-dessus les courbes
\pgfplotsset{every axis/.append style={legend style={fill=white,fill opacity=0.92,text opacity=1,draw=black!15,inner sep=3pt}}}
% étiquettes d'arêtes (syntaxe edge["..."]) avec fond blanc
\tikzset{every edge quotes/.append style={fill=white,inner sep=1.2pt}}
% bulles de cartes mentales : texte à la ligne dans la bulle, bulle un peu plus grande
\tikzset{every concept/.append style={text width=2.0cm,align=center,minimum size=2.3cm,inner sep=2pt}}
%% FIGFIX-END
\begin{document}

\pagegarde{Leçon 39 — Grammaire : verbe modal 可以 ; 的 、 得、 地}{Chinois}{Tle A}{Module 5 : Mass médias et télécommunication}{ZH39}{2 h}{Cette leçon de grammaire présente le verbe modal 可以 (pouvoir/permission) et les trois particules homophones 的、得、地, avec leurs schémas, emplois, exceptions, comparaisons au français et exercices d'application.
\vspace{4pt}
\begin{multicols}{2}\small
\begin{itemize}
\item À la fin de cette leçon, je sais identifier et utiliser le verbe modal 可以 pour exprimer la permission et la possibilité.
\item Je sais distinguer les trois particules 的、得、地 selon leur fonction syntaxique (attribut, complément de degré, adverbial).
\item Je construis correctement les structures [Sujet + 可以 + Verbe + Objet] et [Verbe/Adjectif + 得 + Complément de degré].
\item J'évite les erreurs fréquentes des francophones (confusion 的/得/地, ordre des mots après 可以).
\item Je transforme des phrases affirmatives en interrogatives et négatives avec 可以.
\item J'applique les particules dans des phrases complexes et de courts textes.
\end{itemize}\end{multicols}}

\begin{sommaire}
\begin{multicols}{2}
\begin{enumerate}
\item 1. Observation d'exemples
\item 2. Schéma de la structure
\item 3. Emplois et exceptions
\item 4. Comparaison avec le français et pièges des francophones
\item 5. Exercices d'application et de transformation
\item 6. Compléments utiles à l'examen (bonus)
\item Activité d'intégration
\item Méthodes et exercices résolus
\item Exercices
\item Corrigés des exercices
\item Fiche bilan
\end{enumerate}
\end{multicols}
\end{sommaire}

\begin{objectifs}
\begin{itemize}
\item À la fin de cette leçon, je sais identifier et utiliser le verbe modal 可以 pour exprimer la permission et la possibilité.
\item Je sais distinguer les trois particules 的、得、地 selon leur fonction syntaxique (attribut, complément de degré, adverbial).
\item Je construis correctement les structures [Sujet + 可以 + Verbe + Objet] et [Verbe/Adjectif + 得 + Complément de degré].
\item J'évite les erreurs fréquentes des francophones (confusion 的/得/地, ordre des mots après 可以).
\item Je transforme des phrases affirmatives en interrogatives et négatives avec 可以.
\item J'applique les particules dans des phrases complexes et de courts textes.
\item Je compare les équivalents français (pouvoir, que, -ment) et justifie mes choix grammaticaux.
\item Je réalise des exercices de type Bac (complétion, transformation, traduction).
\item Je produis un court dialogue oral intégrant 可以 et les trois particules.
\end{itemize}
\end{objectifs}

\begin{prerequis}
\begin{itemize}
\item Connaissance des pronoms personnels et de l'ordre SVO de base (Première).
\item Maîtrise des verbes d'action simples (去, 吃, 看) et des adjectifs courants (好, 快, 慢).
\item Notion de base sur les particules de structure (的 comme possessif) abordée en Première.
\item Vocabulaire de la vie quotidienne (boissons, déplacements, devoirs).
\item Capacité à lire le pinyin avec les quatre tons.
\end{itemize}
\end{prerequis}

\newpage

\coursec{Observation d'exemples}

\courssub{Situation d'accroche et problématique}

Imaginez la scène : il est midi, le soleil tape fort sur le marché Mokolo à Yaoundé. Un vendeur de jus, un peu fatigué, s'approche d'une étudiante chinoise en visite et lui demande poliment :

\begin{textebox}[Au marché de Yaoundé]
可以给我一瓶水吗？\\
\emph{Kěyǐ gěi wǒ yī píng shuǐ ma ?}\\
« Puis-je avoir une bouteille d'eau ? »
\end{textebox}

En classe de Terminale, les élèves comparent cette demande polie à la tournure française « Est-ce que je peux\ldots ? ». Ils remarquent aussitôt que le mot \zh{可以} occupe exactement la place et la fonction de « pouvoir » en français : c'est un \cle{verbe modal}. Mais très vite, une autre difficulté apparaît. En lisant leurs cahiers, les élèves croisent trois petits caractères qui se prononcent exactement pareil — \emph{de}, au ton neutre — mais qui s'écrivent différemment et ne se placent pas au même endroit : \zh{的}, \zh{得} et \zh{地}. Le français n'a qu'un seul « de » ; le chinois en a trois, et les confondre est la faute la plus fréquente des francophones.

\textbf{Problématique :} comment choisir correctement entre \zh{的}, \zh{得} et \zh{地}, et comment employer le verbe modal \zh{可以} pour exprimer la permission, la possibilité ou la suggestion ? Pour répondre, commençons par observer des phrases modèles, sans rien expliquer encore : regardons, soulignons, comparons.

\begin{definition}[Le verbe modal 可以]
\zh{可以} (\emph{kěyǐ}) est un \cle{verbe modal} qui se place \textbf{avant le verbe principal} et exprime la \textbf{permission} (« il est permis de »), la \textbf{possibilité} (« il est possible de ») ou la \textbf{suggestion} (« on pourrait »). Il correspond au français « pouvoir ».
\end{definition}

\courssub{Observation 1 : cinq phrases modèles avec 可以}

Lisons d'abord ces cinq phrases. Dans chacune, repérez la place de \zh{可以} : toujours \textbf{entre le sujet et le verbe}.

\begin{exemplebox}[可以 : permission, possibilité, suggestion]
\begin{enumerate}
\item \zh{我可以去北京。} \emph{Wǒ kěyǐ qù Běijīng.} « Je peux aller à Pékin. » (possibilité)
\item \zh{你可以吃苹果。} \emph{Nǐ kěyǐ chī píngguǒ.} « Tu peux manger une pomme. » (permission)
\item \zh{我们可以在这里拍照吗？} \emph{Wǒmen kěyǐ zài zhèlǐ pāizhào ma ?} « Pouvons-nous prendre des photos ici ? » (demande de permission)
\item \zh{今天很热，你可以喝一点水。} \emph{Jīntiān hěn rè, nǐ kěyǐ hē yìdiǎn shuǐ.} « Il fait chaud aujourd'hui, tu peux boire un peu d'eau. » (suggestion)
\item \zh{学生不可以在教室里打电话。} \emph{Xuésheng bù kěyǐ zài jiàoshì lǐ dǎ diànhuà.} « Les élèves ne peuvent pas téléphoner dans la salle de classe. » (interdiction, négation \zh{不可以})
\end{enumerate}
\end{exemplebox}

\textbf{Ce que nous observons :}
\begin{itemize}
\item \zh{可以} ne change jamais de forme : pas de conjugaison, pas d'accord. C'est beaucoup plus simple que le français « peux / peut / pouvons ».
\item La négation se fait avec \zh{不} placé \textbf{avant} \zh{可以} : \zh{不可以} = « ne pas pouvoir, ne pas avoir le droit ».
\item Pour poser une question, on ajoute simplement \zh{吗} à la fin : \zh{可以吗？} = « Est-ce possible ? / Puis-je ? ».
\item Schéma observé : \textbf{Sujet + 可以 + Verbe (+ complément)}.
\end{itemize}

\courssub{Observation 2 : cinq phrases modèles avec 的}

\begin{exemplebox}[的 : attribut du nom et possession]
\begin{enumerate}
\item \zh{这是我的书。} \emph{Zhè shì wǒ de shū.} « C'est mon livre. » (possession)
\item \zh{红色的花很漂亮。} \emph{Hóngsè de huā hěn piàoliang.} « Les fleurs rouges sont très belles. » (adjectif + nom)
\item \zh{喀麦隆的足球很有名。} \emph{Kāmàilóng de zúqiú hěn yǒumíng.} « Le football du Cameroun est très célèbre. » (nom + nom)
\item \zh{我喜欢说中文的老师。} \emph{Wǒ xǐhuan shuō Zhōngwén de lǎoshī.} « J'aime le professeur qui parle chinois. » (proposition + nom)
\item \zh{她穿了一件白色的衣服。} \emph{Tā chuān le yí jiàn báisè de yīfu.} « Elle a mis un vêtement blanc. » (adjectif + nom)
\end{enumerate}
\end{exemplebox}

\begin{propriete}[的 : la particule de l'attribut nominal]
\zh{的} (\emph{de}, ton neutre) se place \textbf{après} un adjectif, un nom, un pronom ou une proposition, et \textbf{avant} le nom qu'il qualifie. Schéma : \textbf{Qualifiant + 的 + Nom}. Il correspond au français « de » (\emph{le livre de Marie}), à l'adjectif possessif (\emph{mon, ton}) ou à la relation adjectivale (\emph{une fleur rouge}).
\end{propriete}

\textbf{Observation essentielle :} en français, l'adjectif se place souvent \emph{après} le nom (\emph{une fleur rouge}) ; en chinois, le qualifiant précède \textbf{toujours} le nom, relié par \zh{的} : \zh{红色的花} = littéralement « rouge-de fleur ».

\courssub{Observation 3 : cinq phrases modèles avec 得}

\begin{exemplebox}[得 : complément de degré après le verbe]
\begin{enumerate}
\item \zh{他跑得很快。} \emph{Tā pǎo de hěn kuài.} « Il court très vite. »
\item \zh{她唱得很好。} \emph{Tā chàng de hěn hǎo.} « Elle chante bien. »
\item \zh{他汉语说得不错。} \emph{Tā Hànyǔ shuō de búcuò.} « Il parle chinois pas mal. »
\item \zh{弟弟高兴得跳了起来。} \emph{Dìdi gāoxìng de tiào le qǐlái.} « Le petit frère était si content qu'il a sauté. »
\item \zh{我累得不想说话。} \emph{Wǒ lèi de bù xiǎng shuōhuà.} « Je suis si fatigué que je n'ai pas envie de parler. »
\end{enumerate}
\end{exemplebox}

\textbf{Ce que nous observons :} \zh{得} se place \textbf{après un verbe ou un adjectif} et introduit un \cle{complément de degré} ou de manière qui évalue l'action : \textbf{Verbe + 得 + Adjectif}. Le français rend cela par un adverbe (« il court \emph{vite} ») ou par « si\ldots que » (« si fatigué \emph{que}\ldots »).

\courssub{Observation 4 : cinq phrases modèles avec 地}

\begin{exemplebox}[地 : l'adjectif devient adverbe]
\begin{enumerate}
\item \zh{她快乐地跳舞。} \emph{Tā kuàilè de tiàowǔ.} « Elle danse joyeusement. »
\item \zh{老师慢慢地写汉字。} \emph{Lǎoshī mànmàn de xiě Hànzì.} « Le professeur écrit lentement les caractères. »
\item \zh{他认真地学习。} \emph{Tā rènzhēn de xuéxí.} « Il étudie sérieusement. »
\item \zh{孩子们高兴地唱歌。} \emph{Háizimen gāoxìng de chànggē.} « Les enfants chantent joyeusement. »
\item \zh{请安静地听。} \emph{Qǐng ānjìng de tīng.} « Écoutez calmement, s'il vous plaît. »
\end{enumerate}
\end{exemplebox}

\textbf{Ce que nous observons :} \zh{地} se place \textbf{entre un adjectif et un verbe} : l'adjectif devient un complément circonstanciel de manière, comme le suffixe français « -ment » (\emph{lent} → \emph{lentement}). Schéma : \textbf{Adjectif + 地 + Verbe}.

\courssub{Analyse comparative : la même phrase avec 的、得、地}

Pour bien sentir la différence, observons le même matériau lexical (\zh{快} « rapide », \zh{跑} « courir ») dans trois constructions :

\begin{center}
\begin{tabularx}{\linewidth}{|X|X|X|}
\hline
\rowcolor{popblueL} Phrase & Structure & Traduction \\
\hline
\zh{快的马跑得快。} \emph{Kuài de mǎ pǎo de kuài.} & 的 avant un nom ; 得 après un verbe & « Un cheval rapide court vite. » \\
\hline
\zh{他跑得很快。} \emph{Tā pǎo de hěn kuài.} & Verbe + 得 + adjectif (degré) & « Il court très vite. » \\
\hline
\zh{他很快地跑了。} \emph{Tā hěn kuài de pǎo le.} & Adjectif + 地 + verbe (manière) & « Il a couru rapidement. » \\
\hline
\end{tabularx}
\end{center}

Autre triplet, avec \zh{快乐} « joyeux » et \zh{唱} « chanter » :

\begin{center}
\begin{tabularx}{\linewidth}{|X|X|X|}
\hline
\rowcolor{popblueL} Phrase & Particule & Traduction \\
\hline
\zh{快乐的孩子们} \emph{kuàilè de háizimen} & 的 (attribut du nom) & « les enfants joyeux » \\
\hline
\zh{孩子们唱得很快乐。} \emph{Háizimen chàng de hěn kuàilè.} & 得 (degré du verbe) & « Les enfants chantent avec joie. » \\
\hline
\zh{孩子们快乐地唱歌。} \emph{Háizimen kuàilè de chànggē.} & 地 (manière) & « Les enfants chantent joyeusement. » \\
\hline
\end{tabularx}
\end{center}

\textbf{Conclusion de l'observation :} le sens global reste proche, mais la \textbf{place} de la particule change la \textbf{fonction grammaticale} : \zh{的} construit un groupe nominal, \zh{得} évalue une action après le verbe, \zh{地} qualifie la manière avant le verbe.

\begin{methode}[Comment choisir entre 的、得、地 ?]
\begin{enumerate}
\item Repérer le mot qui \textbf{précède} et celui qui \textbf{suit} la particule.
\item Si on a \textbf{qualifiant + particule + nom} → c'est \zh{的} : \zh{我的书} « mon livre ».
\item Si on a \textbf{verbe (ou adjectif) + particule + adjectif de degré} → c'est \zh{得} : \zh{跑得快} « courir vite ».
\item Si on a \textbf{adjectif + particule + verbe} → c'est \zh{地} : \zh{快乐地唱} « chanter joyeusement ».
\item Vérifier en traduisant : « de / possessif » → \zh{的} ; « si\ldots que, bien, vite » après le verbe → \zh{得} ; « -ment » avant le verbe → \zh{地}.
\end{enumerate}
\end{methode}

\begin{popfigure}
\begin{tikzpicture}[node distance=0.9cm]
\node[draw=popblue, fill=popblueL, rounded corners, thick, align=center] (ex) {Exemples};
\node[below of=ex, draw=poporange, fill=poporangeL, rounded corners, align=center] (k) {可以 : 我可以去。};
\node[below of=k, draw=popgreen, fill=popgreenL, rounded corners, align=center] (de1) {的 : 我的书。};
\node[below of=de1, draw=poppurple, fill=poppurpleL, rounded corners, align=center] (de2) {得 : 他跑得快。};
\node[below of=de2, draw=poppink, fill=poppinkL, rounded corners, align=center] (de3) {地 : 他快乐地唱。};
\draw[->, thick, popdark] (ex) -- (k);
\draw[->, thick, popdark] (k) -- (de1);
\draw[->, thick, popdark] (de1) -- (de2);
\draw[->, thick, popdark] (de2) -- (de3);
\end{tikzpicture}
\legende{Carte mentale d'observation : un verbe modal (可以) et les trois particules \emph{de}.}
\end{popfigure}

\courssub{Tableau récapitulatif d'observation}

\begin{center}
\begin{tabularx}{\linewidth}{|X|X|X|X|}
\hline
\rowcolor{popblueL} Particule & Place & Schéma & Équivalent français \\
\hline
\zh{可以} \emph{kěyǐ} & Avant le verbe principal & Sujet + 可以 + Verbe & « pouvoir » (permission, possibilité) \\
\hline
\zh{的} \emph{de} & Entre le qualifiant et le nom & Qualifiant + 的 + Nom & « de », possessif, adjectif \\
\hline
\zh{得} \emph{de} & Après le verbe ou l'adjectif & Verbe + 得 + Adjectif & adverbe, « si\ldots que » \\
\hline
\zh{地} \emph{de} & Entre l'adjectif et le verbe & Adjectif + 地 + Verbe & adverbe en « -ment » \\
\hline
\end{tabularx}
\end{center}

\courssub{Exercice résolu d'observation}

\begin{exoresolu}[Identifier la bonne particule]
\textbf{Énoncé.} Dans chaque phrase, indiquez si le caractère souligné est \zh{的}, \zh{得} ou \zh{地}, et justifiez : (1) \zh{这是老师de书。} (2) \zh{他说de很好。} (3) \zh{她高兴de笑了。} (4) \zh{我可以进来吗？}
\tcblower
\textbf{Solution détaillée.}
\begin{enumerate}
\item \zh{这是老师的书。} \emph{Zhè shì lǎoshī de shū.} « C'est le livre du professeur. » → \zh{的}, car on a \textbf{nom (老师) + particule + nom (书)} : possession.
\item \zh{他说得很好。} \emph{Tā shuō de hěn hǎo.} « Il parle très bien. » → \zh{得}, car on a \textbf{verbe (说) + particule + adjectif de degré (很好)}.
\item \zh{她高兴地笑了。} \emph{Tā gāoxìng de xiào le.} « Elle a souri joyeusement. » → \zh{地}, car on a \textbf{adjectif (高兴) + particule + verbe (笑)} : manière.
\item \zh{我可以进来吗？} \emph{Wǒ kěyǐ jìnlái ma ?} « Puis-je entrer ? » → \zh{可以} est un verbe modal placé avant le verbe \zh{进来} ; il exprime ici une demande de permission.
\end{enumerate}
\end{exoresolu}

\begin{attention}[L'erreur n° 1 des francophones]
Ne confondez jamais \zh{的} et \zh{得} après un verbe. La phrase *\zh{他跑的快} est \textbf{fausse} : après le verbe \zh{跑}, pour évaluer l'action, il faut \zh{得} : \zh{他跑得快} « Il court vite ». Astuce : \zh{的} regarde vers le \textbf{nom qui suit} ; \zh{得} regarde vers le \textbf{verbe qui précède} ; \zh{地} regarde vers le \textbf{verbe qui suit}. Autre piège : ne traduisez pas « pouvoir » par \zh{可以} quand il s'agit d'une capacité apprise (savoir nager, parler chinois) : on utilise alors \zh{会} (\emph{huì}), que nous étudierons en section 3.
\end{attention}

\begin{savaistu}[D'où viennent ces particules ?]
En chinois classique, \zh{的} n'existait pas : pour marquer l'attribut et la possession, on utilisait le caractère \zh{之} (\emph{zhī}), que l'on retrouve encore aujourd'hui dans les expressions figées et les titres élégants, par exemple \zh{友谊之光} (\emph{yǒuyì zhī guāng}, « la lumière de l'amitié »). Quant aux trois \emph{de} modernes, ils se prononcent tous \emph{de} (ton neutre) mais s'écrivent avec des clés différentes : \zh{的} contient la clé 白 « blanc », \zh{得} contient la clé 彳 « pas, marcher », et \zh{地} contient la clé 土 « terre ». À l'oral, on ne les distingue pas : seule l'écriture — et la grammaire — permet de choisir le bon caractère. C'est pourquoi les élèves chinois eux-mêmes s'entraînent longtemps à les différencier !
\end{savaistu}

\begin{aretenir}[Ce qu'il faut retenir de cette observation]
\begin{itemize}
\item \zh{可以} (\emph{kěyǐ}) : verbe modal avant le verbe = permission, possibilité, suggestion ; négation \zh{不可以} ; question avec \zh{吗}.
\item \zh{的} : qualifiant + \zh{的} + \textbf{nom} (attribut, possession).
\item \zh{得} : \textbf{verbe} + \zh{得} + adjectif (complément de degré, évaluation).
\item \zh{地} : adjectif + \zh{地} + \textbf{verbe} (complément de manière, « -ment »).
\item Les trois se prononcent \emph{de} (ton neutre) : c'est la \textbf{place dans la phrase} qui décide du caractère.
\end{itemize}
\end{aretenir}

\coursec{Schéma de la structure}

Dans la section précédente, nous avons observé des phrases authentiques contenant le verbe modal \zh{可以} et les trois particules \zh{的}, \zh{得} et \zh{地}. Nous avons remarqué qu'elles apparaissent toujours à des places précises dans la phrase. Le moment est venu de formaliser ces observations : nous allons maintenant écrire les \cle{schémas} (structures générales) de chaque construction, comme on écrit une formule en mathématiques. Une fois ces schémas mémorisés, vous pourrez produire vous-mêmes des phrases correctes, sans hésitation.

\courssub{Le verbe modal 可以 (kěyǐ) : « pouvoir »}

\textbf{Observation.} Reprenons deux phrases de la section 1 :

\zh{学生可以带手机吗？} \emph{Xuésheng kěyǐ dài shǒujī ma?} « Les élèves peuvent-ils apporter leur téléphone portable ? »

\zh{我们不可以迟到。} \emph{Wǒmen bù kěyǐ chídào.} « Nous ne pouvons pas être en retard. »

On constate que \zh{可以} se place \textbf{toujours entre le sujet et le verbe principal}, et qu'il ne change jamais de forme, quel que soit le sujet (\zh{我}, \zh{你}, \zh{他们}...).

\begin{propriete}[Schéma de base avec 可以]
\textbf{Forme affirmative :} Sujet + \zh{可以} + Verbe (+ Objet).

\textbf{Forme négative :} Sujet + \zh{不} + \zh{可以} + Verbe (+ Objet).

\textbf{Forme interrogative :} Sujet + \zh{可以} + Verbe (+ Objet) + \zh{吗} ?

\zh{可以} est un \cle{verbe modal} (aussi appelé \emph{verbe auxiliaire}) : il exprime la \textbf{permission}, la \textbf{possibilité} ou l'\textbf{autorisation}. Il ne se conjugue pas : pas de terminaison, pas d'accord, pas de temps. Il reste identique dans toutes les phrases et se place toujours \textbf{devant le verbe principal}, après le sujet.
\end{propriete}

\begin{exemplebox}[Exemples avec 可以]
\zh{我可以进来吗？} \emph{Wǒ kěyǐ jìnlái ma?} « Puis-je entrer ? »\\
\zh{请问，可以进来吗？} \emph{Qǐngwèn, kěyǐ jìnlái ma?} « Excusez-moi, puis-je entrer ? »\\
\zh{你可以用我的电脑。} \emph{Nǐ kěyǐ yòng wǒ de diànnǎo.} « Tu peux utiliser mon ordinateur. »\\
\zh{我们不可以迟到。} \emph{Wǒmen bù kěyǐ chídào.} « Nous ne pouvons pas être en retard. »\\
\zh{在图书馆不可以大声说话。} \emph{Zài túshūguǎn bù kěyǐ dàshēng shuōhuà.} « À la bibliothèque, on ne peut pas parler fort. »\\
\zh{今天可以上网看新闻。} \emph{Jīntiān kěyǐ shàngwǎng kàn xīnwén.} « Aujourd'hui, on peut aller sur Internet lire les nouvelles. »
\end{exemplebox}

\begin{methode}[Comment former la négation avec 可以]
Pour transformer une phrase affirmative avec \zh{可以} en phrase négative, suivez ces étapes :
\begin{enumerate}
\item Repérez le verbe modal \zh{可以} dans la phrase.
\item Placez l'adverbe de négation \zh{不} \textbf{immédiatement devant} \zh{可以} : \zh{不可以}.
\item Ne changez rien d'autre : ni l'ordre des mots, ni le verbe principal.
\end{enumerate}
Exemple : \zh{他可以看电视。} \emph{Tā kěyǐ kàn diànshì.} (Il peut regarder la télévision.) $\rightarrow$ \zh{他不可以看电视。} \emph{Tā bù kěyǐ kàn diànshì.} (Il ne peut pas regarder la télévision.)
\end{methode}

\begin{attention}[Pas d'inversion après 可以 !]
En français, on inverse parfois le sujet et le verbe pour poser une question (« Peux-tu venir ? »). En chinois, \textbf{l'ordre des mots ne change jamais} : on ajoute simplement \zh{吗} à la fin de la phrase affirmative.\\
\zh{你可以来。} $\rightarrow$ \zh{你可以来吗？} (et jamais *\zh{可以你来吗？})\\
De même, ne placez jamais \zh{不} après \zh{可以} : *\zh{可以不迟到} est fautif ; il faut \zh{不可以迟到}.
\end{attention}

\courssub{La particule 的 (de) : l'attribut du nom}

\textbf{Observation.} Comparez :

\zh{苹果} \emph{píngguǒ} « pomme » $\rightarrow$ \zh{红的苹果} \emph{hóng de píngguǒ} « la pomme rouge »

\zh{老师} \emph{lǎoshī} « professeur » $\rightarrow$ \zh{我们的老师} \emph{wǒmen de lǎoshī} « notre professeur »

La particule \zh{的} introduit un \cle{attribut} (adjectif, nom, pronom) qui précise ou qualifie le nom qui suit.

\begin{propriete}[Schéma de 的]
\textbf{Forme :} (Adjectif / Nom / Pronom) + \zh{的} + Nom.

\zh{的} est la particule de \cle{modification} : tout ce qui précède \zh{的} décrit ou détermine le nom qui le suit. C'est l'équivalent français de « de », de l'accord de l'adjectif ou de l'adjectif possessif.\\
\emph{Remarques :} avec les adjectifs d'une seule syllabe très fréquents (\zh{好}, \zh{大}, \zh{小}), \zh{的} est souvent omis : \zh{好人} « une bonne personne ». Mais avec un adjectif de deux syllabes, \zh{的} est obligatoire : \zh{漂亮的衣服} « de beaux vêtements ». De même, après un pronom personnel suivi d'un terme de parenté ou d'un collectif proche, on omet souvent \zh{的} : \zh{我妈妈} « ma mère », \zh{我们学校} « notre école ».
\end{propriete}

\begin{exemplebox}[Exemples avec 的]
\zh{红的苹果} \emph{hóng de píngguǒ} « la pomme rouge »\\
\zh{我的书} \emph{wǒ de shū} « mon livre »\\
\zh{喀麦隆的首都} \emph{Kāmàilóng de shǒudū} « la capitale du Cameroun »\\
\zh{新来的同学} \emph{xīn lái de tóngxué} « le camarade nouvellement arrivé »\\
\zh{中文的报纸} \emph{Zhōngwén de bàozhǐ} « le journal en chinois »
\end{exemplebox}

\courssub{La particule 得 (de) : le complément de degré}

\textbf{Observation.} Comparez ces deux phrases :

\zh{他做作业。} \emph{Tā zuò zuòyè.} « Il fait ses devoirs. » (simple constat)

\zh{他做得很好。} \emph{Tā zuò de hěn hǎo.} « Il fait (cela) très bien. » (jugement sur la manière)

Dans la deuxième phrase, \zh{得} introduit un commentaire sur \textbf{la manière} ou \textbf{l'intensité} de l'action.

\begin{definition}[Le complément de degré]
Le \cle{complément de degré} est un mot ou un groupe de mots placé après \zh{得} qui indique la \textbf{manière}, l'\textbf{intensité} ou le \textbf{résultat} d'une action ou d'un état. Il est souvent formé d'un adverbe de degré + adjectif (\zh{很好} « très bien », \zh{太快} « trop vite »), ou d'un adjectif renforcé par \zh{极了} (\zh{好极了} « extrêmement bien ») ou \zh{得很} (\zh{累得很} « très fatigué »).
\end{definition}

\begin{propriete}[Schéma de 得]
\textbf{Forme :} Verbe / Adjectif + \zh{得} + Complément de degré.

Exemples de compléments fréquents : \zh{很好} (très bien), \zh{不错} (pas mal), \zh{太快} (trop vite), \zh{完} (terminé), \zh{好极了} (excellent).\\
\emph{Cas particulier :} si le verbe a un objet, on répète le verbe : \zh{他说汉语说得很好。} \emph{Tā shuō Hànyǔ shuō de hěn hǎo.} « Il parle très bien le chinois. » Autre possibilité, placer l'objet en tête : \zh{他汉语说得很好。}
\end{propriete}

\begin{exemplebox}[Exemples avec 得]
\zh{他做得很好。} \emph{Tā zuò de hěn hǎo.} « Il fait très bien. »\\
\zh{她跑得很快。} \emph{Tā pǎo de hěn kuài.} « Elle court très vite. »\\
\zh{我累极了。} \emph{Wǒ lèi jí le.} « Je suis épuisé. »\\
\zh{他高兴得跳了起来。} \emph{Tā gāoxìng de tiào le qǐlái.} « Il était si content qu'il a sauté de joie. »
\end{exemplebox}

\courssub{La particule 地 (de) : l'adverbialisation}

\textbf{Observation.} En français, on transforme un adjectif en adverbe en ajoutant « -ment » : \emph{tranquille} $\rightarrow$ \emph{tranquillement}. Le chinois fait exactement la même chose avec \zh{地} :

\zh{安静} \emph{ānjìng} « tranquille » $\rightarrow$ \zh{安静地读书} \emph{ānjìng de dúshū} « lire tranquillement »

\begin{propriete}[Schéma de 地]
\textbf{Forme :} Adjectif + \zh{地} + Verbe.

\zh{地} transforme un adjectif en \cle{complément circonstanciel de manière} : il décrit \textbf{comment} l'action se déroule. La particule se place toujours \textbf{entre l'adjectif et le verbe}, jamais ailleurs.\\
\emph{Remarque :} avec les adverbes d'une syllabe (\zh{快}, \zh{慢}), \zh{地} est souvent omis : \zh{快走！} « Marche vite ! »
\end{propriete}

\begin{exemplebox}[Exemples avec 地]
\zh{她安静地读书。} \emph{Tā ānjìng de dúshū.} « Elle lit tranquillement. »\\
\zh{他高兴地唱歌。} \emph{Tā gāoxìng de chànggē.} « Il chante joyeusement. »\\
\zh{老师慢慢地解释语法。} \emph{Lǎoshī mànmàn de jiěshì yǔfǎ.} « Le professeur explique lentement la grammaire. »\\
\zh{我们认真地学习汉语。} \emph{Wǒmen rènzhēn de xuéxí Hànyǔ.} « Nous étudions le chinois sérieusement. »
\end{exemplebox}

\courssub{Tableau récapitulatif et vue d'ensemble}

Les trois particules se prononcent toutes \emph{de} (ton neutre), mais s'écrivent avec des caractères différents et ont des fonctions différentes. Ce tableau est à connaître par cœur :

\begin{center}
\begin{tabularx}{\linewidth}{|X|X|X|X|}
\hline
\rowcolor{popblueL} Particule & Schéma & Exemple (caractères + pinyin) & Traduction \\
\hline
\zh{的} & (Adj/Nom) + \zh{的} + Nom & \zh{红的苹果} hóng de píngguǒ & la pomme rouge \\
\hline
\zh{得} & V/Adj + \zh{得} + Compl. de degré & \zh{他做得很好。} Tā zuò de hěn hǎo. & Il fait très bien. \\
\hline
\zh{地} & Adj + \zh{地} + Verbe & \zh{她安静地读书。} Tā ānjìng de dúshū. & Elle lit tranquillement. \\
\hline
\end{tabularx}
\end{center}

\begin{center}
\begin{tabularx}{\linewidth}{|X|X|X|}
\hline
\rowcolor{popblueL} Structure & Exemple en caractères + pinyin & Traduction \\
\hline
Sujet + \zh{可以} + V (+ O) & \zh{我可以用手机。} Wǒ kěyǐ yòng shǒujī. & Je peux utiliser le téléphone. \\
\hline
Sujet + \zh{不可以} + V & \zh{我们不可以迟到。} Wǒmen bù kěyǐ chídào. & Nous ne pouvons pas être en retard. \\
\hline
Sujet + \zh{可以} + V + \zh{吗} ? & \zh{请问可以进来吗？} Qǐngwèn kěyǐ jìnlái ma? & Puis-je entrer, s'il vous plaît ? \\
\hline
\end{tabularx}
\end{center}

\begin{popfigure}
\begin{tikzpicture}[scale=0.8]
\node[draw,rounded corners,fill=popblueL] (k) {Sujet + 可以 + V (+ O)};
\node[draw,rounded corners,below=1cm of k,fill=poporangeL] (de) {Adj/Nom + 的 + Nom};
\node[draw,rounded corners,below=1cm of de,fill=popgreenL] (de2) {V/Adj + 得 + Compl. degré};
\node[draw,rounded corners,below=1cm of de2,fill=poppinkL] (de3) {Adj + 地 + V};
\draw[->,thick,popdark] (k) -- (de);
\draw[->,thick,popdark] (de) -- (de2);
\draw[->,thick,popdark] (de2) -- (de3);
\end{tikzpicture}
\legende{Les quatre schémas de la leçon : le verbe modal 可以 et les trois particules 的、得、地.}
\end{popfigure}

\textbf{Règle d'or de l'ordre des mots.} En chinois, la particule se place \textbf{toujours immédiatement après le mot qu'elle suit} et \textbf{avant le mot modifié} : \zh{的} avant le nom, \zh{得} après le verbe, \zh{地} avant le verbe. Quant à \zh{可以}, il reste \textbf{collé devant le verbe principal}, sans rien entre les deux.

\courssub{Exercice résolu et pièges à éviter}

\begin{exoresolu}[Compléter avec 的、得 ou 地]
Énoncé. Complétez chaque phrase avec la bonne particule (\zh{的}, \zh{得} ou \zh{地}) :\\
1. 她唱\_\_很好听。 2. 这是我\_\_课本。 3. 孩子们高兴\_\_玩游戏。 4. 他跑\_\_很快。
\tcblower
Solution détaillée.\\
1. \zh{她唱得很好听。} \emph{Tā chàng de hěn hǎotīng.} « Elle chante très bien. » — Après le verbe \zh{唱}, on commente la manière : c'est un complément de degré, donc \zh{得}.\\
2. \zh{这是我的课本。} \emph{Zhè shì wǒ de kèběn.} « C'est mon manuel. » — Le pronom \zh{我} détermine le nom \zh{课本} : attribut, donc \zh{的}.\\
3. \zh{孩子们高兴地玩游戏。} \emph{Háizimen gāoxìng de wán yóuxì.} « Les enfants jouent joyeusement. » — L'adjectif \zh{高兴} modifie le verbe \zh{玩} : adverbialisation, donc \zh{地}.\\
4. \zh{他跑得很快。} \emph{Tā pǎo de hěn kuài.} « Il court très vite. » — Complément de degré après le verbe \zh{跑}, donc \zh{得}.\\
Méthode : demandez-vous toujours « quel mot est modifié ? ». Si c'est un \textbf{nom} $\rightarrow$ \zh{的} ; si c'est un \textbf{verbe dont on évalue le résultat} $\rightarrow$ \zh{得} ; si c'est un \textbf{verbe dont on décrit la manière en amont} $\rightarrow$ \zh{地}.
\end{exoresolu}

\begin{attention}[Les trois erreurs classiques des francophones]
\begin{enumerate}
\item Ne jamais insérer \zh{的} entre \zh{可以} et le verbe : *\zh{可以的进去} est une faute grave. \zh{可以} se colle directement au verbe : \zh{可以进去}.
\item Ne pas confondre les trois caractères à l'écrit : ils se prononcent tous \emph{de}, mais \zh{的} (clé 白 « blanc »), \zh{得} (clé 彳 « pas ») et \zh{地} (clé 土 « terre ») s'écrivent différemment. À l'examen, l'emploi du mauvais caractère est sanctionné.
\item Ne pas traduire mécaniquement « très » par \zh{很} partout : dans \zh{他做得很好}, \zh{很} fait partie du complément de degré et ne s'omet pas, sinon la phrase sonne comme une comparaison.
\end{enumerate}
\end{attention}

\begin{savaistu}[可以、会、能 : trois « pouvoirs » différents]
En français, un seul verbe « pouvoir » traduit plusieurs idées. Le chinois distingue : \zh{可以} exprime la \textbf{permission} ou la possibilité extérieure (\zh{我可以走了。} « Je peux partir » = j'ai l'autorisation) ; \zh{会} (huì) exprime la \textbf{capacité acquise}, le « savoir-faire » (\zh{他会说汉语。} « Il sait parler chinois ») ; \zh{能} (néng) exprime la \textbf{capacité physique} ou la possibilité concrète (\zh{他能游一千米。} « Il peut nager mille mètres »). Attention donc : « savoir + infinitif » en français ne se traduit presque jamais par \zh{可以} !
\end{savaistu}

\begin{aretenir}[L'essentiel de la section]
\begin{itemize}
\item \zh{可以} : Sujet + \zh{可以} + V ; négation \zh{不可以} ; question avec \zh{吗} final ; jamais d'inversion, jamais de conjugaison.
\item \zh{的} : (Adj/Nom) + \zh{的} + Nom $\rightarrow$ attribut du nom.
\item \zh{得} : V/Adj + \zh{得} + complément de degré $\rightarrow$ évaluation de l'action.
\item \zh{地} : Adj + \zh{地} + V $\rightarrow$ adverbe de manière.
\item Les trois particules se prononcent \emph{de} mais ne s'écrivent ni ne s'emploient jamais de la même façon.
\end{itemize}
\end{aretenir}

\coursec{Emplois et exceptions}

Dans la section précédente, nous avons établi les \emph{schémas} des structures : \zh{可以} + verbe, \zh{的} comme marque de détermination et de nominalisation, \zh{得} introduisant le complément de degré, \zh{地} marquant le complément circonstanciel de manière. Savoir construire le squelette ne suffit pas : il faut maintenant savoir \textbf{quand} employer chaque forme, \textbf{dans quelles limites}, et \textbf{quelles exceptions} surveiller. C'est l'objet de cette section, la plus importante pour l'examen.

\courssub{Le verbe modal \zh{可以} : permission et possibilité objective}

\begin{textebox}[À la bibliothèque du lycée de Yaoundé]
学生：老师，我可以借这本书吗？\\
\emph{Xuésheng : Lǎoshī, wǒ kěyǐ jiè zhè běn shū ma ?}\\
« Élève : Monsieur, est-ce que je peux emprunter ce livre ? »\\[4pt]
老师：可以，但是你不可以在图书馆里吃东西。\\
\emph{Lǎoshī : Kěyǐ, dànshì nǐ bù kěyǐ zài túshūguǎn lǐ chī dōngxi.}\\
« Professeur : Oui, mais tu ne peux pas manger dans la bibliothèque. »\\[4pt]
学生：明天可以还书吗？\\
\emph{Xuésheng : Míngtiān kěyǐ huán shū ma ?}\\
« Élève : Est-ce qu'on peut rendre le livre demain ? »\\[4pt]
老师：可以。图书馆下午五点关门，五点以前都可以来。\\
\emph{Lǎoshī : Kěyǐ. Túshūguǎn xiàwǔ wǔ diǎn guān mén, wǔ diǎn yǐqián dōu kěyǐ lái.}\\
« Professeur : Oui. La bibliothèque ferme à 17 h, on peut venir avant 17 h. »
\end{textebox}

\textbf{Observation.} Dans ce dialogue, \zh{可以} apparaît quatre fois. Deux emplois se distinguent :

\begin{enumerate}
\item \textbf{La permission} (autorisation donnée ou demandée par une personne) : \zh{我可以借这本书吗？} — l'élève demande l'autorisation au professeur.
\item \textbf{La possibilité objective} (rendue possible par les circonstances, les horaires, les conditions) : \zh{五点以前都可以来} — ce n'est pas une question d'autorisation personnelle, mais une condition objective (l'horaire d'ouverture).
\end{enumerate}

\begin{definition}[\zh{可以} : les deux emplois fondamentaux]
\zh{可以} \emph{kěyǐ} est un verbe modal placé \textbf{avant le verbe principal} (Sujet + \zh{可以} + Verbe). Il exprime :
\begin{itemize}
\item la \cle{permission} : quelqu'un autorise ou interdit une action (« pouvoir » au sens de « avoir le droit de ») ;
\item la \cle{possibilité objective} : les conditions matérielles, le temps, la météo, la loi rendent une action possible (« il est possible de »).
\end{itemize}
Il n'exprime \textbf{jamais} la capacité physique ni le savoir-faire : c'est là la grande différence avec \zh{能} et \zh{会}.
\end{definition}

\begin{exemplebox}[Permission et possibilité objective]
\zh{妈妈说我可以看电视。} \emph{Māma shuō wǒ kěyǐ kàn diànshì.} « Maman dit que je peux regarder la télévision. » (permission)\\[3pt]
\zh{今天不下雨，我们可以去杜阿拉。} \emph{Jīntiān bù xià yǔ, wǒmen kěyǐ qù Dù'ālā.} « Il ne pleut pas aujourd'hui, nous pouvons aller à Douala. » (possibilité objective : le temps le permet)\\[3pt]
\zh{在网上可以买到很多东西。} \emph{Zài wǎngshang kěyǐ mǎi dào hěn duō dōngxi.} « Sur Internet, on peut acheter beaucoup de choses. » (possibilité objective)
\end{exemplebox}

\begin{propriete}[Exception essentielle : la négation \zh{不可以} exprime l'interdiction]
À la forme négative, \zh{不可以} \emph{bù kěyǐ} ne signifie pas seulement « ce n'est pas possible » : il exprime l'\cle{interdiction}, le « il ne faut pas », le « il est défendu de ». C'est la formule standard des règlements (école, bibliothèque, hôpital, panneaux).
\end{propriete}

\begin{exemplebox}[\zh{不可以} = interdit]
\zh{你不可以吸烟。} \emph{Nǐ bù kěyǐ xīyān.} « Tu ne dois pas fumer. / Il est interdit de fumer. »\\[3pt]
\zh{上课的时候不可以打电话。} \emph{Shàngkè de shíhou bù kěyǐ dǎ diànhuà.} « Pendant le cours, il est interdit de téléphoner. »\\[3pt]
\zh{这里不可以停车。} \emph{Zhèlǐ bù kěyǐ tíngchē.} « Il est interdit de stationner ici. »
\end{exemplebox}

\begin{attention}[Nuances de la négation]
\zh{不可以} est fort et formel : c'est l'interdiction écrite ou officielle. Pour une interdiction plus douce ou orale, on utilise \zh{别} \emph{bié} (« ne … pas ! ») : \zh{别吸烟！} \emph{Bié xīyān !} « Ne fume pas ! ». À l'examen, pour traduire « il est interdit de », la réponse attendue est \zh{不可以} + verbe.
\end{attention}

\courssub{\zh{可以} ≠ \zh{能} ≠ \zh{会} : la distinction capitale}

Le français n'a qu'un seul verbe, « pouvoir » (et « savoir »), là où le chinois en a trois. C'est la source d'innombrables erreurs chez les francophones. Observons :

\begin{exemplebox}[Trois phrases, trois modaux, trois sens]
\zh{我可以出去吗？} \emph{Wǒ kěyǐ chūqù ma ?} « Est-ce que je peux sortir ? » (je demande la \textbf{permission})\\[3pt]
\zh{他生病了，今天不能来上课。} \emph{Tā shēngbìng le, jīntiān bù néng lái shàngkè.} « Il est malade, il ne peut pas venir en cours aujourd'hui. » (\textbf{capacité physique} ou condition concrète : son corps ne le permet pas)\\[3pt]
\zh{我会说汉语。} \emph{Wǒ huì shuō Hànyǔ.} « Je sais parler chinois. » (\textbf{compétence acquise} par l'apprentissage)
\end{exemplebox}

\begin{definition}[\zh{可以} / \zh{能} / \zh{会} : trois mondes différents]
\begin{itemize}
\item \zh{可以} \emph{kěyǐ} : la \cle{permission} (une autorité, un règlement) et la possibilité objective. Question : « Ai-je le droit ? »
\item \zh{能} \emph{néng} : la \cle{capacité physique} ou la possibilité concrète du moment (santé, force, temps disponible). Question : « En suis-je physiquement capable ? »
\item \zh{会} \emph{huì} : la \cle{compétence acquise}, le savoir-faire appris (langue, conduite, natation, cuisine). Question : « Sais-je le faire ? »
\end{itemize}
\end{definition}

\begin{popfigure}
\begin{tikzpicture}[node distance=1.6cm, every node/.style={font=\small}]
\node[draw,rounded corners,fill=popblueL] (perm) {Permission : \zh{可以}};
\node[draw,rounded corners,fill=poporangeL,right=2cm of perm] (cap) {Capacité : \zh{能}};
\node[draw,rounded corners,fill=popgreenL,right=2cm of cap] (skill) {Compétence : \zh{会}};
\draw[<->,thick,popdark] (perm) -- (cap);
\draw[<->,thick,popdark] (cap) -- (skill);
\end{tikzpicture}
\legende{Les trois modaux du « pouvoir » français : \zh{可以} (autorisation), \zh{能} (capacité), \zh{会} (savoir-faire).}
\end{popfigure}

\begin{center}
\begin{tabularx}{\linewidth}{|X|X|X|X|}
\hline
\rowcolor{popblueL} Modal & Sens & Exemple (caractères + pinyin) & Traduction \\
\hline
\zh{可以} & permission, possibilité objective & \zh{你可以进来。} \emph{Nǐ kěyǐ jìnlái.} & Tu peux entrer (je t'y autorise). \\
\hline
\zh{能} & capacité physique, possibilité concrète & \zh{他能搬动这个箱子。} \emph{Tā néng bāndòng zhège xiāngzi.} & Il peut soulever cette caisse (il en a la force). \\
\hline
\zh{会} & compétence acquise & \zh{她会开车。} \emph{Tā huì kāichē.} & Elle sait conduire. \\
\hline
\zh{不可以} & interdiction & \zh{你不可以吸烟。} \emph{Nǐ bù kěyǐ xīyān.} & Tu ne dois pas fumer. \\
\hline
\end{tabularx}
\end{center}

\begin{attention}[Piège n°1 : le double modal]
En chinois, \textbf{on ne cumule jamais deux verbes modaux} devant le même verbe. La phrase *\zh{他可以会开车} est fautive : il faut choisir. Pour « il sait conduire », dites \zh{他会开车} ; pour « il a le droit de conduire » (il a son permis), dites \zh{他可以开车}. Autre erreur classique : traduire « je ne peux pas nager » par *\zh{我不可以游泳} alors qu'on parle d'un savoir-faire : il faut \zh{我不会游泳} (je ne sais pas nager). Si c'est une interdiction (piscine fermée), \zh{不可以游泳} est correct.
\end{attention}

\begin{exoresolu}[Choisir le bon modal]
\textbf{Énoncé.} Complétez avec \zh{可以}, \zh{能} ou \zh{会} : (1) 老师，我 \_\_ 去洗手间吗？ (2) 我妹妹五岁，她不 \_\_ 写字。 (3) 他腿疼，今天不 \_\_ 跑步。 (4) 教室里不 \_\_ 吃东西。
\tcblower
\textbf{Solution détaillée.} (1) \zh{可以} : l'élève demande la \emph{permission} au professeur → \zh{我可以去洗手间吗？} « Puis-je aller aux toilettes ? ». (2) \zh{会} : écrire est une \emph{compétence acquise} → \zh{她不会写字} « elle ne sait pas écrire ». (3) \zh{能} : la douleur à la jambe est une \emph{capacité physique} empêchée → \zh{他今天不能跑步} « il ne peut pas courir aujourd'hui ». (4) \zh{可以} : règlement, \emph{interdiction} → \zh{教室里不可以吃东西} « il est interdit de manger en classe ».
\end{exoresolu}

\courssub{\zh{的} : nominalisation et possession après un pronom}

Nous avons vu dans la section 2 que \zh{的} relie un déterminant à un nom (\zh{我的老师}, \zh{红色的花}). Deux emplois particuliers méritent une attention spéciale.

\begin{propriete}[\zh{的} nominalise une proposition verbale]
Placé \textbf{après un verbe ou une proposition verbale}, \zh{的} transforme cette proposition en groupe nominal : « ce qui / ce que … ». Le nom sous-entendu (chose, personne) n'a pas besoin d'être répété. Structure : \textbf{(Sujet +) Verbe + Complément + 的} = « ce que (sujet) verbe ».
\end{propriete}

\begin{exemplebox}[La nominalisation par \zh{的}]
\zh{吃的} \emph{chī de} « ce qu'on mange ; de quoi manger »\\[3pt]
\zh{这是我买的。} \emph{Zhè shì wǒ mǎi de.} « C'est ce que j'ai acheté. »\\[3pt]
\zh{我喜欢吃的菜是恩多雷。} \emph{Wǒ xǐhuan chī de cài shì Ēnduōléi.} « Le plat que j'aime manger, c'est le ndolé. »\\[3pt]
\zh{他说的话很有意思。} \emph{Tā shuō de huà hěn yǒu yìsi.} « Ce qu'il dit est très intéressant. »\\[3pt]
\zh{穿红衣服的是我的哥哥。} \emph{Chuān hóng yīfu de shì wǒ de gēge.} « Celui qui porte un vêtement rouge, c'est mon grand frère. »
\end{exemplebox}

Notez que le français utilise un pronom relatif (« ce que », « celui qui ») : le chinois, lui, se contente de coller \zh{的} après la proposition. C'est un outil d'une économie remarquable.

\begin{propriete}[\zh{的} après un pronom personnel : la possession]
Après un pronom, \zh{的} marque l'\cle{appartenance} : \zh{我的} \emph{wǒ de} « mon/ma/mes », \zh{你的} \emph{nǐ de} « ton/ta », \zh{他的} \emph{tā de} « son/sa », \zh{我们的} « notre », \zh{你们的} « votre », \zh{他们的} « leur ». Exception bien connue : devant un nom de parenté ou de relation proche, \zh{的} est souvent omis : \zh{我妈妈} \emph{wǒ māma} « ma mère », \zh{他朋友} « son ami ». Mais on ne peut jamais l'omettre devant un objet : \zh{我的书} « mon livre ».
\end{propriete}

\begin{attention}[Piège : ne pas confondre les deux \zh{的}]
Dans \zh{这是我买的}, le \zh{的} final \emph{nominalise} (« ce que j'ai acheté ») ; dans \zh{这是我的书}, le \zh{的} marque la \emph{possession} devant un nom explicite. Test simple : si un nom suit immédiatement \zh{的}, c'est un déterminant ; si la phrase se termine par \zh{的}, c'est une nominalisation.
\end{attention}

\courssub{\zh{得} : le complément de degré — obligatoire, mais pas toujours}

\begin{propriete}[Quand \zh{得} est obligatoire, quand il disparaît]
\zh{得} \emph{de} est \textbf{obligatoire} dès qu'un complément de degré ou de manière suit le verbe : \textbf{Verbe + 得 + complément}. En revanche, \textbf{s'il n'y a pas de complément}, \zh{得} disparaît tout simplement : on dit \zh{她跑得很快} « elle court vite », mais \zh{她跑步} « elle court » (jamais *\zh{她得跑步}). Autre règle : si le verbe a un objet, on répète le verbe : \zh{他唱歌唱得很好} « il chante bien ».
\end{propriete}

\begin{propriete}[Les compléments fréquents après \zh{得}]
Après \zh{得}, on trouve très souvent un \cle{adverbe de degré} + adjectif : \zh{很好} « très bien », \zh{真快} « vraiment vite », \zh{特别清楚} « particulièrement clairement », ou encore \zh{得} + adjectif + \zh{极了} \emph{jí le} « extrêmement » : \zh{她唱得好极了} « elle chante extrêmement bien ». Attention : \zh{极了} peut aussi se coller \textbf{directement} à un adjectif, sans \zh{得} : \zh{天气热极了} « il fait extrêmement chaud ».
\end{propriete}

\begin{exemplebox}[\zh{得} en action]
\zh{她跑得真快。} \emph{Tā pǎo de zhēn kuài.} « Elle court vraiment vite. »\\[3pt]
\zh{他汉语说得特别流利。} \emph{Tā Hànyǔ shuō de tèbié liúlì.} « Il parle chinois particulièrement couramment. »\\[3pt]
\zh{喀麦隆的足球队踢得好极了。} \emph{Kāmàilóng de zúqiúduì tī de hǎo jí le.} « L'équipe de football du Cameroun joue extrêmement bien. »\\[3pt]
\zh{你昨天睡得好吗？} \emph{Nǐ zuótiān shuì de hǎo ma ?} « As-tu bien dormi hier ? »
\end{exemplebox}

\courssub{\zh{地} : requis à l'écrit, souvent omis à l'oral}

\begin{propriete}[\zh{地} à l'écrit et à l'oral]
\zh{地} \emph{de} relie un adverbe de manière au verbe : \textbf{Adverbe + 地 + Verbe}. À l'\textbf{oral rapide}, les locuteurs l'omettent souvent (\zh{他高兴笑了} s'entend), mais à l'\textbf{écrit formel} — et donc à l'examen MINESEC — \zh{地} est \textbf{requis} après un adverbe de deux syllabes ou plus : \zh{他高兴地笑了} « il a ri joyeusement ». Après un adverbe monosyllabique très court (\zh{快走} « marche vite »), on ne le met pas.
\end{propriete}

\begin{exemplebox}[\zh{地} à l'écrit]
\zh{他高兴地笑了。} \emph{Tā gāoxìng de xiào le.} « Il a ri joyeusement. »\\[3pt]
\zh{同学们认真地听老师讲课。} \emph{Tóngxuémen rènzhēn de tīng lǎoshī jiǎngkè.} « Les élèves écoutent attentivement le professeur. »\\[3pt]
\zh{她慢慢地走进教室。} \emph{Tā mànmàn de zǒujìn jiàoshì.} « Elle est entrée lentement dans la salle de classe. »
\end{exemplebox}

\begin{methode}[Le test de traduction pour choisir entre \zh{得} et \zh{地}]
Face à une phrase à traduire ou à compléter :
\begin{enumerate}
\item Traduisez mentalement le complément par « \textbf{de manière à} / d'une manière … » : s'il précède le verbe en chinois et décrit \emph{comment} l'action se déroule → \zh{地} (\zh{认真地听} « écouter de manière attentive »).
\item Traduisez par « \textbf{si … que} / tellement … que » : si le complément \emph{suit} le verbe et évalue son résultat ou son degré → \zh{得} (\zh{跑得很快} « courir si vite que ça se voit »).
\item Si le mot-outil relie un déterminant à un nom → \zh{的}.
\end{enumerate}
Repère de position : \zh{地} est \textbf{avant} le verbe, \zh{得} est \textbf{après} le verbe, \zh{的} est \textbf{avant un nom}.
\end{methode}

\begin{center}
\begin{tabularx}{\linewidth}{|X|X|X|X|}
\hline
\rowcolor{popblueL} Mot-outil & Position & Fonction & Exemple + traduction \\
\hline
\zh{的} & avant un nom & détermination, possession, nominalisation & \zh{我买的书} « le livre que j'ai acheté » \\
\hline
\zh{地} & avant un verbe & manière (adverbe) & \zh{高兴地笑} « rire joyeusement » \\
\hline
\zh{得} & après un verbe / adjectif & degré, résultat, évaluation & \zh{跑得真快} « courir vraiment vite » \\
\hline
\end{tabularx}
\end{center}

\begin{exoresolu}[\zh{的}, \zh{得} ou \zh{地} ?]
\textbf{Énoncé.} Complétez avec \zh{的}, \zh{得} ou \zh{地} : (1) 这是我姐姐买 \_\_ 衣服。 (2) 她唱歌唱 \_\_ 很好听。 (3) 孩子们快乐 \_\_ 玩游戏。 (4) 他 \_\_ 手机很漂亮。
\tcblower
\textbf{Solution détaillée.} (1) \zh{的} : la proposition \zh{我姐姐买} détermine le nom \zh{衣服} → \zh{这是我姐姐买的衣服。} « Ce sont les vêtements que ma sœur aînée a achetés. » (2) \zh{得} : complément de degré \emph{après} le verbe répété → \zh{她唱歌唱得很好听。} « Elle chante très bien (agréablement). » (3) \zh{地} : adverbe de manière \emph{avant} le verbe \zh{玩} → \zh{孩子们快乐地玩游戏。} « Les enfants jouent joyeusement. » (4) \zh{的} : possession après le pronom \zh{他} devant le nom \zh{手机} → \zh{他的手机很漂亮。} « Son téléphone est très beau. »
\end{exoresolu}

\begin{attention}[Récapitulatif des erreurs fréquentes des francophones]
\begin{itemize}
\item Double modal : *\zh{他可以会开车} → choisir \zh{可以} \emph{ou} \zh{会}.
\item Traduire « je ne sais pas nager » par *\zh{我不可以游泳} → \zh{我不会游泳}.
\item Oublier \zh{得} devant le complément : *\zh{她跑很快} → \zh{她跑得很快}.
\item Mettre \zh{得} avant le verbe ou \zh{地} après : retenir les positions (\zh{地} + V ; V + \zh{得}).
\item Oublier \zh{地} à l'écrit après un adverbe long : *\zh{他高兴笑了} → \zh{他高兴地笑了}.
\item Croire que \zh{不可以} signifie « impossible » : c'est une \emph{interdiction}.
\end{itemize}
\end{attention}

\begin{savaistu}[\zh{可以} et \zh{得} à Taiwan]
À Taiwan, on écrit les caractères traditionnels (繁體字), mais \zh{可以} s'y écrit exactement de la même façon qu'en simplifié : ses caractères n'ont pas de forme traditionnelle différente. En revanche, l'usage de \zh{得} y est encore plus fréquent à l'oral que sur le continent : les Taïwanais ponctuent volontiers leurs évaluations de \zh{得} (\zh{好吃得不得了} \emph{hǎochī de bùdéliǎo} « tellement bon qu'on n'en peut plus »). Quant à \zh{地}, il reste partout la marque de l'écrit soigné : un élève qui maîtrise le trio \zh{的}、\zh{地}、\zh{得} passe immédiatement pour un sinisant sérieux, car même de nombreux Chinois les confondent dans leurs SMS !
\end{savaistu}

\begin{aretenir}[L'essentiel de la section]
\zh{可以} = permission et possibilité objective ; sa négation \zh{不可以} = interdiction. \zh{能} = capacité physique ; \zh{会} = savoir-faire appris. Jamais deux modaux ensemble. \zh{的} après une proposition nominalise (« ce que … ») et après un pronom marque la possession (\zh{我的}). \zh{得} est obligatoire devant tout complément de degré qui suit le verbe (\zh{跑得真快}, \zh{唱得好极了}) et absent sans complément ; \zh{极了} peut aussi se coller directement à un adjectif (\zh{热极了}). \zh{地} est requis à l'écrit entre l'adverbe de manière et le verbe (\zh{高兴地笑了}). Positions : \zh{的} + nom ; \zh{地} + verbe ; verbe + \zh{得}.
\end{aretenir}

\coursec{Comparaison avec le français et pièges des francophones}

Dans la section précédente, nous avons étudié les emplois et les exceptions du verbe modal \zh{可以} et des trois particules \zh{的}、\zh{得}、\zh{地}. Nous savons maintenant \emph{quand} les utiliser. Reste le point le plus délicat pour un élève francophone : ces structures ne coïncident \emph{jamais exactement} avec le français. Le mot français « de » ne correspond pas toujours à \zh{的}, l'adverbe en « -ment » ne se place pas au même endroit que le groupe avec \zh{地}, et « pouvoir » ne se traduit pas mécaniquement par \zh{可以} n'importe où dans la phrase. Cette section confronte systématiquement les deux langues pour transformer vos erreurs typiques en réflexes corrects.

\courssub{可以 ↔ « pouvoir / avoir la permission de » + infinitif}

\begin{definition}[Particule à ton neutre]
Une \cle{particule à ton neutre} est une syllabe qui n'a pas de ton propre : elle se prononce brièvement et légèrement, en s'appuyant sur la syllabe précédente (on parle parfois de « 5\textsuperscript{e} ton »). \zh{的} (de), \zh{得} (de) et \zh{地} (de) sont trois particules à ton neutre : à l'oral, elles se prononcent toutes \emph{de}, exactement de la même façon. Seule l'écriture les distingue.
\end{definition}

En français, « pouvoir » est un verbe qui se conjugue et marque le temps : \emph{je peux, je pouvais, je pourrai}. En chinois, \zh{可以} (kěyǐ) est un \cle{verbe modal} : il est \textbf{invariable} et \textbf{ne marque pas le temps}. C'est le contexte (un complément de temps comme \zh{昨天} ou \zh{明天}, ou un marqueur d'aspect comme \zh{了}) qui indique le passé ou le futur.

\begin{center}
\begin{tabularx}{\linewidth}{|X|X|X|}
\hline
\rowcolor{popblueL} Français & Chinois (caractères + pinyin) & Remarque \\
\hline
Je peux y aller. & 我可以去。Wǒ kěyǐ qù. & présent \\
\hline
Hier, j'ai pu y aller. & 昨天我可以去。Zuótiān wǒ kěyǐ qù. & \zh{昨天} marque le passé, pas \zh{可以} \\
\hline
Demain, je pourrai y aller. & 明天我可以去。Míngtiān wǒ kěyǐ qù. & \zh{明天} marque le futur \\
\hline
\end{tabularx}
\end{center}

\begin{attention}[Erreur fréquente : la place de 可以]
Le français dit « Il peut aller » : le verbe modal vient \emph{avant} le verbe principal, et le chinois fait de même : \zh{我可以去。} (Wǒ kěyǐ qù.) est correct. Mais le francophone qui « pense français » produit parfois \zh{我去可以} (*Wǒ qù kěyǐ), en plaçant le modal après le verbe : c'est \textbf{incorrect}. De même, ne placez jamais l'objet avant \zh{可以} dans une phrase déclarative neutre : on dit \zh{我可以吃苹果。} (Wǒ kěyǐ chī píngguǒ, « Je peux manger la pomme »), jamais \zh{我苹果可以吃}. La négation se place devant le modal : \zh{不可以} (bù kěyǐ, « ne pas avoir le droit de »).
\end{attention}

\begin{exemplebox}[Exemples traduits : 可以]
\zh{你可以帮我吗？} \emph{Nǐ kěyǐ bāng wǒ ma?} « Peux-tu m'aider ? »\\
\zh{这里可以拍照吗？} \emph{Zhèlǐ kěyǐ pāizhào ma?} « Peut-on prendre des photos ici ? »\\
\zh{学生不可以在教室里打电话。} \emph{Xuésheng bù kěyǐ zài jiàoshì lǐ dǎ diànhuà.} « Les élèves n'ont pas le droit de téléphoner en classe. » (permission refusée : \zh{不可以})
\end{exemplebox}

\begin{propriete}[可以, 能, 会 : trois « pouvoir » à ne pas confondre]
Le français n'a qu'un mot « pouvoir » ; le chinois en a trois, aux nuances distinctes :
\begin{itemize}
\item \zh{可以} (kěyǐ) : la \textbf{permission}, l'autorisation → \zh{我可以进来吗？} (Wǒ kěyǐ jìnlái ma? « Puis-je entrer ? ») ;
\item \zh{能} (néng) : la \textbf{capacité concrète}, la possibilité matérielle → \zh{他今天能来。} (Tā jīntiān néng lái. « Il peut venir aujourd'hui » : rien ne l'en empêche) ;
\item \zh{会} (huì) : la \textbf{compétence acquise}, le savoir-faire → \zh{我会开车。} (Wǒ huì kāichē. « Je sais conduire. »).
\end{itemize}
Dans beaucoup de contextes, \zh{可以} et \zh{能} sont interchangeables ; mais pour une permission demandée ou refusée, \zh{可以} est le choix le plus sûr à l'examen.
\end{propriete}

\courssub{的 ↔ « de » possessif ou relatif « qui/que »}

Le français exprime la possession avec « de » \emph{avant} le nom possesseur : « le livre \textbf{de} mon frère ». Le chinois place \zh{的} \emph{après} le possesseur : \zh{我哥哥的书} (wǒ gēge de shū), littéralement « mon frère + 的 + livre ». L'ordre est donc \textbf{inversé} par rapport au français.

\begin{propriete}[Ordre avec 的]
Schéma : \cle{Possesseur / Déterminant + 的 + Nom}. Le déterminant (adjectif, pronom, proposition relative) précède toujours le nom, relié par \zh{的}. En français, l'adjectif épithète suit souvent le nom (« un livre \emph{intéressant} ») ; en chinois, il le précède toujours : \zh{有意思的书} (yǒu yìsi de shū), « un livre intéressant ». Exception : entre un pronom personnel et un nom de parenté ou de collectivité proche, \zh{的} peut être omis : \zh{我妈妈} (wǒ māma, « ma mère »), \zh{我们学校} (wǒmen xuéxiào, « notre école »).
\end{propriete}

\begin{exemplebox}[Exemples traduits : 的]
\zh{我的书} \emph{wǒ de shū} « mon livre »\\
\zh{喀麦隆的首都} \emph{Kāmàilóng de shǒudū} « la capitale du Cameroun »\\
\zh{我昨天买的手机} \emph{wǒ zuótiān mǎi de shǒujī} « le téléphone que j'ai acheté hier » (proposition relative : le français utilise « que », le chinois place toute la relative \emph{avant} le nom et la termine par \zh{的})
\end{exemplebox}

\begin{attention}[Piège 2 : la relative chinoise est « à l'envers »]
Français : « le téléphone \emph{que j'ai acheté hier} » — la relative suit le nom. Chinois : \zh{我昨天买的手机} — la relative \emph{précède} le nom. Ne dites jamais \zh{手机我昨天买的} pour « le téléphone que j'ai acheté hier ». Réflexe : tout ce qui décrit le nom se met devant lui, suivi de \zh{的}.
\end{attention}

\courssub{得 ↔ « verbe + si/très + adjectif »}

Le français dit « il court \emph{si vite} », « elle chante \emph{très bien} » : le degré est exprimé par un adverbe (si, très) placé \emph{avant} l'adjectif, et le tout suit le verbe. Le chinois utilise la particule \zh{得} \emph{immédiatement après le verbe}, suivie de l'adjectif : \zh{他跑得快。} (Tā pǎo de kuài.) « Il court vite. » Il n'y a \textbf{pas de mot « de »} en français, alors que le chinois exige \zh{得} : c'est la source de l'oubli le plus fréquent.

\begin{exemplebox}[Exemples traduits : 得]
\zh{他跑得快。} \emph{Tā pǎo de kuài.} « Il court vite. »\\
\zh{她汉语说得很好。} \emph{Tā Hànyǔ shuō de hěn hǎo.} « Elle parle très bien le chinois. »\\
\zh{他高兴得跳了起来。} \emph{Tā gāoxìng de tiào le qǐlái.} « Il était si content qu'il a sauté. »
\end{exemplebox}

\begin{attention}[Piège 3 : oublier 得]
On ne peut pas juxtaposer verbe et adjectif de manière : \zh{他跑快} (*Tā pǎo kuài) est \textbf{incorrect}. Dès qu'un complément de degré ou de manière suit le verbe, \zh{得} est obligatoire : \zh{他跑得快。} Astuce : si en français vous pouvez dire « si/très + adjectif » après le verbe, mettez \zh{得} en chinois. Autre difficulté : si le verbe a un objet, on répète le verbe devant \zh{得} : \zh{他开车开得很好。} (Tā kāichē kāi de hěn hǎo. « Il conduit très bien. »), ou on avance l'objet : \zh{她汉语说得很好。}
\end{attention}

\courssub{地 ↔ adjectif + « -ment »}

\begin{propriete}[Place de 地 : l'inverse du français]
En français, l'adverbe en « -ment » se place \emph{après} le verbe : « elle parle \emph{doucement} ». En chinois, le groupe « adjectif + \zh{地} » se place \emph{avant} le verbe : \zh{她轻声地说话。} (Tā qīngshēng de shuōhuà.) Schéma : \cle{Adjectif + 地 + Verbe}. Remarque : avec les adjectifs monosyllabiques redoublés, \zh{地} est facultatif mais très fréquent : \zh{慢慢地走} (mànmàn de zǒu, « marcher lentement »).
\end{propriete}

\begin{exemplebox}[Exemples traduits : 地]
\zh{她轻声地说话。} \emph{Tā qīngshēng de shuōhuà.} « Elle parle doucement. »\\
\zh{孩子们高兴地唱歌。} \emph{Háizimen gāoxìng de chànggē.} « Les enfants chantent joyeusement. »\\
\zh{老师慢慢地解释这个语法点。} \emph{Lǎoshī mànmàn de jiěshì zhège yǔfǎ diǎn.} « Le professeur explique lentement ce point de grammaire. »
\end{exemplebox}

\begin{attention}[Piège 4 : 地 après le verbe]
Sous l'influence du français (« chanter joyeusement »), le francophone écrit \zh{他唱地高兴} ou \zh{她说话轻声地} : c'est \textbf{incorrect}. \zh{地} ne suit jamais le verbe ; il introduit le verbe : \zh{他高兴地唱。} Attention aussi à ne pas confondre avec \zh{得} : \zh{他唱得好} (il chante bien : évaluation \emph{après} le verbe) et \zh{他高兴地唱} (il chante joyeusement : manière \emph{avant} le verbe) sont tous deux corrects, mais ne veulent pas dire la même chose.
\end{attention}

\courssub{Schéma de correspondance français ↔ chinois}

\begin{popfigure}
\begin{tikzpicture}[node distance=1cm]
\node[draw,rounded corners,fill=poporangeL] (fr) {Français};
\node[draw,rounded corners,below=1.6cm of fr,fill=popblueL] (zh) {Chinois};
\draw[->,thick,popdark] (fr.west) to[bend right=15] node[left,align=center,font=\small] {pouvoir + inf.} (zh.west);
\draw[->,thick,popdark] (fr) -- node[right,align=center,font=\small] {nom + de} (zh);
\draw[->,thick,popdark] (fr.east) to[bend left=15] node[right,align=center,font=\small] {adj + -ment\\verbe + si/très + adj} (zh.east);
\node[below=0.4cm of zh,align=center,font=\small] {可以 + V \quad|\quad X + 的 + N \quad|\quad V + 得 + Adj \quad|\quad Adj + 地 + V};
\end{tikzpicture}
\legende{Chaque construction française correspond à un schéma chinois précis : traduire, c'est réordonner.}
\end{popfigure}

\begin{center}
\begin{tabularx}{\linewidth}{|X|X|X|X|}
\hline
\rowcolor{popblueL} Notion française & Particule & Schéma chinois & Exemple \\
\hline
pouvoir + infinitif & 可以 & Sujet + 可以 + Verbe & 我可以去。Je peux y aller. \\
\hline
possession « de » & 的 & Possesseur + 的 + Nom & 我的书。Mon livre. \\
\hline
relatif « qui/que » & 的 & Proposition + 的 + Nom & 我买的书。Le livre que j'ai acheté. \\
\hline
verbe + si/très + adj. & 得 & Verbe + 得 + Adjectif & 他跑得快。Il court vite. \\
\hline
adverbe en « -ment » & 地 & Adjectif + 地 + Verbe & 她轻声地说话。Elle parle doucement. \\
\hline
\end{tabularx}
\end{center}

\courssub{Piège 1 : 的、得、地 se prononcent tous « de »}

À l'oral, aucune différence : les trois particules se disent \emph{de} (ton neutre). À l'écrit, choisir le mauvais caractère est une faute. Le réflexe à acquérir :

\begin{itemize}
\item devant un \textbf{nom} → \zh{的} (\zh{我的老师}, mon professeur) ;
\item après un \textbf{verbe}, devant un adjectif de degré → \zh{得} (\zh{跑得快}, courir vite) ;
\item devant un \textbf{verbe}, après un adjectif de manière → \zh{地} (\zh{高兴地说}, dire joyeusement).
\end{itemize}

\begin{savaistu}[Le caractère 得 a une double vie]
\zh{得} existe aussi comme verbe plein, prononcé \emph{dé} (2\textsuperscript{e} ton), qui signifie « obtenir » : \zh{得到} (dédào, « obtenir »), \zh{他得了第一名} (Tā dé le dì-yī míng, « il a obtenu la première place »). Mais dans \zh{跑得快}, c'est une simple particule au ton neutre, sans sens propre. Même caractère, deux fonctions : le contexte décide.
\end{savaistu}

\begin{methode}[Traduire sans faute : la méthode en trois étapes]
\begin{enumerate}
\item \textbf{Traduire mentalement en français} et identifier la construction : possession ? degré après le verbe ? manière avant le verbe ? verbe modal ?
\item \textbf{Choisir la bonne particule} : nom → \zh{的} ; verbe + adjectif → \zh{得} ; adjectif + verbe → \zh{地} ; modal → \zh{可以} avant le verbe.
\item \textbf{Réordonner selon le schéma chinois}, jamais selon l'ordre français : possesseur + \zh{的} + nom ; adverbe + \zh{地} + verbe ; sujet + \zh{可以} + verbe + objet.
\end{enumerate}
\end{methode}

\begin{exoresolu}[Transformation : du français au chinois]
Traduisez en chinois : « Le professeur de mon école explique très clairement, et les élèves peuvent poser des questions. »
\tcblower
\textbf{Étape 1 — Découpage.} (a) « le professeur de mon école » : possession → \zh{的} ; (b) « explique très clairement » : manière avant le verbe → adjectif + \zh{地} + verbe ; (c) « peuvent poser » : modal → \zh{可以} + verbe.\\
\textbf{Étape 2 — Réordonnancement.} Mon école + \zh{的} + professeur → \zh{我们学校的老师} ; clairement + expliquer → \zh{很清楚地解释} ; élèves + pouvoir + poser des questions → \zh{学生可以问问题}.\\
\textbf{Phrase finale :} \zh{我们学校的老师解释得很清楚，学生可以问问题。}\\
\emph{Wǒmen xuéxiào de lǎoshī jiěshì de hěn qīngchu, xuéshēng kěyǐ wèn wèntí.}\\
Remarque : on pouvait aussi dire \zh{老师很清楚地解释} (manière, avec \zh{地}) ; la version avec \zh{得} insiste sur le résultat (« il explique de façon très claire »). Les deux sont correctes, mais les particules ne sont pas interchangeables.
\end{exoresolu}

\begin{exercice}[À vous de jouer]
Corrigez les phrases suivantes, toutes fautives : 1) \zh{我去可以商店。} 2) \zh{他跑快。} 3) \zh{她唱歌地高兴。} 4) \zh{书的我哥哥。} Puis traduisez : « Mon frère peut conduire, il conduit prudemment et il conduit très bien. »
\end{exercice}

\begin{corrige}[Exercice « À vous de jouer »]
1) \zh{我可以去商店。} (le modal \zh{可以} précède le verbe). 2) \zh{他跑得快。} (\zh{得} obligatoire devant le complément de degré). 3) \zh{她高兴地唱歌。} (\zh{地} se place avant le verbe, pas après). 4) \zh{我哥哥的书。} (possesseur + \zh{的} + nom). Traduction : \zh{我哥哥可以开车，他小心地开车，他开车开得很好。} \emph{Wǒ gēge kěyǐ kāichē, tā xiǎoxīn de kāichē, tā kāichē kāi de hěn hǎo.} Notez : quand le verbe a un objet (\zh{开车}), on répète le verbe devant \zh{得} : \zh{开车开得很好}. Pour « peut conduire » au sens de « sait conduire », \zh{会开车} serait aussi correct.
\end{corrige}

\begin{aretenir}[Les quatre réflexes anti-erreur]
\begin{itemize}
\item \zh{可以} est invariable, ne marque pas le temps, et se place \textbf{avant} le verbe, \textbf{après} le sujet ; négation : \zh{不可以}.
\item \zh{的} : devant un \textbf{nom} (possession, description, relatif) — ordre inversé par rapport au français « de ».
\item \zh{得} : \textbf{après le verbe}, devant l'adjectif de degré (« courir vite ») — jamais de juxtaposition verbe + adjectif ; verbe à objet : on répète le verbe (\zh{开车开得很好}).
\item \zh{地} : \textbf{avant le verbe}, après l'adjectif de manière (« -ment ») — position inverse du français.
\end{itemize}
\end{aretenir}

\coursec{Exercices d'application et de transformation}

Dans la section précédente, nous avons comparé le chinois et le français et repéré les pièges typiques des francophones : oubli du mot-outil, mauvais choix entre \zh{的}, \zh{得} et \zh{地}, place incorrecte de \zh{可以} dans la phrase. Il est temps de vérifier que tout est acquis. Cette section propose un parcours progressif : d'abord la \cle{complétion} (choisir le bon mot), puis la \cle{transformation} (affirmatif → négatif → interrogatif), puis la \cle{traduction} (français → chinois), et enfin la \cle{production libre} guidée. Chaque exercice est précédé d'un rappel de méthode et suivi d'un corrigé détaillé.

\begin{popfigure}
\begin{tikzpicture}
\node[draw,rounded corners,fill=popblueL,align=center] (ex) {Exercices\\de grammaire};
\node[draw,rounded corners,below left=1.2cm and 0.8cm of ex,fill=poporangeL,align=center] (c1) {Complétion\\(choisir le mot)};
\node[draw,rounded corners,below right=1.2cm and 0.8cm of ex,fill=popgreenL,align=center] (c2) {Transformation\\(négatif, question)};
\node[draw,rounded corners,below=2.4cm of ex,fill=poppurpleL,align=center] (c3) {Traduction et\\production guidée};
\draw[->,thick,popblue] (ex) -- (c1);
\draw[->,thick,popblue] (ex) -- (c2);
\draw[->,thick,popblue] (ex) -- (c3);
\end{tikzpicture}
\legende{Les trois grandes familles d'exercices de grammaire au Bac.}
\end{popfigure}

\courssub{Exercice 1 : compléter avec \zh{可以} / \zh{不可以} / \zh{可以…吗}}

\begin{textebox}[Exemple d'item type Bac]
Consigne : complétez la phrase avec le mot qui convient.\\
« Tu \_\_\_\_ (pouvoir) partir maintenant. »\\
Réponse : 你可以走了。\\
\emph{Nǐ kěyǐ zǒu le.}\\
« Tu peux partir maintenant. »
\end{textebox}

\begin{methode}[Choisir entre les trois formes]
\begin{enumerate}
\item La phrase est affirmative et exprime une permission accordée : utilisez \zh{可以} (\emph{kěyǐ}), placé \textbf{avant le verbe principal} : Sujet + 可以 + Verbe.
\item La phrase interdit quelque chose : utilisez \zh{不可以} (\emph{bù kěyǐ}) ; la négation \zh{不} se place \textbf{avant} le modal : Sujet + 不可以 + Verbe.
\item La phrase demande une permission : gardez \zh{可以} et ajoutez \zh{吗} (\emph{ma}) à la fin, avec le point d'interrogation chinois ？ : Sujet + 可以 + Verbe + 吗？
\end{enumerate}
\end{methode}

\begin{attention}[Sandhi tonal de \zh{可以}]
Vérifiez toujours l'accord tonal : \zh{可以} s'écrit \emph{kě yǐ} (3\textsuperscript{e} ton + 3\textsuperscript{e} ton), mais se prononce \emph{ké yǐ} par sandhi du troisième ton (deux troisièmes tons qui se suivent : le premier devient un deuxième ton). À l'écrit du Bac, on conserve les marques d'origine : \emph{kěyǐ}. Ne l'écrivez jamais \emph{kéyǐ} dans une copie.
\end{attention}

\begin{exercice}[Exercice 1 — À vous]
Complétez avec \zh{可以}, \zh{不可以} ou \zh{可以…吗} :
\begin{enumerate}
\item 教室里\_\_\_\_吃东西。（Interdiction）
\item 我\_\_\_\_用一下你的手机\_\_\_\_？（Demande de permission）
\item 考完试以后，你们\_\_\_\_回家。（Permission）
\item 上课的时候\_\_\_\_说话。（Interdiction）
\item 妈妈，我\_\_\_\_看电视\_\_\_\_？（Demande de permission）
\end{enumerate}
\end{exercice}

\begin{corrige}[Exercice 1]
\begin{enumerate}
\item 教室里不可以吃东西。\emph{Jiàoshì lǐ bù kěyǐ chī dōngxi.} « On ne peut pas manger dans la salle de classe. »
\item 我可以用一下你的手机吗？\emph{Wǒ kěyǐ yòng yíxià nǐ de shǒujī ma ?} « Puis-je utiliser ton téléphone un moment ? »
\item 考完试以后，你们可以回家。\emph{Kǎo wán shì yǐhòu, nǐmen kěyǐ huí jiā.} « Après l'examen, vous pouvez rentrer à la maison. »
\item 上课的时候不可以说话。\emph{Shàngkè de shíhou bù kěyǐ shuōhuà.} « Pendant le cours, on ne peut pas parler. »
\item 妈妈，我可以看电视吗？\emph{Māma, wǒ kěyǐ kàn diànshì ma ?} « Maman, est-ce que je peux regarder la télévision ? »
\end{enumerate}
\end{corrige}

\courssub{Exercice 2 : choisir la bonne particule — \zh{的}, \zh{得} ou \zh{地}}

Avant l'exercice, revoici le tableau de décision qui sert de grille de lecture.

\begin{center}
\begin{tabularx}{\linewidth}{|X|X|X|X|}
\hline
\rowcolor{popblueL} Particule & Pinyin & Fonction & Schéma \\
\hline
的 & de & attribut, possession & Nom/Adj + 的 + Nom \\
\hline
得 & de & complément de degré & Verbe + 得 + Adj \\
\hline
地 & de & complément de manière & Adj + 地 + Verbe \\
\hline
\end{tabularx}
\end{center}

\begin{methode}[Trouver la bonne particule en trois questions]
\begin{enumerate}
\item Le mot qui suit est-il un \textbf{nom} ? → \zh{的} (这是谁的笔？).
\item Le mot qui précède est-il un \textbf{verbe} et celui qui suit un \textbf{adjectif} ? → \zh{得} (他做得真好。).
\item Le mot qui précède est-il un \textbf{adjectif} et celui qui suit un \textbf{verbe} ? → \zh{地} (她开心地跳。).
\end{enumerate}
\end{methode}

\begin{exemplebox}[Les trois particules en action]
\begin{enumerate}
\item 我可以吃吗？\emph{Wǒ kěyǐ chī ma ?} « Est-ce que je peux manger ? » (modal + question)
\item 这是谁的笔？\emph{Zhè shì shéi de bǐ ?} « À qui est ce stylo ? » (possession)
\item 他做得真好。\emph{Tā zuò de zhēn hǎo.} « Il le fait vraiment bien. » (degré)
\item 她开心地跳。\emph{Tā kāixīn de tiào.} « Elle saute joyeusement. » (manière)
\item 我们不可以迟到。\emph{Wǒmen bù kěyǐ chídào.} « Nous ne pouvons pas être en retard. » (interdiction)
\end{enumerate}
\end{exemplebox}

\begin{exoresolu}[Exercice 2 — Dix phrases à trous]
Énoncé : complétez avec \zh{的}, \zh{得} ou \zh{地}.\\
1) 这是老师\_\_\_\_书。 2) 他跑\_\_\_\_很快。 3) 孩子们高兴\_\_\_\_唱歌。 4) 我\_\_\_\_家在雅温得。 5) 她说\_\_\_\_不错。 6) 他慢慢\_\_\_\_走。 7) 红\_\_\_\_花很漂亮。 8) 他汉语说\_\_\_\_真好。 9) 我们认真\_\_\_\_学习。 10) 这是谁\_\_\_\_电脑？
\tcblower
Solution détaillée :
\begin{enumerate}
\item 这是老师\textbf{的}书。\emph{Zhè shì lǎoshī de shū.} — avant le nom \zh{书} → possession.
\item 他跑\textbf{得}很快。\emph{Tā pǎo de hěn kuài.} — verbe \zh{跑} + adjectif \zh{快} → degré.
\item 孩子们高兴\textbf{地}唱歌。\emph{Háizimen gāoxìng de chànggē.} — adjectif \zh{高兴} + verbe \zh{唱歌} → manière.
\item 我\textbf{的}家在雅温得。\emph{Wǒ de jiā zài Yǎwēndé.} — possession devant \zh{家}.
\item 她说\textbf{得}不错。\emph{Tā shuō de búcuò.} — verbe + adjectif → degré.
\item 他慢慢\textbf{地}走。\emph{Tā mànmàn de zǒu.} — manière devant le verbe \zh{走}.
\item 红\textbf{的}花很漂亮。\emph{Hóng de huā hěn piàoliang.} — attribut devant le nom \zh{花}.
\item 他汉语说\textbf{得}真好。\emph{Tā Hànyǔ shuō de zhēn hǎo.} — attention : avec un complément d'objet (\zh{汉语}), le verbe se répète avant \zh{得}.
\item 我们认真\textbf{地}学习。\emph{Wǒmen rènzhēn de xuéxí.} — manière.
\item 这是谁\textbf{的}电脑？\emph{Zhè shì shéi de diànnǎo ?} — possession interrogative.
\end{enumerate}
\end{exoresolu}

\begin{attention}[Piège fréquent n°1 : tout écrire « de »]
Les trois particules se prononcent toutes \emph{de} (ton neutre), mais s'écrivent différemment. À l'écrit, confondre \zh{的}, \zh{得} et \zh{地} est une faute de caractère sanctionnée au Bac. Astuce : \zh{的} contient la clé 白 (blanc), \zh{得} contient 彳 (pas, marcher → résultat de l'action), \zh{地} contient 土 (terre → le « sol » sur lequel l'action se déroule, la manière).
\end{attention}

\courssub{Exercice 3 : transformer avec \zh{可以}}

\begin{methode}[Transformation affirmative → négative → interrogative]
\begin{enumerate}
\item \textbf{Gardez le sujet} et tout le reste de la phrase dans le même ordre.
\item Pour la \textbf{négation}, ajoutez \zh{不} \textbf{avant} \zh{可以} : Sujet + 不可以 + Verbe.
\item Pour la \textbf{question}, gardez la forme affirmative et ajoutez \zh{吗} à la fin : Sujet + 可以 + Verbe + 吗？
\item La réponse courte : 可以。/ 不可以。 (jamais un simple « oui » isolé comme en français).
\end{enumerate}
\end{methode}

\begin{exercice}[Exercice 3 — À vous]
Transformez chaque phrase en phrase négative, puis en question :
\begin{enumerate}
\item 你可以进来。
\item 学生可以带手机。
\item 我们可以去杜阿拉。
\item 他可以请假。
\item 你们可以用词典。
\end{enumerate}
\end{exercice}

\begin{corrige}[Exercice 3]
\begin{enumerate}
\item 你不可以进来。\emph{Nǐ bù kěyǐ jìnlái.} « Tu ne peux pas entrer. » / 你可以进来吗？\emph{Nǐ kěyǐ jìnlái ma ?} « Puis-je entrer ? »
\item 学生不可以带手机。\emph{Xuésheng bù kěyǐ dài shǒujī.} « Les élèves ne peuvent pas apporter leur téléphone. » / 学生可以带手机吗？\emph{Xuésheng kěyǐ dài shǒujī ma ?} « Les élèves peuvent-ils apporter leur téléphone ? »
\item 我们不可以去杜阿拉。\emph{Wǒmen bù kěyǐ qù Dù'ālā.} « Nous ne pouvons pas aller à Douala. » / 我们可以去杜阿拉吗？\emph{Wǒmen kěyǐ qù Dù'ālā ma ?} « Pouvons-nous aller à Douala ? »
\item 他不可以请假。\emph{Tā bù kěyǐ qǐngjià.} « Il ne peut pas demander un congé. » / 他可以请假吗？\emph{Tā kěyǐ qǐngjià ma ?} « Peut-il demander un congé ? »
\item 你们不可以用词典。\emph{Nǐmen bù kěyǐ yòng cídiǎn.} « Vous ne pouvez pas utiliser le dictionnaire. » / 你们可以用词典吗？\emph{Nǐmen kěyǐ yòng cídiǎn ma ?} « Pouvez-vous utiliser le dictionnaire ? »
\end{enumerate}
\end{corrige}

\begin{attention}[Piège fréquent n°2 : la place de \zh{不}]
Ne dites jamais \zh{可以不} pour « ne pas pouvoir » : la négation porte sur le modal, donc \zh{不} précède \zh{可以}. Comparez : \zh{他不可以去。} \emph{Tā bù kěyǐ qù.} « Il ne peut pas y aller » (interdiction) et \zh{他可以不去。} \emph{Tā kěyǐ bú qù.} « Il peut ne pas y aller » (il a le droit de ne pas y aller). La place de \zh{不} change le sens !
\end{attention}

\courssub{Exercice 4 : de l'adverbe français en « -ment » à la structure Adj + \zh{地} + V}

En français, « lentement », « joyeusement », « sérieusement » se construisent avec le suffixe \emph{-ment}. En chinois, c'est la particule \zh{地} qui joue ce rôle : on place l'adjectif \textbf{avant} le verbe, suivi de \zh{地}.

\begin{center}
\begin{tabularx}{\linewidth}{|X|X|X|}
\hline
\rowcolor{popblueL} Français & Chinois (caractères + pinyin) & Traduction littérale \\
\hline
parler lentement & 慢慢地说 \emph{mànmàn de shuō} & « lent + 地 + parler » \\
\hline
travailler sérieusement & 认真地工作 \emph{rènzhēn de gōngzuò} & « sérieux + 地 + travailler » \\
\hline
écouter attentivement & 注意地听 \emph{zhùyì de tīng} & « attentif + 地 + écouter » \\
\hline
\end{tabularx}
\end{center}

\begin{exercice}[Exercice 4 — À vous]
Réécrivez en chinois avec la structure Adj + \zh{地} + V :
\begin{enumerate}
\item Elle chante joyeusement. (高兴 / 唱歌)
\item Le professeur explique patiemment. (耐心 / 解释)
\item Nous étudions sérieusement. (认真 / 学习)
\item Il conduit prudemment. (小心 / 开车)
\item Les élèves écoutent attentivement. (注意 / 听)
\end{enumerate}
\end{exercice}

\begin{corrige}[Exercice 4]
\begin{enumerate}
\item 她高兴地唱歌。\emph{Tā gāoxìng de chànggē.} « Elle chante joyeusement. »
\item 老师耐心地解释。\emph{Lǎoshī nàixīn de jiěshì.} « Le professeur explique patiemment. »
\item 我们认真地学习。\emph{Wǒmen rènzhēn de xuéxí.} « Nous étudions sérieusement. »
\item 他小心地开车。\emph{Tā xiǎoxīn de kāichē.} « Il conduit prudemment. »
\item 同学们注意地听。\emph{Tóngxuémen zhùyì de tīng.} « Les élèves écoutent attentivement. »
\end{enumerate}
Remarque : contrairement au français où l'adverbe suit souvent le verbe, en chinois le complément de manière \textbf{précède toujours} le verbe.
\end{corrige}

\courssub{Exercice 5 : traduction du français vers le chinois}

\begin{exercice}[Exercice 5 — À vous]
Traduisez en chinois (permission, attribut, degré, manière) :
\begin{enumerate}
\item Est-ce que je peux poser une question ? (permission)
\item Le livre de mon ami est sur la table. (attribut)
\item Elle parle chinois très bien. (degré)
\item Il répond poliment au professeur. (manière)
\item On ne peut pas fumer à l'école. (interdiction)
\end{enumerate}
\end{exercice}

\begin{corrige}[Exercice 5]
\begin{enumerate}
\item 我可以问一个问题吗？\emph{Wǒ kěyǐ wèn yí ge wèntí ma ?} — classificateur \zh{个} devant \zh{问题}.
\item 我朋友的书在桌子上。\emph{Wǒ péngyou de shū zài zhuōzi shàng.} — \zh{的} de possession ; lieu exprimé par \zh{在} + nom + \zh{上}.
\item 她汉语说得很好。\emph{Tā Hànyǔ shuō de hěn hǎo.} — verbe répété à cause du complément d'objet \zh{汉语}.
\item 他有礼貌地回答老师。\emph{Tā yǒu lǐmào de huídá lǎoshī.} — manière : \zh{有礼貌} + \zh{地} + verbe.
\item 在学校不可以抽烟。\emph{Zài xuéxiào bù kěyǐ chōuyān.} — complément de lieu en tête, puis \zh{不可以}.
\end{enumerate}
\end{corrige}

\begin{attention}[Piège fréquent n°3 : calquer l'ordre français]
Pour « Elle parle chinois très bien », le français place « très bien » à la fin ; le chinois exige la structure \zh{汉语说得很好} avec répétition du verbe : Sujet + Objet + Verbe + 得 + Adjectif. La forme \zh{她说汉语很好} est compréhensible mais ne teste pas la structure exigée au Bac.
\end{attention}

\courssub{Exercice 6 : production guidée — un dialogue en six répliques}

\begin{experience}[Rédiger un mini-dialogue]
Par groupes de deux, écrivez un dialogue de \textbf{six répliques} entre un élève et un professeur (ou un élève et sa mère) qui contient :
\begin{itemize}
\item au moins une question avec \zh{可以…吗} ;
\item au moins une interdiction avec \zh{不可以} ;
\item au moins une particule \zh{的} (possession) ;
\item au moins une particule \zh{得} (degré) ;
\item au moins une particule \zh{地} (manière).
\end{itemize}
Puis jouez le dialogue devant la classe en soignant les tons.
\end{experience}

\begin{textebox}[Modèle de dialogue corrigé]
学生：老师，我可以问一个问题吗？\\
\emph{Xuésheng : Lǎoshī, wǒ kěyǐ wèn yí ge wèntí ma ?}\\
« Élève : Monsieur, puis-je poser une question ? »\\
老师：可以，但是上课的时候不可以用手机。\\
\emph{Lǎoshī : Kěyǐ, dànshì shàngkè de shíhou bù kěyǐ yòng shǒujī.}\\
« Professeur : Oui, mais pendant le cours on ne peut pas utiliser le téléphone. »\\
学生：好的。老师，您的汉语说得真好！\\
\emph{Xuésheng : Hǎo de. Lǎoshī, nín de Hànyǔ shuō de zhēn hǎo !}\\
« Élève : D'accord. Monsieur, vous parlez vraiment bien chinois ! »\\
老师：谢谢！你也很认真地学习。\\
\emph{Lǎoshī : Xièxie ! Nǐ yě hěn rènzhēn de xuéxí.}\\
« Professeur : Merci ! Toi aussi, tu étudies très sérieusement. »\\
学生：这是我从中国朋友那儿借的书。\\
\emph{Xuésheng : Zhè shì wǒ cóng Zhōngguó péngyou nàr jiè de shū.}\\
« Élève : Voici le livre que j'ai emprunté à mon ami chinois. »\\
老师：很好！你读得越来越多了。\\
\emph{Lǎoshī : Hěn hǎo ! Nǐ dú de yuè lái yuè duō le.}\\
« Professeur : Très bien ! Tu lis de plus en plus. »
\end{textebox}

\courssub{Auto-évaluation : ma grille en dix points}

\begin{definition}[Grille d'auto-évaluation (10 items)]
Après chaque production, attribuez-vous 1 point par critère maîtrisé :
\begin{enumerate}
\item Structure correcte : \zh{可以} placé avant le verbe principal.
\item Négation correcte : \zh{不} placé avant \zh{可以}.
\item Question correcte : \zh{吗} en fin de phrase, avec ？
\item Particule adéquate : \zh{的} devant un nom.
\item Particule adéquate : \zh{得} entre verbe et adjectif.
\item Particule adéquate : \zh{地} entre adjectif et verbe.
\item Ordre des mots : complément de manière avant le verbe.
\item Ponctuation chinoise pleine chasse : ，。？ utilisés correctement.
\item Tonalité : \emph{kěyǐ} prononcé \emph{ké yǐ} (sandhi), particules \zh{的}, \zh{得}, \zh{地} au ton neutre.
\item Caractères : traits corrects, aucune confusion 的 / 得 / 地 à l'écrit.
\end{enumerate}
Objectif : 8/10 minimum avant l'examen.
\end{definition}

\begin{savaistu}[Les particules au Baccalauréat]
Au Bac, les items de grammaire représentent une part importante de la note globale de l'épreuve de chinois. La maîtrise des particules \zh{的}, \zh{得} et \zh{地} est un critère discriminateur : les correcteurs distinguent immédiatement les copies qui les emploient correctement de celles qui les confondent. Un entraînement régulier par petites séries de dix phrases à trous, comme l'exercice 2, est la méthode la plus efficace pour progresser.
\end{savaistu}

\begin{aretenir}[Bilan de la section]
\begin{itemize}
\item \zh{可以} : permission ; négation \zh{不可以} ; question \zh{可以…吗}.
\item \zh{的} + nom (attribut, possession) ; verbe + \zh{得} + adjectif (degré) ; adjectif + \zh{地} + verbe (manière).
\item Les trois particules se prononcent \emph{de} (ton neutre) mais s'écrivent différemment : vigilance maximale à l'écrit.
\item Méthode de transformation : garder le sujet, \zh{不} avant le modal, \zh{吗} à la fin.
\item Auto-évaluez chaque production avec la grille en dix points.
\end{itemize}
\end{aretenir}

\coursec{Compléments utiles à l'examen (bonus)}

Après avoir consolidé les structures de base à travers les exercices d'application et de transformation de la section précédente, nous abordons maintenant les points de vigilance et les stratégies spécifiques qui font la différence lors de l'épreuve écrite du Baccalauréat (Module 5 : Mass médias et télécommunication). Cette section rassemble les formes contractées du registre formel, les pièges de la dictée sur les trois \cle{de}, l'analyse d'un article de presse type et une méthode mnémotechnique infaillible.

\courssub{Formes contractées : 可以 → 可 (kě) à l'écrit littéraire et administratif}

Dans les textes officiels, les panneaux d'affichage, les titres de presse ou les règlements intérieurs (très fréquents dans le Module 5), le verbe modal \zh{可以} (pouvoir, être autorisé à) se contracte souvent en son premier caractère \zh{可} (kě). Le sens reste la permission ou la possibilité, mais le registre devient formel, concis, voire impératif.

\begin{textebox}[Exemples de formes contractées dans l'espace public]
\zh{可 以 入 内} (Kěyǐ rùnèi) — « Entrée autorisée » (Panneau). \\
\zh{未 成年 者 可 以 独自 入场} (Wèichéngniánzhě kěyǐ dúzì rùchǎng) → \zh{未成年者可独自入场} (Les mineurs peuvent entrer seuls). \\
\zh{本 规定 可 以 修改} (Běn guīdìng kěyǐ xiūgǎi) — « Ce règlement peut être modifié. »
\end{textebox}

\begin{propriete}[Règle : 可以  vs  可]
\emph{Structure :} Sujet + \textbf{可} + Verbe (souvent sans objet explicite ou avec objet court). \\
\emph{Registre :} Écrit formel, administratif, journalistique, panneaux, contrats. \\
\emph{Emploi :} Exprime la permission (\emph{may}) ou la possibilité objective (\emph{can be}). \\
\emph{Exception :} À l'oral et dans l'écrit courant (messages, dialogues, rédaction personnelle), on utilise \textbf{toujours} \zh{可以} (kěyǐ). Ne dites jamais \zh{我可去} à l'oral.
\end{propriete}

\begin{attention}[Piège : 可以 (permission) vs 可 (possible / littéral)]
Ne confondez pas :
\begin{itemize}
    \item \zh{可以} (kěyǐ) : Verbe modal complet. \zh{我可以去。} (Je peux y aller / J'ai la permission).
    \item \zh{可} (kě) : Adverbe / Verbe auxiliaire littéral signifiant « peut / possible / digne de ». S'utilise souvent dans des constructions figées : \zh{可爱} (kě'ài, mignon, litt. « digne d'amour »), \zh{可怜} (kělián, pitoyable), \zh{可行} (kěxíng, faisable).
    \item Dans un sujet d'examen, si vous voyez \zh{违者可罚款} (Les contrevenants s'exposent à une amende), \zh{可} = \zh{可以} (permission/possibilité légale). Si vous voyez \zh{这件事可大可小} (Cette affaire peut être grave ou mineure), \zh{可} = « peut / être susceptible de ».
\end{itemize}
\end{attention}

\begin{exemplebox}[Exemples traduits : 可以  vs  可]
\zh{可 以 说 是 最好 的。} \\
\emph{Kěyǐ shuō shì zuìhǎo de.} \\
« On peut dire que c'est le meilleur. » (Oral / Écrit standard)

\zh{违规 者 可 罚款 五百 元。} \\
\emph{Wéiguī zhě kě fákuǎn wǔbǎi yuán.} \\
« Les contrevenants peuvent être amendés de 500 yuans. » (Règlement / Affiche officielle)

\zh{这 个 方案 可 行。} \\
\emph{Zhège fāng'àn kěxíng.} \\
« Ce plan est faisable. » (\zh{可} = adverbe « possible », ne se substitue pas à \zh{可以} ici).
\end{exemplebox}

\courssub{可以 dans les annonces officielles et la presse (Module 5 : Mass médias)}

Le Module 5 impose la compréhension de textes journalistiques. Les annonces de service public, les éditoriaux et les communiqués utilisent massivement \zh{可以} (ou \zh{可}) pour exprimer les \textbf{droits des citoyens}, les \textbf{possibilités offertes par les nouvelles technologies} ou les \textbf{conséquences légales}.

\begin{textebox}[Article de presse simulé : « La 5G change la vie à Yaoundé »]
\zh{【雅温得日报】五 G 网络 正式 商用，市民 可以 享受 极速 体验} \\
\zh{雅温得 讯} （记者 李 华）\zh{昨天}，\zh{喀麦隆} 首都 \zh{雅温得} 正式 开通 \zh{五 G} 商用 网络。\zh{据} 运营商 介绍，\zh{用户} \zh{只需} 更换 \zh{支持 五 G 的 手机}，\zh{即可} 享受 \zh{下载 速度 超过} 1\zh{Gbps} \zh{的 极速 网络}。 \\

\zh{专家 指出}，\zh{五 G 的 低延迟 特性} \zh{可以} \zh{大大 提升} 远程 医疗 \zh{和} 在线 教育 \zh{的 质量}。\zh{对于} 偏远 地区 \zh{的} 学生，\zh{这 意味着} 他们 \zh{可以} \zh{实时 观看} 北京 或 巴黎 \zh{的 名校 课程}。\zh{此外}，\zh{智慧 城市} 项目 \zh{也 可以} 更 好 地 监控 交通 流量，\zh{减少} 拥堵。 \\

\zh{不过}，\zh{网络 安全 专家 提醒}：\zh{个人 信息 保护} \zh{不容忽视}。\zh{用户 在 享受 便利 的 同时}，\zh{应当 注意 设置 隐私 权限}，\zh{避免 数据 泄露}。\zh{相关 部门 表示}，\zh{将 严格 执法}，\zh{违规 收集 用户 数据 者}，\zh{可 处 以 重罚}。 \\
\emph{[Yǎwēndé Rìbào] Wǔ G wǎngluò zhèngshì shāngyòng, shímín kěyǐ xiǎngshòu jí sù tǐyàn. Yǎwēndé xùn (jìzhě Lǐ Huá) zuótiān, Kāmàilóng shǒudū Yǎwēndé zhèngshì kāitōng Wǔ G shāngyòng wǎngluò. Jù yùnxiāngshāng jièshào, yònghù zhǐxū gēnghuàn zhīchí Wǔ G de shǒujī, jí kě xiǎngshòu xiàzài sùdù chāoguò 1 Gbps de jí sù wǎngluò. Zhuānjiā zhǐchū, Wǔ G de dī yánchí tèxìng kěyǐ dàdà tíshēng yuǎnchéng yīliáo hé zàixiàn jiàoyù de zhìliàng. Duìyú piānyuǎn dìqū de xuéshēng, zhè yìwèizhe tāmen kěyǐ shíshí guānkàn Běijīng huò Bālí de míngxiào kèchéng. Cǐwài, zhìhuì chéngshì xiàngmù yě kěyǐ gèng hǎo de jiānkòng jiāotōng liúliàng, jiǎnshǎo yōngdǔ. Bùguò, wǎngluò ānquán zhuānjiā tíxǐng: gèrén xìnxī bǎohù bùrónghūshì. Yònghù zài xiǎngshòu biànlì de tóngshí, yīngdāng zhùyì shèzhì yǐnsī quánxiàn, bìmiǎn shùjù xièlòu. Xiāngguān bùmén biǎoshì, jiāng yánzhí zhífǎ, wéiguī shōují yònghù shùjù zhě, kě chǔ yǐ zhòngfá.}
\end{textebox}

\begin{center}
\begin{tabularx}{\linewidth}{|X|X|X|X|}
\hline
\rowcolor{popblueL} \textbf{Caractères} & \textbf{Pinyin} & \textbf{Nature} & \textbf{Traduction} \\
\hline
商用 & shāngyòng & V & Exploiter commercialement / Mise en service commerciale \\
\hline
运营商 & yùnxiāngshāng & N & Opérateur (télécom) \\
\hline
即可 & jí kě & Adv + V & Aussitôt / dès lors peut (formel : 就可以) \\
\hline
低延迟 & dī yánchí & N & Faible latence \\
\hline
偏远 & piānyuǎn & Adj & Éloigné, reculé \\
\hline
实时 & shíshí & Adv & En temps réel \\
\hline
智慧城市 & zhìhuì chéngshì & N & Ville intelligente \\
\hline
监控 & jiānkòng & V & Surveiller, monitorer \\
\hline
不容忽视 & bùróng hūshì & Idiome & Ne pas négliger / Impératif de ne pas ignorer \\
\hline
隐私权限 & yǐnsī quánxiàn & N & Permissions de confidentialité \\
\hline
违规 & wéiguī & V/Adj & Enfreindre les règles / Irrégulier \\
\hline
可处以重罚 & kě chǔ yǐ zhòngfá & Structure & Peuvent écoper de lourdes amendes (formel) \\
\hline
\end{tabularx}
\end{center}

\begin{exercice}[Questions sur le texte (Compréhension écrite type Bac)]
\begin{enumerate}
    \item Relevez deux phrases où \zh{可以} exprime une \textbf{possibilité offerte par la technologie} (pas une permission administrative).
    \item Trouvez l'équivalent contracté de \zh{可以} dans la dernière phrase. Quelle en est la nuance ?
    \item Identifiez les trois particules \zh{的}、\zh{得}、\zh{地} dans le texte et justifiez leur emploi grammatical pour chacune.
\end{enumerate}
\end{exercice}

\courssub{Maîtriser les trois « de » (的、得、地) en dictée et en traduction}

C'est le point noir n°1 des francophones à l'écrit. En dictée, vous n'entendez qu'un seul son \emph{/də/} (ton neutre). La distinction se fait \textbf{uniquement par l'analyse syntaxique du contexte}.

\begin{methode}[Mnémotechnique « D - D - D » : Associer chaque particule à un mot-clé français]
\begin{enumerate}
    \item \textbf{的 (de) $\rightarrow$ « D » de \cle{Déterminant} / \cle{De} (possessif/attribut).} \\
    \emph{Règle :} Se place \textbf{APRÈS} un adjectif, un nom, un pronom ou une proposition pour qualifier un \textbf{NOM} qui suit. \\
    \emph{Test :} On peut souvent le remplacer par « \emph{le/la/lesquel(le)} » ou « \emph{de} » (possession).
    \item \textbf{得 (de) $\rightarrow$ « D » de \cle{Degré} / \cle{D'après le verbe}.} \\
    \emph{Règle :} Se place \textbf{APRÈS un VERBE (ou adjectif verbal)} pour introduire un \textbf{COMPLÉMENT DE DEGRÉ / RÉSULTAT / MANIÈRE INTENSIVE}. \\
    \emph{Test :} Le verbe précède, l'évaluation suit. Question : « Comment ? » / « À quel point ? ».
    \item \textbf{地 (de) $\rightarrow$ « D » de \cle{Devant le verbe} / \cle{Devant l'action}.} \\
    \emph{Règle :} Se place \textbf{AVANT un VERBE} (souvent après un adjectif + \zh{地}) pour indiquer la \textbf{MANIÈRE} dont l'action se déroule. \\
    \emph{Test :} Il modifie le verbe qui suit. Question : « Comment fait-il ? ».
\end{enumerate}
\end{methode}

\begin{center}
\begin{tabularx}{\linewidth}{|X|X|X|X|}
\hline
\rowcolor{poporangeL} \textbf{Particule} & \textbf{Position} & \textbf{Fonction} & \textbf{Exemple type} \\
\hline
\textbf{的} (de) & [Modificateur] + \textbf{的} + \textbf{[NOM]} & Attribut / Possession & \zh{红色的苹果} (pomme rouge), \zh{我的书} (mon livre) \\
\hline
\textbf{得} (de) & \textbf{[VERBE/ADJ]} + \textbf{得} + [Complément degré] & Degré / Résultat / Manière intensive & \zh{跑得快} (court vite), \zh{累得睡着了} (fatigué au point de s'endormir) \\
\hline
\textbf{地} (de) & [Adjectif] + \textbf{地} + \textbf{[VERBE]} & Manière / Circonstance & \zh{开心地笑} (rire joyeusement), \zh{认真地听} (écouter attentivement) \\
\hline
\end{tabularx}
\end{center}

\begin{savaistu}[Le caractère 得 (de/dé/dei) : 11 traits, radical 彳]
Le caractère \zh{得} s'écrit en \textbf{11 traits}. Son radical (clé) est \zh{彳} (chì, « pas », « petit pas »), qui évoque l'idée de marche, d'action, de progression. C'est un indice visuel : \zh{得} suit souvent une \textbf{action} (verbe) pour en mesurer le « pas », le degré.
\emph{Ordre des traits :} 1.\zh{撇}(piě) 2.\zh{撇}(piě) 3.\zh{竖}(shù) 4.\zh{横}(héng) 5.\zh{横}(héng) 6.\zh{竖}(shù) 7.\zh{提}(tí) 8.\zh{横}(héng) 9.\zh{撇}(piě) 10.\zh{捺}(nà) 11.\zh{点}(diǎn).
Attention : \zh{得} a trois prononciations : \zh{de} (particule), \zh{dé} (obtenir, \zh{得到}), \zh{děi} (devoir, \zh{得不到}).
\end{savaistu}

\begin{attention}[Pièges de dictée : Le contexte décide]
\begin{itemize}
    \item \zh{他 唱 得 好} (Tā chàng \textbf{de} hǎo) : \zh{唱} est un verbe $\rightarrow$ \textbf{得} (degré).
    \item \zh{他 的 歌 唱 得 好} (Tā \textbf{de} gē chàng de hǎo) : 1er \zh{的} (possessif, avant nom \zh{歌}) ; 2e \zh{得} (après verbe \zh{唱}).
    \item \zh{他 高高兴兴 地 唱 歌} (Tā gāoxìngxìng \textbf{de} chàng gē) : \zh{地} avant verbe \zh{唱} (manière).
    \item \zh{开心 的 歌} (Kāixīn \textbf{de} gē) : \zh{的} avant nom \zh{歌} (attribut).
    \item \zh{开心 地 唱} (Kāixīn \textbf{de} chàng) : \zh{地} avant verbe \zh{唱} (manière).
\end{itemize}
\emph{Astuce d'examen :} Si le mot qui suit est un \textbf{NOM} $\rightarrow$ \zh{的}. Si le mot qui \textbf{PRÉCÈDE} est un \textbf{VERBE/ADJ} et que ce qui suit évalue l'action $\rightarrow$ \zh{得}. Si le mot qui \textbf{SUIT} est un \textbf{VERBE} d'action $\rightarrow$ \zh{地}.
\end{attention}

\courssub{Stratégies de repérage dans les sujets d'examen (Bac Blanc / Bac)}

Les correcteurs du MINESEC glissent souvent des indices lexicaux pour guider le choix de la structure.

\begin{propriete}[Indices contextuels au Bac]
\begin{itemize}
    \item \textbf{Mot \zh{请} (qǐng, « veuillez / s'il vous plaît ») :} Presque toujours suivi de \zh{可以} / \zh{能} / \zh{请} + V. Ex: \zh{请 可以 在这里 停车} (incorrect, pléonasme) $\rightarrow$ \zh{请 在这里 停车} ou \zh{可以 在这里 停车}. Mais \zh{请问 可以 进来 吗 ?} (Puis-je entrer ?).
    \item \textbf{Mot \zh{违规} (wéiguī), \zh{违者} (wéizhě), \zh{者} (zhě) :} Signale un texte règlementaire $\rightarrow$ attend \zh{可} (contracté) ou \zh{将} (jiāng) + \zh{被} / \zh{处以}.
    \item \textbf{Mot \zh{应该} (yīnggāi) / \zh{必须} (bìxū) :} Obligation $\rightarrow$ exclut \zh{可以} (permission).
    \item \textbf{Adjectif bisyllabique + \_\_\_ + Verbe :} Trou à remplir $\rightarrow$ \zh{地}. Ex: \zh{认真 \_\_\_ 学习} $\rightarrow$ \zh{地}.
    \item \textbf{Verbe + \_\_\_ + Adjectif (rapide, lent, bien, mal) :} Trou à remplir $\rightarrow$ \zh{得}. Ex: \zh{跑 \_\_\_ 快} $\rightarrow$ \zh{得}.
    \item \textbf{Nom / Pronom + \_\_\_ + Nom :} Trou à remplir $\rightarrow$ \zh{的}. Ex: \zh{我 \_\_\_ 书} $\rightarrow$ \zh{的}.
\end{itemize}
\end{propriete}

\begin{exoresolu}[Exercice type Bac : Compléter avec 可以 / 可 / 的 / 得 / 地]
\textbf{Énoncé :} Complétez le texte suivant.
\zh{《图书馆 管理 规定》} \\
\zh{第一 条}：\zh{本 馆 读者 \_\_\_\_\_\_ 凭 证 入馆。} \\
\zh{第二 条}：\zh{请 保持 安静，\_\_\_\_\_\_ 大声 喧哗。} \\
\zh{第三 条}：\zh{书籍 \_\_\_\_\_\_ 整齐地 摆放，\_\_\_\_\_\_ 方便 他人 查阅。} \\
\zh{第四 条}：\zh{损坏 书籍 者，\_\_\_\_\_\_ 按 价 赔偿。} \\

\textbf{Solution détaillée :}
\begin{enumerate}
    \item \zh{可以} (ou \zh{可} pour le style très formel « Règlement »). \emph{Contexte : Permission d'entrée. Sujet « lecteurs ». Registre officiel $\rightarrow$ \zh{可} est élégant, \zh{可以} correct.}
    \item \zh{不可以} / \zh{不可} / \zh{请勿}. \emph{Interdiction. \zh{请勿} (qǐngwù) est le standard des panneaux. \zh{不可} (bùkě) forme contractée littéraire.}
    \item 1er trou : \zh{应当} / \zh{必须} / \zh{要}. \emph{Contexte : Obligation règlementaire. Structure : Sujet + Modalité + Manière (整齐地) + Verbe (摆放).}
    \item 2e trou : \zh{以便} (yǐbiàn) / \zh{从而} (cóng'ér). \emph{Piège : « 方便他人查阅 » exprime une finalité. On n'utilise pas \zh{得} ni \zh{的} ni \zh{地} ici. Structure : V + \zh{以便} / \zh{从而} + V.}
    \item \zh{可} / \zh{可以} / \zh{应当} / \zh{须}. \emph{Registre règlementaire : \zh{损坏书籍者，可按价赔偿} (Les endommageurs peuvent/doivent indemniser selon le prix). \zh{可} exprime ici la conséquence légale (permission/possibilité juridique d'exiger réparation).}
\end{enumerate}
\end{exoresolu}

\begin{popfigure}
\begin{tikzpicture}[node distance=1.2cm, every node/.style={font=\small}]
\node[draw=popblue, thick, rounded corners, fill=popblueL, minimum width=3cm, minimum height=1cm] (bonus) {\textbf{Bonus Examen}};
\node[draw=poporange, thick, rounded corners, fill=poporangeL, minimum width=3.5cm, minimum height=1cm, below=of bonus] (mnemo) {\textbf{Mnémotechnique D-D-D}};
\node[draw=popgreen, thick, rounded corners, fill=popgreenL, minimum width=3cm, minimum height=1cm, left=of mnemo, xshift=-1cm] (d1) {\zh{的} = \cle{Déterminant} (Nom)};
\node[draw=poppurple, thick, rounded corners, fill=poppurpleL, minimum width=3cm, minimum height=1cm, right=of mnemo, xshift=1cm] (d2) {\zh{得} = \cle{Degré} (Verbe+)};
\node[draw=poppink, thick, rounded corners, fill=poppinkL, minimum width=3.5cm, minimum height=1cm, below=of mnemo] (d3) {\zh{地} = \cle{Devant Verbe} (Manière)};
\draw[->, thick, popblue] (bonus) -- (mnemo);
\draw[->, thick, poporange] (mnemo) -- (d1);
\draw[->, thick, poporange] (mnemo) -- (d2);
\draw[->, thick, poporange] (mnemo) -- (d3);
\end{tikzpicture}
\legende{Carte mentale de révision : Les 3 « de » et l'entrée « Bonus Examen »}
\end{popfigure}

\begin{aretenir}[Check-list finale pour l'épreuve écrite (Module 5)]
\begin{itemize}
    \item \textbf{Repérage registre :} Panneau / Titre / Règlement $\rightarrow$ \zh{可} (au lieu de \zh{可以}).
    \item \textbf{Dictée / Complétion :} Appliquer la méthode \textbf{D-D-D} (Déterminant / Degré / Devant Verbe) \textbf{systématiquement}.
    \item \textbf{Traduction Thème (Fr $\rightarrow$ Ch) :} « Il court vite » $\nrightarrow$ \zh{他跑快} (phrase incomplète) $\rightarrow$ \zh{他跑得快}. « Ma voiture rouge » $\rightarrow$ \zh{我的红色的车} (ou \zh{我的红车}).
    \item \textbf{Expression écrite :} Utiliser \zh{可以} pour vos opinions/droits (\zh{我认为我们可以...}), \zh{可} pour citations de lois/règlements dans un essai argumentatif.
\end{itemize}
\end{aretenir}

\newpage

\coursec{Activité d'intégration}

\coursec{Leçon 39 — Le verbe modal 可以 et les particules 的、得、地}

\courssub{Objectifs de l'activité}

À l'issue de cette activité d'intégration, tu seras capable de :
\begin{itemize}
\item demander la permission et exprimer la possibilité avec le verbe modal \zh{可以} (\emph{kěyǐ}) dans une situation réelle d'achat ;
\item identifier et décrire un objet grâce à la particule \zh{的} (\emph{de}) dans une structure déterminative ;
\item évaluer une qualité ou un résultat avec la structure \zh{得} (\emph{de}) après un verbe ;
\item décrire la manière d'accomplir une action avec la structure \zh{地} (\emph{de}) avant un verbe ;
\item rédiger puis jouer un dialogue de huit répliques entre un client et un vendeur au marché ;
\item t'auto-évaluer à l'aide d'une grille portant sur l'usage correct des quatre structures.
\end{itemize}

\courssub{Situation}

Tu es à Douala, au marché central, un samedi matin. Tu veux acheter des mangues pour ta famille. Le vendeur, Monsieur Li, est un commerçant chinois installé au Cameroun depuis plusieurs années ; il parle bien chinois et un peu français. Tu veux profiter de l'occasion pour pratiquer ton chinois : demander la permission de goûter, décrire le fruit que tu cherches, commenter la qualité du service, et écouter le vendeur expliquer comment il prépare les fruits.

Voici le début du dialogue que tu vas compléter et jouer :

\begin{textebox}[Au marché de Douala — extrait]
客户：你好！请问，我可以尝一尝这个芒果吗？\\
\emph{Kèhù : Nǐ hǎo! Qǐngwèn, wǒ kěyǐ cháng yi cháng zhège mángguǒ ma?}\\
« Client : Bonjour ! Puis-je goûter cette mangue, s'il vous plaît ? »\\
商贩：当然可以！这是今天刚到的芒果。\\
\emph{Shāngfàn : Dāngrán kěyǐ! Zhè shì jīntiān gāng dào de mángguǒ.}\\
« Vendeur : Bien sûr ! Ce sont des mangues qui viennent d'arriver aujourd'hui. »
\end{textebox}

\courssub{Consignes}

\textbf{Tâche 1 — Compréhension et observation (10 min).} Lis l'extrait du dialogue à voix haute avec ton binôme. Souligne en bleu la structure avec \zh{可以}, puis réponds par écrit : qui demande la permission ? qui la donne ? Quelle expression le vendeur utilise-t-il pour répondre positivement ?

\textbf{Tâche 2 — Lexique (10 min).} À l'aide du tableau de vocabulaire de la leçon, choisis : deux fruits vendus au marché, deux adjectifs de qualité (par exemple \zh{甜} « sucré », \zh{新鲜} « frais »), deux adverbes de manière (par exemple \zh{快} « vite », \zh{小心} « avec précaution »). Vérifie le pinyin et les tons de chaque mot.

\textbf{Tâche 3 — Rédaction du dialogue (25 min).} En binôme, rédigez un dialogue complet de \textbf{huit répliques} (quatre par personnage) qui respecte le plan suivant :
\begin{enumerate}
\item Réplique 1 (client) : salutation et demande de permission avec \zh{可以} + verbe.
\item Réplique 2 (vendeur) : réponse positive ou négative avec \zh{可以} / \zh{不可以}.
\item Réplique 3 (client) : description du fruit souhaité avec \zh{的} (structure : adjectif ou verbe + \zh{的} + nom).
\item Réplique 4 (vendeur) : présentation du produit avec \zh{的}.
\item Réplique 5 (client) : commentaire sur la qualité du service avec verbe + \zh{得} + adjectif.
\item Réplique 6 (vendeur) : réponse avec verbe + \zh{得} + adjectif.
\item Réplique 7 (vendeur) : description de la manière de couper ou préparer le fruit avec adjectif + \zh{地} + verbe.
\item Réplique 8 (client) : remerciement et au revoir.
\end{enumerate}

\textbf{Tâche 4 — Relecture grammaticale (10 min).} Échangez votre dialogue avec un autre binôme. Vérifiez : la place de \zh{可以} (avant le verbe principal), la place de \zh{得} (après le verbe), la place de \zh{地} (avant le verbe), la présence de \zh{的} dans les groupes déterminatifs. Corrigez au crayon.

\textbf{Tâche 5 — Jeu de rôle (15 min).} Chaque binôme joue la scène devant la classe. Pendant les prestations des autres, remplis la grille d'auto-évaluation ci-dessous.

\begin{center}
\begin{tabularx}{\linewidth}{|X|X|X|}\hline
\rowcolor{popblueL} Critère & Oui & Non \\ \hline
\zh{可以} est placé avant le verbe principal & & \\ \hline
\zh{的} relie bien le déterminant au nom & & \\ \hline
\zh{得} est placé après le verbe, avant l'adjectif & & \\ \hline
\zh{地} est placé après l'adjectif, avant le verbe & & \\ \hline
Les tons sont correctement prononcés & & \\ \hline
Le dialogue comporte huit répliques & & \\ \hline
\end{tabularx}
\end{center}

\courssub{Ressources à mobiliser}

\begin{aretenir}[Les quatre structures de la leçon]
\begin{itemize}
\item \textbf{Permission / possibilité} : Sujet + \zh{可以} + Verbe (+ complément). Négation : \zh{不可以}. Question : \zh{可以……吗？}
\item \textbf{Détermination} : Déterminant (adjectif, verbe, nom, pronom) + \zh{的} + Nom. Exemple : \zh{今天刚到的芒果} « les mangues arrivées aujourd'hui ».
\item \textbf{Évaluation du résultat} : Verbe + \zh{得} + Adjectif. Exemple : \zh{说得很好} « parler très bien ».
\item \textbf{Manière} : Adjectif + \zh{地} + Verbe. Exemple : \zh{小心地切} « couper avec précaution ».
\end{itemize}
\end{aretenir}

\begin{propriete}[Bien distinguer 的、得、地 : trois particules, un même son \emph{de}]
\begin{center}
\begin{tabularx}{\linewidth}{|X|X|X|}\hline
\rowcolor{popblueL} Structure & Exemple en caractères + pinyin & Traduction \\ \hline
Déterminant + \zh{的} + Nom & \zh{新鲜的芒果} \emph{xīnxiān de mángguǒ} & la mangue fraîche \\ \hline
Verbe + \zh{得} + Adjectif & \zh{他说得很快。} \emph{Tā shuō de hěn kuài.} & Il parle (très) vite. \\ \hline
Adjectif + \zh{地} + Verbe & \zh{他很快地切。} \emph{Tā hěn kuài de qiē.} & Il coupe rapidement. \\ \hline
\end{tabularx}
\end{center}
Moyen mnémotechnique : \zh{的} s'attache à un \textbf{nom} (avant lui), \zh{得} juge un \textbf{résultat} (après le verbe), \zh{地} décrit une \textbf{manière} (avant le verbe).
\end{propriete}

\begin{center}
\begin{tabularx}{\linewidth}{|X|X|X|X|}\hline
\rowcolor{popblueL} Caractères & Pinyin & Nature & Traduction \\ \hline
尝 & cháng & verbe & goûter \\ \hline
芒果 & mángguǒ & nom & mangue \\ \hline
新鲜 & xīnxiān & adjectif & frais \\ \hline
甜 & tián & adjectif & sucré \\ \hline
切 & qiē & verbe & couper \\ \hline
小心 & xiǎoxīn & adjectif & prudent, avec précaution \\ \hline
服务 & fúwù & nom / verbe & service ; servir \\ \hline
便宜 & piányi & adjectif & bon marché \\ \hline
\end{tabularx}
\end{center}

\begin{attention}[Pièges fréquents des francophones]
\begin{itemize}
\item Ne dis pas \zh{我可以尝吗这个芒果} : le complément d'objet se place après le verbe, et la particule interrogative \zh{吗} se met en fin de phrase : \zh{我可以尝一尝这个芒果吗？}
\item \zh{得} et \zh{地} se prononcent tous les deux \emph{de}, mais \zh{得} suit le verbe (évaluation) et \zh{地} précède le verbe (manière). Ne les intervertis jamais à l'écrit : à l'examen, la confusion entre les trois caractères est sanctionnée.
\item Dans \zh{他切得很快} (\emph{Tā qiē de hěn kuài.}), l'adjectif vient après \zh{得} ; dans \zh{他很快地切} (\emph{Tā hěn kuài de qiē.}), il vient avant \zh{地}. Compare : « Il coupe vite » (résultat constaté) et « Il coupe rapidement » (manière d'agir).
\item \zh{的} ne se traduit pas toujours : \zh{红的芒果} (\emph{hóng de mángguǒ}) = « la mangue rouge », sans mot supplémentaire en français.
\item N'ajoute pas \zh{很} après \zh{得} quand l'adjectif est déjà évalué par le contexte ; en revanche, sans \zh{得}, l'adjectif seul exige souvent \zh{很} : \zh{芒果很甜。} « La mangue est sucrée. »
\end{itemize}
\end{attention}

\courssub{Proposition de production et corrigé commenté}

\begin{exoresolu}[Dialogue modèle : Au marché de Douala]
Rédige un dialogue de huit répliques entre un client et un vendeur, en mobilisant \zh{可以}, \zh{的}, \zh{得} et \zh{地}.
\tcblower
\textbf{Production modèle :}

客户：你好！请问，我可以尝一尝这个芒果吗？\\
\emph{Nǐ hǎo! Qǐngwèn, wǒ kěyǐ cháng yi cháng zhège mángguǒ ma?}\\
« Bonjour ! Puis-je goûter cette mangue, s'il vous plaît ? »

商贩：当然可以！这是今天刚到的芒果，很新鲜。\\
\emph{Dāngrán kěyǐ! Zhè shì jīntiān gāng dào de mángguǒ, hěn xīnxiān.}\\
« Bien sûr ! Ce sont des mangues arrivées aujourd'hui, très fraîches. »

客户：我想买甜的、不太贵的芒果。\\
\emph{Wǒ xiǎng mǎi tián de, bú tài guì de mángguǒ.}\\
« Je voudrais acheter des mangues sucrées et pas trop chères. »

商贩：您看，这些黄的芒果又甜又便宜。\\
\emph{Nín kàn, zhèxiē huáng de mángguǒ yòu tián yòu piányi.}\\
« Regardez, ces mangues jaunes sont à la fois sucrées et bon marché. »

客户：你的中文说得真好！\\
\emph{Nǐ de Zhōngwén shuō de zhēn hǎo!}\\
« Tu parles vraiment bien chinois ! »

商贩：谢谢！我在杜阿拉住了很多年，法语也说得不错。\\
\emph{Xièxie! Wǒ zài Dù'ālā zhù le hěn duō nián, Fǎyǔ yě shuō de búcuò.}\\
« Merci ! J'habite à Douala depuis de nombreuses années, je parle aussi pas mal français. »

商贩：您看，我小心地切，芒果切得又快又好。\\
\emph{Nín kàn, wǒ xiǎoxīn de qiē, mángguǒ qiē de yòu kuài yòu hǎo.}\\
« Regardez, je coupe avec précaution, et la mangue est coupée vite et bien. »

客户：太好了！我买两公斤。谢谢，再见！\\
\emph{Tài hǎo le! Wǒ mǎi liǎng gōngjīn. Xièxie, zàijiàn!}\\
« Parfait ! J'en prends deux kilos. Merci, au revoir ! »

\textbf{Commentaire des choix de langue :}
\begin{itemize}
\item Réplique 1 : \zh{可以} précède le verbe \zh{尝} ; la redondance \zh{尝一尝} adoucit la demande, comme « goûter un peu ».
\item Répliques 2, 3 et 4 : \zh{的} construit les groupes déterminatifs (\zh{今天刚到的芒果}, \zh{甜的芒果}, \zh{黄的芒果}) ; on omet le nom après \zh{的} quand le contexte est clair (\zh{甜的} = « les sucrées »).
\item Répliques 5 et 6 : \zh{说得很好} et \zh{说得不错} évaluent le résultat de l'action de parler ; attention, \zh{得} suit toujours le verbe.
\item Réplique 7 : \zh{小心地切} décrit la manière (adjectif + \zh{地} + verbe), puis \zh{切得又快又好} évalue le résultat : les deux structures coexistent dans la même phrase, ce qui montre leur complémentarité.
\item Formules de politesse : \zh{请问} pour introduire une question, \zh{您} (forme respectueuse de \zh{你}) employé par le vendeur envers le client.
\end{itemize}
\end{exoresolu}

\courssub{Exercice d'application et de transformation}

\begin{exercice}[À toi de jouer]
\begin{enumerate}
\item Complète avec \zh{的}, \zh{得} ou \zh{地} :\\
a) 他汉语说\_\_\_很好。\\
b) 她慢慢\_\_\_走。\\
c) 这是我喜欢\_\_\_水果。\\
d) 老师高兴\_\_\_说：« 你们考得不错！ »
\item Mets \zh{可以} à la bonne place, puis traduis : 我\_\_\_用\_\_\_你的词典\_\_\_吗？
\item Transforme en phrase d'évaluation avec \zh{得} : \zh{他切芒果很快。}
\end{enumerate}
\end{exercice}

\begin{corrige}[Exercice d'application]
\begin{enumerate}
\item a) \zh{得} : \zh{他汉语说得很好。} (\emph{Tā Hànyǔ shuō de hěn hǎo.}) « Il parle très bien chinois. » — évaluation après le verbe.\\
b) \zh{地} : \zh{她慢慢地走。} (\emph{Tā mànmàn de zǒu.}) « Elle marche lentement. » — manière avant le verbe.\\
c) \zh{的} : \zh{这是我喜欢的水果。} (\emph{Zhè shì wǒ xǐhuan de shuǐguǒ.}) « C'est le fruit que j'aime. » — déterminant avant le nom.\\
d) \zh{地} puis \zh{得} : \zh{老师高兴地说：« 你们考得不错！ »} (\emph{Lǎoshī gāoxìng de shuō : « Nǐmen kǎo de búcuò! »}) « Le professeur dit joyeusement : "Vous avez bien réussi l'examen !" »
\item \zh{我可以用你的词典吗？} (\emph{Wǒ kěyǐ yòng nǐ de cídiǎn ma?}) « Puis-je utiliser ton dictionnaire ? » — \zh{可以} se place avant le verbe principal \zh{用}, jamais entre le verbe et son complément.
\item \zh{他芒果切得很快。} ou \zh{他切芒果切得很快。} (\emph{Tā qiē mángguǒ qiē de hěn kuài.}) « Il coupe les mangues très vite. » — quand le verbe a un complément d'objet, on répète le verbe avant \zh{得} ou on avance l'objet en tête de phrase.
\end{enumerate}
\end{corrige}

\courssub{Bilan de l'activité}

Cette activité t'a permis de mobiliser dans une même situation de communication les quatre structures de la leçon : \zh{可以} pour demander et accorder la permission, \zh{的} pour identifier et décrire les objets, \zh{得} pour évaluer une qualité ou un résultat, et \zh{地} pour décrire la manière d'agir. Tu as également pratiqué l'expression orale en jouant la scène et développé ton esprit critique grâce à la grille d'auto-évaluation. Pour aller plus loin, réécris le dialogue en changeant de contexte (au restaurant, à la librairie, à la gare) en conservant les quatre structures : c'est un excellent entraînement pour l'épreuve écrite du baccalauréat.

\newpage

\coursec{Méthodes et exercices résolus}

\courssub{Observer : trois « de » et un « pouvoir » dans un dialogue}

\begin{textebox}[Au cybercafé de Douala]
\zh{甲：请问，我可以用这台电脑吗？}\\
\emph{Jiǎ : Qǐngwèn, wǒ kěyǐ yòng zhè tái diànnǎo ma?}\\
« A : Excusez-moi, est-ce que je peux utiliser cet ordinateur ? »\\
\zh{乙：可以，但是你不可以在这里抽烟。}\\
\emph{Yǐ : Kěyǐ, dànshì nǐ bù kěyǐ zài zhèlǐ chōuyān.}\\
« B : Oui, mais tu ne peux pas fumer ici. »\\
\zh{甲：好的。你的中文说得很好！}\\
\emph{Jiǎ : Hǎo de. Nǐ de Zhōngwén shuō de hěn hǎo!}\\
« A : D'accord. Tu parles très bien le chinois ! »\\
\zh{乙：谢谢！我每天都认真地学习。}\\
\emph{Yǐ : Xièxie! Wǒ měitiān dōu rènzhēn de xuéxí.}\\
« B : Merci ! J'étudie sérieusement tous les jours. »
\end{textebox}

\begin{aretenir}[Ce qu'il faut observer]
Dans ce dialogue, repère trois mots qui se prononcent tous « de » mais s'écrivent différemment : \zh{你的中文} (\zh{的} devant un nom), \zh{说得很好} (\zh{得} après un verbe), \zh{认真地学习} (\zh{地} devant un verbe). Repère aussi le verbe modal \zh{可以} (kěyǐ, « pouvoir, avoir la permission »), toujours placé entre le sujet et le verbe : \zh{我可以用}, \zh{你不可以抽烟}.
\end{aretenir}

\courssub{Construire une phrase avec le verbe modal 可以}

\begin{definition}[Le verbe modal 可以]
\zh{可以} (kěyǐ) est un \textbf{verbe modal} (verbe auxiliaire) qui exprime la \textbf{permission} (« avoir le droit de ») ou la \textbf{possibilité} (« il est possible de »). Comme tous les verbes modaux chinois (\zh{会}, \zh{能}, \zh{想}, \zh{要}), il se place \textbf{avant le verbe principal} et ne se conjugue jamais.
\end{definition}

\begin{propriete}[Schéma de la structure]
\begin{center}
\textbf{Sujet + 可以 + Verbe + (Complément)}
\end{center}
\begin{itemize}
\item Affirmation : \zh{我可以用这台电脑。} \emph{Wǒ kěyǐ yòng zhè tái diànnǎo.} « Je peux utiliser cet ordinateur. »
\item Négation : Sujet + \zh{不可以} (ou \zh{不能}) + Verbe : \zh{你不可以在这里抽烟。} \emph{Nǐ bù kěyǐ zài zhèlǐ chōuyān.} « Tu ne peux pas fumer ici. »
\item Question avec \zh{吗} : \zh{我可以用这台电脑吗？} « Est-ce que je peux utiliser cet ordinateur ? »
\item Question affirmative-négative : \zh{可以不可以} + Verbe : \zh{你可以不可以用这台电脑？} « Peux-tu utiliser cet ordinateur, oui ou non ? »
\item Réponse courte : \zh{可以。} / \zh{不可以。} (on ne répond jamais par le seul verbe principal).
\end{itemize}
\end{propriete}

\begin{methode}[Autorisation et possibilité avec 可以]
Pour exprimer la permission ou la possibilité avec \zh{可以} (kěyǐ) :
\begin{enumerate}
\item Identifier le sujet (la personne qui a la permission).
\item Placer \zh{可以} immédiatement \textbf{après le sujet} et \textbf{avant le verbe principal} : Sujet + 可以 + Verbe + (Complément).
\item Pour la négation, utiliser \zh{不可以} ou \zh{不能} : Sujet + 不可以 + Verbe.
\item Pour la question, deux formes : Sujet + 可以 + Verbe + 吗？ ou bien la forme affirmative-négative 可以不可以 + Verbe？
\item Réponse courte : 可以。/ 不可以。
\end{enumerate}
Réflexe : \zh{可以} ne se conjugue pas, ne prend jamais \zh{了}, et reste toujours à la même place, comme « pouvoir » devant l'infinitif en français.
\end{methode}

\begin{exemplebox}[可以 en contexte camerounais]
\zh{明天我们可以去市场。} \emph{Míngtiān wǒmen kěyǐ qù shìchǎng.} « Demain, nous pouvons aller au marché. »\\
\zh{学生不可以在教室里打电话。} \emph{Xuésheng bù kěyǐ zài jiàoshì lǐ dǎ diànhuà.} « Les élèves ne peuvent pas téléphoner dans la salle de classe. »\\
\zh{你可以说英语，也可以说法语。} \emph{Nǐ kěyǐ shuō Yīngyǔ, yě kěyǐ shuō Fǎyǔ.} « Tu peux parler anglais, et tu peux aussi parler français. »
\end{exemplebox}

\begin{exoresolu}[Au cybercafé de Douala]
\textbf{Énoncé.} Traduis en chinois : 1) « Est-ce que je peux utiliser cet ordinateur ? » 2) « Tu ne peux pas fumer ici. » 3) « Demain, nous pouvons aller au marché. »
\tcblower
\textbf{Solution détaillée.}
\begin{enumerate}
\item Sujet 我 + 可以 + verbe 用 + complément 这台电脑 + 吗.\\
\zh{我可以用这台电脑吗？} \emph{Wǒ kěyǐ yòng zhè tái diànnǎo ma?} « Est-ce que je peux utiliser cet ordinateur ? »\\
Justification : 可以 se place avant le verbe 用 ; le classificateur d'un ordinateur est 台.
\item Négation : 不可以 (ou 不能) + complément de lieu 在这里 + verbe 抽烟.\\
\zh{你不可以在这里抽烟。} \emph{Nǐ bù kěyǐ zài zhèlǐ chōuyān.} « Tu ne peux pas fumer ici. »\\
Justification : le complément de lieu introduit par 在 se place entre 可以 et le verbe, jamais en fin de phrase.
\item 明天 (temps, en tête) + 我们 + 可以 + 去 + 市场.\\
\zh{明天我们可以去市场。} \emph{Míngtiān wǒmen kěyǐ qù shìchǎng.} « Demain, nous pouvons aller au marché. »\\
Justification : le complément de temps se place en tête de phrase ou après le sujet, mais jamais après le verbe.
\end{enumerate}
Conclusion : la structure Sujet + 可以 + Verbe reste invariable ; seuls les compléments de temps et de lieu obéissent à des règles de position précises.
\end{exoresolu}

\courssub{Employer 的 pour la détermination}

\begin{definition}[Le mot-outil 的]
\zh{的} (de) est une \textbf{particule structurale} qui relie un élément \textbf{déterminant} (possesseur, adjectif, proposition) au \textbf{nom} qu'il détermine. Schéma : \textbf{Déterminant + 的 + Nom}. Le nom principal vient toujours \textbf{après} \zh{的} : c'est l'inverse du français (« le livre \emph{de} mon ami » → \zh{我朋友的书}).
\end{definition}

\begin{methode}[Le mot-outil 的]
\zh{的} (de) relie un élément déterminant au nom qu'il détermine :
\begin{enumerate}
\item Schéma : Déterminant + 的 + Nom déterminé (possession, qualité, appartenance).
\item Possession : 我的书 (wǒ de shū) « mon livre » ; 喀麦隆的首都 (Kāmàilóng de shǒudū) « la capitale du Cameroun ».
\item Adjectif monosyllabique seul : 的 est souvent omis (好朋友, 新书) ; adjectif dissyllabique ou développé : 的 obligatoire (漂亮的衣服).
\item Après un pronom personnel + membre de la famille ou relation proche, 的 peut tomber : 我妈妈 « ma mère ».
\item Attention : l'ordre est l'inverse du français : ce qui suit 的 est le nom principal.
\end{enumerate}
\end{methode}

\begin{exoresolu}[La famille et l'école]
\textbf{Énoncé.} Traduis en chinois : 1) « le téléphone de mon père » 2) « une très belle ville » 3) « les élèves de notre école ».
\tcblower
\textbf{Solution détaillée.}
\begin{enumerate}
\item Déterminant 我爸爸 + 的 + nom 手机 : \zh{我爸爸的手机} \emph{wǒ bàba de shǒujī} « le téléphone de mon père ». (我爸爸 : 的 omis entre pronom et parent proche ; 的 obligatoire devant 手机.)
\item Adjectif développé 非常漂亮 + 的 + 城市 : \zh{非常漂亮的城市} \emph{fēicháng piàoliang de chéngshì} « une très belle ville ». L'adverbe 非常 rend 的 obligatoire.
\item 我们学校 + 的 + 学生 : \zh{我们学校的学生} \emph{wǒmen xuéxiào de xuésheng} « les élèves de notre école ».
\end{enumerate}
Conclusion : dans tous les cas, le nom principal vient \textbf{après} 的, contrairement au français où il précède « de ».
\end{exoresolu}

\courssub{Employer 得 pour le complément de degré}

\begin{definition}[Le mot-outil 得]
\zh{得} (de) est une particule structurale qui introduit le \textbf{complément de degré ou de manière évaluative}, placé \textbf{après le verbe}. Il sert à évaluer comment une action est réalisée : \zh{他说得很好} « il parle bien ». Schéma : \textbf{Verbe + 得 + Adjectif}.
\end{definition}

\begin{methode}[Le mot-outil 得]
\zh{得} (de) introduit le complément de degré, placé \textbf{après le verbe} :
\begin{enumerate}
\item Schéma : Verbe + 得 + Adjectif (comment l'action est faite) : 他说得很好 « il parle bien ».
\item Si le verbe a un complément d'objet, on répète le verbe : 他说法语说得很好, ou on avance l'objet : 他法语说得很好.
\item Négation : Verbe + 得 + 不 + Adjectif : 他跑得很快 / 跑得不快.
\item Question : Verbe + 得 + 怎么样？ « comment... ? » : 你汉语说得怎么样？
\item Ne jamais mettre le complément de manière avant le verbe comme en français (« il bien parle » est impossible).
\end{enumerate}
\end{methode}

\begin{exoresolu}[Au terrain de football]
\textbf{Énoncé.} Traduis en chinois : 1) « Il joue très bien au football. » 2) « Elle ne chante pas bien. » 3) « Comment écris-tu les caractères ? »
\tcblower
\textbf{Solution détaillée.}
\begin{enumerate}
\item Le verbe 踢 a un objet 足球 : on répète le verbe. \zh{他踢足球踢得非常好。} \emph{Tā tī zúqiú tī de fēicháng hǎo.} « Il joue très bien au football. » Variante admise : 他足球踢得非常好。
\item Négation : 唱 + 得 + 不 + 好 : \zh{她唱得不好。} \emph{Tā chàng de bù hǎo.} « Elle ne chante pas bien. »
\item Question avec 怎么样 : \zh{你写汉字写得怎么样？} \emph{Nǐ xiě Hànzì xiě de zěnmeyàng?} « Comment écris-tu les caractères ? » (répétition du verbe 写 car il y a un objet 汉字).
\end{enumerate}
Conclusion : 得 suit toujours le verbe ; la présence d'un objet oblige à répéter le verbe ou à avancer l'objet.
\end{exoresolu}

\courssub{Employer 地 pour le complément de manière}

\begin{definition}[Le mot-outil 地]
\zh{地} (de) est une particule structurale qui relie un \textbf{adverbe de manière} au \textbf{verbe}, placé \textbf{avant le verbe}. Schéma : \textbf{Adverbe de manière + 地 + Verbe} : \zh{他认真地学习} « il étudie sérieusement ». Dans cet emploi grammatical, \zh{地} se prononce « de » (et non « dì » comme dans \zh{土地} « la terre »).
\end{definition}

\begin{methode}[Le mot-outil 地]
\zh{地} (de) relie un adverbe de manière au verbe, placé \textbf{avant le verbe} :
\begin{enumerate}
\item Schéma : Adverbe de manière + 地 + Verbe : 他认真地学习 « il étudie sérieusement ».
\item L'adverbe est en général dissyllabique ou redoublé : 慢慢地走 « marcher lentement », 高高兴兴地回家 « rentrer joyeusement à la maison ».
\item Différence avec 得 : 地 décrit la manière \textbf{avant} l'action (comment on agit), 得 évalue le résultat \textbf{après} l'action (à quel point on le fait bien).
\item 地 se prononce « de » (et non « dì ») dans cet emploi grammatical.
\end{enumerate}
\end{methode}

\begin{exoresolu}[En classe de chinois]
\textbf{Énoncé.} Traduis en chinois : 1) « Les élèves écoutent attentivement le professeur. » 2) « Elle parle lentement. » 3) « Nous préparons sérieusement l'examen. »
\tcblower
\textbf{Solution détaillée.}
\begin{enumerate}
\item Manière 认真 + 地 + verbe 听 : \zh{学生们认真地听老师讲课。} \emph{Xuéshengmen rènzhēn de tīng lǎoshī jiǎngkè.} « Les élèves écoutent attentivement le professeur. »
\item \zh{她慢慢地说。} \emph{Tā mànmàn de shuō.} « Elle parle lentement. » (adverbe redoublé 慢慢地, 地 facultatif mais conseillé).
\item \zh{我们认真地准备考试。} \emph{Wǒmen rènzhēn de zhǔnbèi kǎoshì.} « Nous préparons sérieusement l'examen. »
\end{enumerate}
Conclusion : la manière exprimée avec 地 se place avant le verbe, comme l'adverbe français ; c'est la position qui distingue 地 de 得.
\end{exoresolu}

\courssub{Choisir entre 的、得、地 et transformer des phrases}

\begin{propriete}[Tableau récapitulatif des trois « de »]
\begin{center}
\begin{tabularx}{\linewidth}{|X|X|X|X|}
\hline
\rowcolor{popblueL} Particule & Position et fonction & Exemple (caractères + pinyin) & Traduction \\
\hline
的 (de) & devant un nom : détermination & 我的老师 wǒ de lǎoshī & mon professeur \\
\hline
地 (de) & devant un verbe : manière & 认真地学习 rènzhēn de xuéxí & étudier sérieusement \\
\hline
得 (de) & après un verbe : degré, évaluation & 说得很好 shuō de hěn hǎo & parler très bien \\
\hline
\end{tabularx}
\end{center}
\end{propriete}

\begin{methode}[Le réflexe des trois « de »]
Les trois mots-outils se prononcent « de » mais ont des fonctions différentes :
\begin{enumerate}
\item \textbf{的} : devant un \textbf{nom} (détermination) → 我的老师, 漂亮的城市.
\item \textbf{地} : devant un \textbf{verbe} (manière) → 认真地学习.
\item \textbf{得} : après un \textbf{verbe} (degré, évaluation) → 说得很好.
\item Test rapide : regarder le mot qui suit ou précède. Nom → 的 ; verbe après → 地 ; verbe avant → 得.
\item Transformation type examen : « Il parle bien le chinois » (évaluation → 得) : 他汉语说得很好 ; « Il parle prudemment » (manière → 地) : 他小心地说.
\end{enumerate}
\end{methode}

\begin{popfigure}
\begin{tikzpicture}[node distance=0.9cm]
\node[draw=popblue, fill=popblueL, rounded corners, align=center, minimum width=3.4cm] (de1) {\zh{的} de\\ \textbf{+ nom}\\ détermination};
\node[draw=popgreen, fill=popgreenL, rounded corners, align=center, minimum width=3.4cm, right=1.1cm of de1] (de2) {\zh{地} de\\ \textbf{+ verbe}\\ manière (avant)};
\node[draw=poporange, fill=poporangeL, rounded corners, align=center, minimum width=3.4cm, right=1.1cm of de2] (de3) {\zh{得} de\\ verbe \textbf{+}\\ degré (après)};
\node[below=0.35cm of de1, align=center] {\zh{我的书}\\ \emph{mon livre}};
\node[below=0.35cm of de2, align=center] {\zh{认真地听}\\ \emph{écouter attentivement}};
\node[below=0.35cm of de3, align=center] {\zh{说得很好}\\ \emph{parler très bien}};
\end{tikzpicture}
\legende{Les trois particules « de » : même prononciation, trois positions, trois fonctions.}
\end{popfigure}

\begin{exoresolu}[À compléter avec 的、得 ou 地]
\textbf{Énoncé.} Complète avec 的、得 ou 地 :\\
1) 这是我\_\_老师。 2) 他高兴\_\_回家了。 3) 她汉语说\_\_很好。 4) 雅温得是喀麦隆\_\_首都。 5) 弟弟认真\_\_做作业。
\tcblower
\textbf{Solution détaillée.}
\begin{enumerate}
\item 的 : devant le nom 老师 (détermination). \zh{这是我的老师。} \emph{Zhè shì wǒ de lǎoshī.} « C'est mon professeur. »
\item 地 : devant le verbe 回家 (manière). \zh{他高兴地回家了。} \emph{Tā gāoxìng de huí jiā le.} « Il est rentré joyeusement à la maison. »
\item 得 : après le verbe 说 (évaluation). \zh{她汉语说得很好。} \emph{Tā Hànyǔ shuō de hěn hǎo.} « Elle parle très bien le chinois. »
\item 的 : devant le nom 首都. \zh{雅温得是喀麦隆的首都。} \emph{Yǎwēndé shì Kāmàilóng de shǒudū.} « Yaoundé est la capitale du Cameroun. »
\item 地 : devant le verbe 做. \zh{弟弟认真地做作业。} \emph{Dìdi rènzhēn de zuò zuòyè.} « Le petit frère fait sérieusement ses devoirs. »
\end{enumerate}
Conclusion : identifier la nature du mot voisin (nom ou verbe) et la position (avant ou après le verbe) suffit à choisir le bon « de ».
\end{exoresolu}

\begin{attention}[Les 5 erreurs qui coûtent des points au Bac]
\begin{enumerate}
\item \textbf{Confondre les trois « de ».} Écrire 的 partout : *他说\textcolor{popink}{的}很好 est faux ; après un verbe, il faut 得 : 他说得很好. Mémorise : 的 + nom, 地 + verbe, verbe + 得.
\item \textbf{Placer 可以 après le verbe.} Calque du français « je veux pouvoir partir » : *我去可以 est impossible. Toujours Sujet + 可以 + Verbe : 我可以去.
\item \textbf{Mettre le complément de lieu en fin de phrase.} *我可以抽烟在这里 est faux ; le lieu avec 在 se place avant le verbe : 我不可以在这里抽烟.
\item \textbf{Oublier de répéter le verbe avec 得.} *他说汉语得很好 est fautif ; avec un objet, on répète le verbe : 他说汉语说得很好, ou on avance l'objet : 他汉语说得很好.
\item \textbf{Inverser l'ordre avec 的.} Calque du français « le livre de mon ami » : *书的我朋友 est faux ; le déterminant précède : 我朋友的书. Le nom principal vient toujours après 的.
\end{enumerate}
\end{attention}

\courssub{Compléments utiles à l'examen}

\begin{aretenir}[可以 ou 能 ?]
Les deux verbes modaux se traduisent par « pouvoir », mais :
\begin{itemize}
\item \zh{可以} insiste sur la \textbf{permission} (le règlement, l'autorisation) : \zh{这里不可以抽烟。} « Il est interdit de fumer ici. »
\item \zh{能} (néng) insiste sur la \textbf{capacité ou la possibilité matérielle} : \zh{他今天不能来，他很忙。} « Il ne peut pas venir aujourd'hui, il est très occupé. »
\item À l'examen, pour une interdiction, préfère \zh{不可以} ou \zh{不能} ; les deux sont acceptés dans la négation.
\end{itemize}
\end{aretenir}

\begin{savaistu}[Trois caractères, une seule prononciation]
En chinois parlé, \zh{的}, \zh{得} et \zh{地} se prononcent tous « de » (ton neutre) : impossible de les distinguer à l'oreille ! C'est l'écriture qui porte la grammaire. À l'écrit, les Chinois eux-mêmes font parfois la confusion, comme les francophones hésitent entre « a » et « à », « et » et « est ». Au Bac, la dictée et la traduction sanctionnent le choix du bon caractère : entraîne-toi à les écrire sans faute.
\end{savaistu}

\begin{exercice}[Entraînement type Bac]
\begin{enumerate}
\item Traduis en chinois : « Les élèves ne peuvent pas utiliser le téléphone en classe. »
\item Traduis en chinois : « Ma sœur chante très bien. »
\item Traduis en français : \zh{他高兴地告诉我们这个好消息。}
\item Complète avec 的、得 ou 地 : 她慢慢\_\_走。/ 哥哥跑\_\_很快。/ 这是中国\_\_地图。
\item Transforme avec 得 : \zh{他汉语说得很好。} → forme négative.
\end{enumerate}
\end{exercice}

\begin{corrige}[Entraînement type Bac]
\begin{enumerate}
\item \zh{学生们不可以在教室里用手机。} \emph{Xuéshengmen bù kěyǐ zài jiàoshì lǐ yòng shǒujī.} (可以 avant le verbe ; lieu 在教室里 avant 用.)
\item \zh{我妹妹唱得很好。} \emph{Wǒ mèimei chàng de hěn hǎo.} (évaluation après le verbe → 得 ; pas d'objet, donc pas de répétition.)
\item « Il nous a annoncé joyeusement cette bonne nouvelle. » (高兴地 + verbe 告诉 : manière avant l'action → 地.)
\item 地 : \zh{她慢慢地走。} ; 得 : \zh{哥哥跑得很快。} ; 的 : \zh{这是中国的地图。}
\item \zh{他汉语说得不好。} \emph{Tā Hànyǔ shuō de bù hǎo.} « Il ne parle pas bien le chinois. » (négation : 得 + 不 + adjectif.)
\end{enumerate}
\end{corrige}

\newpage

\coursec{Exercices}

\courssub{Je vérifie mes connaissances}

\begin{exercice}[QCM : le bon choix]
Pour chaque phrase, choisis la bonne réponse (A, B ou C).

\begin{enumerate}
\item 学生不\underline{\hspace{1.5cm}}在教室里打电话。\\
A. 可以 \quad B. 得 \quad C. 地
\item 他汉语说\underline{\hspace{1.5cm}}很好。\\
A. 的 \quad B. 得 \quad C. 地
\item 她高兴\underline{\hspace{1.5cm}}唱起了歌。\\
A. 的 \quad B. 得 \quad C. 地
\item 这是我\underline{\hspace{1.5cm}}新手机。\\
A. 的 \quad B. 得 \quad C. 地
\end{enumerate}
\end{exercice}

\begin{exercice}[Vrai ou faux ? Justifie ta réponse]
Dis si chaque phrase est correcte (V) ou incorrecte (F), puis justifie en français.

\begin{enumerate}
\item \zh{我可以明天不来上课吗？} \emph{(Wǒ kěyǐ míngtiān bù lái shàngkè ma ?)}
\item \zh{他跑得很快。} \emph{(Tā pǎo de hěn kuài.)}
\item \zh{我地朋友住在杜阿拉。} \emph{(Wǒ de péngyou zhù zài Dù'ālā.)}
\item \zh{不可以在这里吸烟。} \emph{(Bù kěyǐ zài zhèlǐ xīyān.)}
\end{enumerate}
\end{exercice}

\begin{exercice}[Vocabulaire : complète le tableau]
Complète le tableau en donnant le pinyin et la traduction française de chaque mot.

\begin{center}
\begin{tabularx}{\linewidth}{|X|X|X|X|}
\hline
\rowcolor{popblueL} Caractères & Pinyin & Nature & Traduction \\
\hline
可以 & \ldots\ldots\ldots & verbe modal & \ldots\ldots\ldots \\
\hline
允许 & \ldots\ldots\ldots & \ldots\ldots\ldots & \ldots\ldots\ldots \\
\hline
禁止 & \ldots\ldots\ldots & \ldots\ldots\ldots & \ldots\ldots\ldots \\
\hline
慢 & \ldots\ldots\ldots & \ldots\ldots\ldots & \ldots\ldots\ldots \\
\hline
认真地 & \ldots\ldots\ldots & \ldots\ldots\ldots & \ldots\ldots\ldots \\
\hline
\end{tabularx}
\end{center}
\end{exercice}

\begin{exercice}[Associe le pinyin aux caractères]
Relie chaque groupe de caractères à son pinyin.

\begin{center}
\begin{tabularx}{\linewidth}{|X|X|}
\hline
\rowcolor{popblueL} Caractères & Pinyin \\
\hline
1. 说得好 & a. rènzhēn de xuéxí \\
\hline
2. 认真地学习 & b. kěyǐ jìnlái ma \\
\hline
3. 可以进来吗 & c. shuō de hǎo \\
\hline
4. 我的老师 & d. wǒ de lǎoshī \\
\hline
\end{tabularx}
\end{center}
\end{exercice}

\courssub{Je m'entraîne}

\begin{exercice}[Traduis en français]
Traduis les phrases suivantes en français.

\begin{enumerate}
\item \zh{在图书馆里不可以大声说话。} \emph{(Zài túshūguǎn li bù kěyǐ dàshēng shuōhuà.)}
\item \zh{她的汉字写得非常漂亮。} \emph{(Tā de Hànzì xiě de fēicháng piàoliang.)}
\item \zh{孩子们开心地玩着手机。} \emph{(Háizi men kāixīn de wánzhe shǒujī.)}
\item \zh{老师，我可以借您的词典吗？} \emph{(Lǎoshī, wǒ kěyǐ jiè nín de cídiǎn ma ?)}
\end{enumerate}
\end{exercice}

\begin{exercice}[Transforme les phrases]
Transforme chaque phrase affirmative : (a) en phrase négative, (b) en question avec \zh{吗}.

Exemple : \zh{他可以上网。} → \zh{他不可以上网。} / \zh{他可以上网吗？}

\begin{enumerate}
\item \zh{我们可以看电视。} \emph{(Wǒmen kěyǐ kàn diànshì.)}
\item \zh{你可以用我的电脑。} \emph{(Nǐ kěyǐ yòng wǒ de diànnǎo.)}
\item \zh{学生可以带手机去学校。} \emph{(Xuésheng kěyǐ dài shǒujī qù xuéxiào.)}
\item \zh{我们可以在这里拍照。} \emph{(Wǒmen kěyǐ zài zhèlǐ pāizhào.)}
\end{enumerate}
\end{exercice}

\begin{exercice}[Texte à trous : 的、得 ou 地]
Complète le texte suivant avec la bonne particule : \zh{的}, \zh{得} ou \zh{地}.

\begin{textebox}[À la médiathèque]
玛丽是我们班\underline{\hspace{1cm}}新同学。她汉语说\underline{\hspace{1cm}}很好，因为她每天认真\underline{\hspace{1cm}}学习。今天下午，她高兴\underline{\hspace{1cm}}对我说：« 图书馆\underline{\hspace{1cm}}电脑可以用，我们一起去做作业吧！» 她打字打\underline{\hspace{1cm}}真快！\\
\emph{(Mǎlì shì wǒmen bān \ldots\ xīn tóngxué. Tā Hànyǔ shuō \ldots\ hěn hǎo, yīnwèi tā měitiān rènzhēn \ldots\ xuéxí. Jīntiān xiàwǔ, tā gāoxìng \ldots\ duì wǒ shuō : « Túshūguǎn \ldots\ diànnǎo kěyǐ yòng, wǒmen yìqǐ qù zuò zuòyè ba ! » Tā dǎzì dǎ \ldots\ zhēn kuài !)}
\end{textebox}
\end{exercice}

\begin{exercice}[Remets les mots dans l'ordre]
Remets les mots dans le bon ordre pour former des phrases correctes, puis traduis-les en français.

\begin{enumerate}
\item \zh{可以 / 我 / 吗 / 出去 / ？}
\item \zh{得 / 他 / 跑 / 很 / 快}
\item \zh{地 / 认真 / 听 / 我们 / 老师 / 讲课}
\item \zh{的 / 这 / 是 / 我 / 手机 / 新}
\end{enumerate}
\end{exercice}

\courssub{Je me prépare au Bac}

\begin{exercice}[Compréhension de texte et questions de langue (10 points)]
Lis le texte suivant, puis réponds aux questions.

\begin{textebox}[Le téléphone portable au lycée]
在喀麦隆的很多中学里，学生上课的时候不可以使用手机。校长说：« 手机很有用，但是上课的时候学生应该认真地听老师讲课。» 下课以后，学生可以用手机给家长打电话，也可以上网查资料。有的学生觉得学校的规定太严格了，可是老师们认为这个规定很好。一个雅温得的学生说：« 我汉语学得不太好，所以我想下课以后用手机学习汉语。这样我可以学得更开心！»\\
\emph{(Zài Kāmàilóng de hěn duō zhōngxué li, xuésheng shàngkè de shíhou bù kěyǐ shǐyòng shǒujī. Xiàozhǎng shuō : « Shǒujī hěn yǒuyòng, dànshì shàngkè de shíhou xuésheng yīnggāi rènzhēn de tīng lǎoshī jiǎngkè. » Xiàkè yǐhòu, xuésheng kěyǐ yòng shǒujī gěi jiāzhǎng dǎ diànhuà, yě kěyǐ shàngwǎng chá zīliào. Yǒu de xuésheng juéde xuéxiào de guīdìng tài yángé le, kěshì lǎoshī men rènwéi zhège guīdìng hěn hǎo. Yí ge Yǎwēndé de xuésheng shuō : « Wǒ Hànyǔ xué de bù tài hǎo, suǒyǐ wǒ xiǎng xiàkè yǐhòu yòng shǒujī xuéxí Hànyǔ. Zhèyàng wǒ kěyǐ xué de gèng kāixīn ! »)}\\
« Dans beaucoup de collèges et lycées du Cameroun, les élèves ne peuvent pas utiliser leur téléphone pendant les cours. Le proviseur dit : "Le téléphone est très utile, mais pendant les cours, les élèves doivent écouter attentivement le professeur." Après les cours, les élèves peuvent utiliser leur téléphone pour appeler leurs parents, et peuvent aussi chercher des informations sur Internet. Certains élèves trouvent le règlement de l'école trop strict, mais les professeurs pensent que ce règlement est très bien. Un élève de Yaoundé dit : "Je n'apprends pas très bien le chinois, donc je voudrais utiliser mon téléphone pour étudier le chinois après les cours. Comme ça, je pourrai apprendre avec encore plus de joie !" »
\end{textebox}

\textbf{Questions de compréhension (réponds en français) :}
\begin{enumerate}
\item Que peuvent faire les élèves avec leur téléphone après les cours ? (2 points)
\item Pourquoi l'élève de Yaoundé veut-il utiliser son téléphone ? (2 points)
\end{enumerate}

\textbf{Questions de langue :}
\begin{enumerate}
\item Relève dans le texte deux phrases avec \zh{可以} ou \zh{不可以} et traduis-les. (2 points)
\item Relève une phrase avec \zh{得} et une avec \zh{地}, puis explique la fonction de chaque particule. (2 points)
\item Complète le dialogue avec \zh{可以}, \zh{不可以} ou \zh{可以吗} :\\
— 老师，我\underline{\hspace{1.5cm}}去喝水\underline{\hspace{1.5cm}}？\\
— 现在\underline{\hspace{1.5cm}}，下课以后再去。 (3 points)
\end{enumerate}
\end{exercice}

\begin{exercice}[Thème et version (10 points)]
\textbf{A. Thème : traduis en chinois (caractères obligatoires). (6 points)}
\begin{enumerate}
\item On ne peut pas fumer à l'hôpital.
\item Mon ami parle très bien le chinois.
\item Les élèves écoutent attentivement le professeur.
\end{enumerate}

\textbf{B. Version : traduis en français. (4 points)}

\begin{textebox}[Version]
妈妈高兴地对我说：« 你的房间打扫得真干净！现在你可以看电视了。»\\
\emph{(Māma gāoxìng de duì wǒ shuō : « Nǐ de fángjiān dǎsǎo de zhēn gānjìng ! Xiànzài nǐ kěyǐ kàn diànshì le. »)}
\end{textebox}
\end{exercice}

\begin{exercice}[Expression écrite (10 points)]
\textbf{Sujet :} Tu es élève en Terminale. Tu veux utiliser le Wi-Fi du lycée pour faire tes recherches. Rédige en chinois un message adressé au proviseur (\zh{校长}) pour demander la permission.

\textbf{Consignes :}
\begin{itemize}
\item Écris entre 60 et 80 caractères chinois.
\item Utilise au moins une fois \zh{可以} ou \zh{可以吗}.
\item Utilise au moins une particule parmi \zh{的}, \zh{得}, \zh{地}.
\item Termine par une formule de politesse, par exemple \zh{谢谢您！} \emph{(Xièxie nín !)}
\end{itemize}
\end{exercice}



\newpage

\coursec{Corrigés des exercices}

\courssub{Je vérifie mes connaissances}

\begin{corrige}[Exercice 1 — QCM : le bon choix]
\begin{enumerate}
\item \textbf{A. 可以} : \zh{学生不可以在教室里打电话。} \emph{(Xuésheng bù kěyǐ zài jiàoshì li dǎ diànhuà.)} « Les élèves ne peuvent pas téléphoner dans la salle de classe. » La négation \zh{不} se place devant le verbe modal \zh{可以}.
\item \textbf{B. 得} : \zh{他汉语说得很好。} \emph{(Tā Hànyǔ shuō de hěn hǎo.)} « Il parle très bien chinois. » \zh{得} introduit le complément de degré après le verbe \zh{说}.
\item \textbf{C. 地} : \zh{她高兴地唱起了歌。} \emph{(Tā gāoxìng de chàng qǐ le gē.)} « Joyeuse, elle s'est mise à chanter. » \zh{地} suit l'adverbe \zh{高兴} devant le verbe.
\item \textbf{A. 的} : \zh{这是我的新手机。} \emph{(Zhè shì wǒ de xīn shǒujī.)} « C'est mon nouveau téléphone portable. » \zh{的} marque la possession (\zh{我的} = « mon »).
\end{enumerate}
\end{corrige}

\begin{corrige}[Exercice 2 — Vrai ou faux ?]
\begin{enumerate}
\item \textbf{V} : « Puis-je ne pas venir en cours demain ? » Structure correcte : Sujet + \zh{可以} + complément de temps + verbe + \zh{吗}. La négation \zh{不} porte ici sur le verbe \zh{来}, ce qui est possible.
\item \textbf{V} : « Il court très vite. » Verbe + \zh{得} + complément de degré \zh{很快} : structure correcte.
\item \textbf{F} : il faut écrire \zh{我的朋友住在杜阿拉。} \emph{(Wǒ de péngyou zhù zài Dù'ālā.)} « Mon ami habite à Douala. » Devant un nom, on utilise la particule de possession \zh{的}, jamais \zh{地} (réservé aux adverbes devant un verbe).
\item \textbf{V} : « Il est interdit de fumer ici. » Phrase d'interdiction sans sujet : \zh{不可以} + complément de lieu + verbe. Structure correcte, typique des panneaux et règlements.
\end{enumerate}
\end{corrige}

\begin{corrige}[Exercice 3 — Vocabulaire : complète le tableau]
\begin{center}
\begin{tabularx}{\linewidth}{|X|X|X|X|}
\hline
\rowcolor{popblueL} Caractères & Pinyin & Nature & Traduction \\
\hline
可以 & kěyǐ & verbe modal & pouvoir (permission), il est permis de \\
\hline
允许 & yǔnxǔ & verbe & permettre, autoriser \\
\hline
禁止 & jìnzhǐ & verbe & interdire \\
\hline
慢 & màn & adjectif & lent \\
\hline
认真地 & rènzhēn de & adverbe & sérieusement, attentivement \\
\hline
\end{tabularx}
\end{center}
\end{corrige}

\begin{corrige}[Exercice 4 — Associe le pinyin aux caractères]
1 -- c (\zh{说得好} : « parler bien », verbe + \zh{得} + adjectif) ; 2 -- a (\zh{认真地学习} : « étudier sérieusement », adverbe + \zh{地} + verbe) ; 3 -- b (\zh{可以进来吗} : « puis-je entrer ? ») ; 4 -- d (\zh{我的老师} : « mon professeur », possession avec \zh{的}).
\end{corrige}

\courssub{Je m'entraîne}

\begin{corrige}[Exercice 5 — Traduis en français]
\begin{enumerate}
\item « À la bibliothèque, on ne peut pas parler fort. » (\zh{不可以} = interdiction ; \zh{大声} = « à voix haute ».)
\item « Elle écrit les caractères chinois très joliment. » (litt. « Ses caractères chinois sont écrits très beaux » : verbe \zh{写} + \zh{得} + complément de degré.)
\item « Les enfants jouent joyeusement avec leurs téléphones portables. » (\zh{开心地} : adverbe + \zh{地} devant le verbe \zh{玩}.)
\item « Monsieur le professeur, puis-je emprunter votre dictionnaire ? » (\zh{可以} + verbe + \zh{吗} : demande de permission polie ; \zh{您} = « vous » de politesse.)
\end{enumerate}
\end{corrige}

\begin{corrige}[Exercice 6 — Transforme les phrases]
Rappel de la règle : la négation de \zh{可以} est \zh{不可以} (\zh{不} devant le modal) ; la question se forme en ajoutant \zh{吗} à la fin, sans changer l'ordre des mots.
\begin{enumerate}
\item \zh{我们不可以看电视。} « Nous ne pouvons pas regarder la télévision. » / \zh{我们可以看电视吗？} « Pouvons-nous regarder la télévision ? »
\item \zh{你不可以用我的电脑。} « Tu ne peux pas utiliser mon ordinateur. » / \zh{你可以用我的电脑吗？} « Puis-je utiliser ton ordinateur ? » (litt. « Tu peux utiliser mon ordinateur ? »)
\item \zh{学生不可以带手机去学校。} « Les élèves ne peuvent pas apporter leur téléphone à l'école. » / \zh{学生可以带手机去学校吗？} « Les élèves peuvent-ils apporter leur téléphone à l'école ? »
\item \zh{我们不可以在这里拍照。} « Nous ne pouvons pas prendre de photos ici. » / \zh{我们可以在这里拍照吗？} « Pouvons-nous prendre des photos ici ? »
\end{enumerate}
\end{corrige}

\begin{corrige}[Exercice 7 — Texte à trous : 的、得 ou 地]
\zh{玛丽是我们班的新同学。她汉语说得很好，因为她每天认真地学习。今天下午，她高兴地对我说：« 图书馆的电脑可以用，我们一起去做作业吧！» 她打字打得真快！}\\
\emph{(Mǎlì shì wǒmen bān de xīn tóngxué. Tā Hànyǔ shuō de hěn hǎo, yīnwèi tā měitiān rènzhēn de xuéxí. Jīntiān xiàwǔ, tā gāoxìng de duì wǒ shuō : « Túshūguǎn de diànnǎo kěyǐ yòng, wǒmen yìqǐ qù zuò zuòyè ba ! » Tā dǎzì dǎ de zhēn kuài !)}\\
« Marie est la nouvelle camarade de notre classe. Elle parle très bien chinois, parce qu'elle étudie sérieusement tous les jours. Cet après-midi, elle m'a dit joyeusement : "Les ordinateurs de la bibliothèque sont utilisables, allons faire nos devoirs ensemble !" Elle tape vraiment vite ! »\\
Justifications : \zh{班的新同学} et \zh{图书馆的电脑} = possession/appartenance (\zh{的}) ; \zh{说得很好} et \zh{打得真快} = complément de degré après le verbe (\zh{得}) ; \zh{认真地学习} et \zh{高兴地对我说} = adverbe devant le verbe (\zh{地}).
\end{corrige}

\begin{corrige}[Exercice 8 — Remets les mots dans l'ordre]
\begin{enumerate}
\item \zh{我可以出去吗？} \emph{(Wǒ kěyǐ chūqù ma ?)} « Puis-je sortir ? » (Sujet + \zh{可以} + verbe + \zh{吗}.)
\item \zh{他跑得很快。} \emph{(Tā pǎo de hěn kuài.)} « Il court très vite. » (Verbe + \zh{得} + complément de degré.)
\item \zh{我们认真地听老师讲课。} \emph{(Wǒmen rènzhēn de tīng lǎoshī jiǎngkè.)} « Nous écoutons attentivement le professeur faire cours. » (Adverbe + \zh{地} + verbe.)
\item \zh{这是我的新手机。} \emph{(Zhè shì wǒ de xīn shǒujī.)} « C'est mon nouveau téléphone portable. » (\zh{的} de possession entre le pronom et le nom.)
\end{enumerate}
\end{corrige}

\courssub{Je me prépare au Bac}

\begin{corrige}[Exercice 9 — Compréhension de texte et questions de langue (10 points)]
\textbf{Compréhension :}
\begin{enumerate}
\item Après les cours, les élèves peuvent téléphoner à leurs parents et chercher des informations sur Internet (\zh{给家长打电话}、\zh{上网查资料}). (2 points)
\item Parce qu'il n'apprend pas très bien le chinois et qu'il veut utiliser son téléphone pour étudier le chinois après les cours, afin d'apprendre avec plus de plaisir. (2 points)
\end{enumerate}
\textbf{Langue :}
\begin{enumerate}
\item Exemples attendus : \zh{学生上课的时候不可以使用手机。} « Les élèves ne peuvent pas utiliser leur téléphone pendant les cours. » ; \zh{这样我可以学得更开心！} « Comme ça, je pourrai apprendre avec encore plus de joie ! » (2 points : 1 point par phrase relevée et traduite.)
\item Avec \zh{得} : \zh{我汉语学得不太好} — \zh{得} introduit le complément de degré (\zh{不太好}) qui évalue l'action du verbe \zh{学}. Avec \zh{地} : \zh{认真地听老师讲课} — \zh{地} relie l'adverbe \zh{认真} au verbe \zh{听} qu'il modifie. (2 points)
\item — 老师，我\textbf{可以}去喝水\textbf{吗}？ — 现在\textbf{不可以}，下课以后再去。\\
« — Monsieur, puis-je aller boire de l'eau ? — Pas maintenant, tu iras après les cours. » (3 points : 1 point par blanc.)
\end{enumerate}
\textbf{Points de vigilance :} ne pas confondre \zh{得} (après le verbe) et \zh{地} (avant le verbe) ; dans les réponses en français, restituer la nuance de permission de \zh{可以} (« pouvoir », « il est permis de »).
\end{corrige}

\begin{corrige}[Exercice 10 — Thème et version (10 points)]
\textbf{A. Thème (6 points, 2 points par phrase) :}
\begin{enumerate}
\item \zh{在医院里不可以吸烟。} \emph{(Zài yīyuàn li bù kěyǐ xīyān.)} — On accepte aussi \zh{医院里不可以抽烟。}
\item \zh{我的朋友汉语说得很好。} \emph{(Wǒ de péngyou Hànyǔ shuō de hěn hǎo.)} — Attention : le verbe \zh{说} doit être répété devant \zh{得} ; on ne dit pas \zh{说得很好汉语}.
\item \zh{学生们认真地听老师讲课。} \emph{(Xuésheng men rènzhēn de tīng lǎoshī jiǎngkè.)}
\end{enumerate}
\textbf{B. Version (4 points) :} « Maman m'a dit joyeusement : "Ta chambre est vraiment bien rangée ! Maintenant, tu peux regarder la télévision." »\\
Points de langue : \zh{高兴地} (adverbe + \zh{地}), \zh{打扫得真干净} (verbe + \zh{得} + complément de degré), \zh{可以看电视了} (\zh{可以} + verbe, \zh{了} marquant le changement de situation : la permission est maintenant accordée).\\
\textbf{Barème indicatif :} thème 6 points (2 points par phrase : 1 point pour la structure, 1 point pour les caractères et le vocabulaire) ; version 4 points (fidélité du sens 2 points, qualité du français 2 points).\\
\textbf{Points de vigilance :} place de \zh{不} devant \zh{可以} ; répétition du verbe dans la construction avec \zh{得} ; \zh{地} et non \zh{的} devant un verbe.
\end{corrige}

\begin{corrige}[Exercice 11 — Expression écrite (10 points)]
\textbf{Modèle de production :}
\begin{textebox}[Modèle de message]
校长先生，您好！我是毕业班的学生。我的汉语学得不太好，所以我想用学校的Wi-Fi上网查资料，认真地准备考试。请问我可以用学校的Wi-Fi吗？谢谢您！\\
\emph{(Xiàozhǎng xiānsheng, nín hǎo ! Wǒ shì bìyèbān de xuésheng. Wǒ de Hànyǔ xué de bù tài hǎo, suǒyǐ wǒ xiǎng yòng xuéxiào de Wi-Fi shàngwǎng chá zīliào, rènzhēn de zhǔnbèi kǎoshì. Qǐngwèn wǒ kěyǐ yòng xuéxiào de Wi-Fi ma ? Xièxie nín !)}\\
« Monsieur le Proviseur, bonjour ! Je suis élève en classe de Terminale. Je n'apprends pas très bien le chinois, c'est pourquoi je voudrais utiliser le Wi-Fi de l'école pour chercher des informations sur Internet et préparer sérieusement l'examen. Puis-je utiliser le Wi-Fi de l'école, s'il vous plaît ? Merci beaucoup ! »
\end{textebox}
\textbf{Barème indicatif :} respect des consignes et du nombre de caractères (2 points) ; emploi correct de \zh{可以}/\zh{可以吗} (2 points) ; emploi correct d'au moins une particule \zh{的}, \zh{得} ou \zh{地} (2 points) ; exactitude des caractères et de l'ordre des mots (2 points) ; salutation et formule de politesse (2 points).\\
\textbf{Points de vigilance :} ouvrir par une salutation (\zh{校长先生，您好！}) et fermer par \zh{谢谢您！} ; vérifier que \zh{得} suit bien le verbe répété (\zh{学得不太好}) et que \zh{地} précède le verbe (\zh{认真地准备}) ; ne pas oublier \zh{吗} dans la question de permission.
\end{corrige}

\newpage

\coursec{Fiche bilan}

\begin{popfigure}
\begin{tikzpicture}[mindmap, grow cyclic, every node/.style={concept, font=\small},
  root concept/.append style={concept color=popblue, font=\bfseries\small, minimum size=2.6cm},
  level 1/.append style={level distance=3.4cm, sibling angle=60, font=\footnotesize},
  level 2/.append style={level distance=2.2cm, font=\scriptsize}]
\node[root concept, align=center]{Leçon 39\\可以 et 的、得、地}
  child[concept color=poporange]{ node[align=center]{可以 kěyǐ\\modal}
    child { node[align=center]{pouvoir /\\permission} }
    child { node {S + 可以 + V} }
  }
  child[concept color=popgreen]{ node {的 de}
    child { node {nom + 的 + nom} }
    child { node {adj + 的 + nom} }
  }
  child[concept color=poppurple]{ node {得 de}
    child { node {V + 得 + compl.} }
    child { node {degré, manière} }
  }
  child[concept color=poppink]{ node {地 de}
    child { node {adj + 地 + V} }
    child { node {adverbe} }
  }
  child[concept color=popgold]{ node[align=center]{Pièges\\francophones}
    child { node[align=center]{3 « de »\\≠ français} }
    child { node[align=center]{ordre\\des mots} }
  }
  child[concept color=popblueL!80!popblue]{ node[align=center]{Médias\\Module 5}
    child { node[align=center]{网络、电视\\可以…} }
  };
\end{tikzpicture}
\legende{Carte mentale de la leçon 39 : le verbe modal 可以 et les trois particules 的、得、地.}
\end{popfigure}

\begin{bilan}
\textbf{1. Lexique clé (mass médias et télécommunication)}
\begin{center}
\begin{tabularx}{\linewidth}{|X|X|X|X|}
\hline
\rowcolor{popblueL} Caractères & Pinyin & Nature & Traduction \\
\hline
可以 & kěyǐ & verbe modal & pouvoir, avoir le droit de \\
\hline
网络 & wǎngluò & nom & Internet, le réseau \\
\hline
电视 & diànshì & nom & télévision \\
\hline
手机 & shǒujī & nom & téléphone portable \\
\hline
消息 & xiāoxi & nom & nouvelle, information \\
\hline
节目 & jiémù & nom & émission, programme \\
\hline
报纸 & bàozhǐ & nom & journal (papier) \\
\hline
快 & kuài & adjectif & rapide, vite \\
\hline
认真 & rènzhēn & adjectif & sérieux, appliqué \\
\hline
\end{tabularx}
\end{center}

\textbf{2. Règles de grammaire à connaître par cœur}
\begin{center}
\begin{tabularx}{\linewidth}{|X|X|X|}
\hline
\rowcolor{popblueL} Structure & Exemple (caractères + pinyin) & Traduction \\
\hline
Sujet + 可以 + verbe & 我可以上网。Wǒ kěyǐ shàngwǎng. & Je peux aller sur Internet. \\
\hline
Négation : 不可以 + verbe & 上课不可以打电话。Shàngkè bù kěyǐ dǎ diànhuà. & En classe, on ne peut pas téléphoner. \\
\hline
Question : 可以…吗 ? & 我可以用你的手机吗？Wǒ kěyǐ yòng nǐ de shǒujī ma ? & Puis-je utiliser ton portable ? \\
\hline
X + 的 + nom (possession, qualification) & 喀麦隆的新闻 Kāmàilóng de xīnwén & les nouvelles du Cameroun \\
\hline
Verbe + 得 + complément (manière, degré) & 他说得很快。Tā shuō de hěn kuài. & Il parle très vite. \\
\hline
Adjectif + 地 + verbe (adverbe) & 她认真地看报纸。Tā rènzhēn de kàn bàozhǐ. & Elle lit le journal attentivement. \\
\hline
\end{tabularx}
\end{center}

\textbf{3. Méthodes express}
\begin{itemize}
  \item Pour choisir entre 的、得、地, regarde ce qui \emph{suit} la particule : un \cle{nom} $\rightarrow$ 的 ; un \cle{complément de manière ou de degré} après un verbe $\rightarrow$ 得 ; un \cle{verbe} précédé d'un adverbe $\rightarrow$ 地.
  \item Astuce : 的 s'accroche au nom, 得 s'accroche au verbe qui précède, 地 s'accroche au verbe qui suit.
  \item 可以 se place toujours \emph{avant} le verbe principal et \emph{après} le sujet ; jamais en fin de phrase.
  \item Pour nier la permission : 不可以 (interdiction) ; 不能 insiste sur l'incapacité, 可以 sur l'autorisation.
  \item À l'écrit, les trois particules se prononcent toutes « de » (ton neutre) mais s'écrivent différemment : vérifier le caractère avant de rédiger.
\end{itemize}
\end{bilan}

\begin{aretenir}[Checklist avant le Bac]
\begin{itemize}
  \item $\square$ Je sais conjuguer une phrase avec 可以 : Sujet + 可以 + verbe (+ complément).
  \item $\square$ Je sais former la négation 不可以 et la question 可以…吗.
  \item $\square$ Je distingue 可以 (permission) de 能 (capacité) et de 会 (savoir-faire acquis).
  \item $\square$ Je place correctement 的 entre le déterminant et le nom : \zh{我的手机} (wǒ de shǒujī, « mon portable »).
  \item $\square$ Je sais que 得 suit le verbe et introduit un complément de manière ou de degré : \zh{跑得很快} (pǎo de hěn kuài, « courir très vite »).
  \item $\square$ Je sais que 地 transforme un adjectif en adverbe devant un verbe : \zh{慢慢地走} (mànmàn de zǒu, « marcher lentement »).
  \item $\square$ Je n'écris jamais « de » au hasard : je vérifie si c'est 的, 得 ou 地 selon la structure.
  \item $\square$ Je respecte l'ordre des mots chinois : le complément de manière avec 得 ne se place jamais avant le verbe.
  \item $\square$ Je connais le vocabulaire des médias : 网络、电视、手机、报纸、节目、消息.
  \item $\square$ Je sais construire au moins trois phrases originales sur les médias avec 可以, 得 et 地.
  \item $\square$ Je prononce les trois particules avec le ton neutre « de » et je soigne les tons de 可以 (kěyǐ, 3\textsuperscript{e} + 3\textsuperscript{e} ton).
  \item $\square$ Je relis mes productions pour chasser les oublis de 的 dans les groupes nominaux longs.
\end{itemize}
\end{aretenir}

\findecours

\end{document}
